\documentclass[10pt]{article}
\usepackage{amsmath,amssymb,amsthm,mathtools}
\usepackage[margin=1in]{geometry}
\usepackage{booktabs}
\usepackage{makecell}
\usepackage{hyperref}

\newtheorem{theorem}{Theorem}
\newtheorem{proposition}[theorem]{Proposition}
\newtheorem{corollary}[theorem]{Corollary}
\theoremstyle{definition}
\newtheorem{definition}[theorem]{Definition}
\newtheorem{example}[theorem]{Example}

\newcommand{\T}{\mathsf{T}}
\newcommand{\Cl}[1]{\ulcorner #1 \urcorner}
\newcommand{\Ecl}{\mathrel{\equiv_{\mathrm{cl}}}}
\newcommand{\Econd}{\mathrel{\equiv_{\mathrm{cond}}}}
\newcommand{\Ecomp}{\mathrel{\equiv_{\mathrm{comp}}}}
\newcommand{\Epi}{\mathrel{\equiv_{\mathrm{pt}}}}
\newcommand{\V}[1]{\textsc{#1}}

\title{What Can a Prover Certify About Autoformalization Faithfulness?\\
Measuring the Collapse of Closed Equivalence and Certifying Component Differences}
\newif\ifanonymous
\anonymoustrue % set \anonymoustrue for double-blind submission
\ifanonymous
\author{Anonymous}
\date{}
\pdfinfoomitdate=1
\pdftrailerid{}
\hypersetup{pdfauthor={},pdfcreator={},pdfproducer={}}
\else
\author{Anonymous}
\date{\today}
\fi

\begin{document}
\maketitle

\begin{abstract}
Reference-based checks for autoformalization ask whether a candidate statement is provably
equivalent to a trusted reference. Any two true statements are provably equivalent, so such a
check loses discriminating power as its prover gets better. On all 1,140 labelled ProofNetVerif
pairs, closed equivalence accepts 2\% of the unfaithful pairs with fixed tactics and 39\% with
Claude Opus~5.5, and raising the budget of one model takes this from 2\% to 44\%. We introduce a
component criterion whose comparison goals, over the hypotheses and conclusions of canonical
Lean telescopes, are fixed before proof search, so that the prover need not be restricted.
Kernel-checked certificates back every decided component relation, rejections included,
while regrouping certificates tie the compared telescopes to the original statements. These
certificates establish formal differences, not faithfulness to the informal problem; which
components are compared depends on canonicalization and alignment, and faithfulness needs
external evidence. With the same prover our false accepts stay at 4--15 of 694, at a measured
cost in coverage: 56\% of the faithful pairs are accepted, against 86\% by closed equivalence.
Applied to benchmark references, the same goals give kernel-checked evidence for 34 errors in
ProofNet\#, miniF2F and PutnamBench, 20 of them proved false or vacuous. Accepts remain necessary
but not sufficient: constant conclusions and the absorption of a conclusion into the hypotheses
defeat every check that accepts any proof of the component equivalences, so certified
differences complement checks that restrict which proofs count.
\end{abstract}

%======================================================================
\section{Introduction}\label{sec:intro}

A formalization is faithful when it says what the informal statement says. Reference-based
checks replace the informal statement by a trusted formal reference $R$ and ask whether a
candidate $S$ is equivalent to it: BEq+~\cite{beq} proves $R\to S$ and $S\to R$ with restricted
tactics, shadow benchmarks~\cite{shadowbench} check $S$ against expert-written consequences of
$R$, and related tests compare consequences or downstream proofs~\cite{bpf,t2}. All of them face
the same fact: if $R$ and $S$ are both theorems then $R\leftrightarrow S$ is a theorem, so a
complete prover declares every true candidate equivalent to every true reference. This is
folklore, and the field has answered it by restriction. BEq+ requires each proof to use the other
statement, shadow benchmarks only allow direct projections, the case for theory-level
autoformalization argues that equivalence should be checked in a restricted fragment~\cite{tla},
and the recently introduced SemantEq~\cite{semanteq} defines equivalence through chains of
pre-specified, audited transformation steps and checks the resulting proof terms, rather than
accepting arbitrary proofs of mutual implication. Other pipelines use language-model judges, alone or as reinforcement-learning
rewards~\cite{beyondcomp,formarl,reform,criticlean}, or BEq itself as the reward~\cite{stepfun}.

Restriction trades recall for precision, and judges carry no guarantee. We ask three questions.

\emph{RQ1: how badly does unrestricted (closed) equivalence fail with the provers people actually
use, and how does the failure grow with the prover?} For one prover and a stronger one that
proves a superset, the closed false-accept set can only grow; we call this logical fact the
\emph{monotonicity paradox}. Under two of four prompts, SemantEq reports that a language model
produces proofs in both directions for about two fifths of its non-equivalent pairs, which is why
it audits proof terms rather than relying on compilation alone~\cite{semanteq}; we measure how
the effect grows across models, ordered by the closed proofs they find, and across the budgets of
one model. \emph{RQ2: without restricting the prover, can a check give kernel-checked evidence of
a formal difference and locate it in a fixed component?} SemantEq secures the meaning of an
accept by restricting the proofs it may use; a rejection there means only that no proof passed
its audits. \emph{RQ3: what does such evidence settle, what does it cost, and where does it still
fail?} This paper answers the three on real benchmark data and is explicit about what is not
solved.

We do not claim to have found that closed equivalence collapses on true statements, nor that
comparing hypotheses and conclusions separately is new: checks of premises, of conclusions and
of the two directions of an entailment are in use (\S\ref{sec:related}), and fixing the rule
before any prover runs is also how SemantEq keeps the meaning of an accept independent of the
prover~\cite{semanteq}. Our claim is narrower
and testable: a component-level decision procedure whose verdicts are bound to the original
statements by kernel certificates can separate three cases that a closed check conflates (a
difference located in the aligned hypotheses or conclusions, a difference certified only at the
level of the closed statements, and pairs it cannot tell apart), at a coverage cost that can be
measured.

The claim rests on five properties together, none of which suffices alone: the
comparison goals are generated from the two statements before any prover runs, so the prover
decides fixed goals and never chooses a witness that seems to separate the statements; every
decided verdict is a kernel proof, rejects included, so a reject certifies that the components
differ rather than that no proof was found; regrouping certificates bind the closure of each side's
components to its original statement; the outcomes keep apart a located difference, a
difference certified only between the closed statements, an alignment that ends early, an
undecided pair and the accepts that still collapse; and on 1,140 labelled pairs we count each
outcome, its cost against closed equivalence with the same prover, and where it fails.

The guarantee has three
layers, which we keep apart throughout. The kernel guarantees that every decided goal holds in
Lean with Mathlib. The canonicalization and alignment rules decide which components the goals
are about; the regrouping certificates guarantee that the components recompose to the original
statements, not that they are the components a mathematician would compare. A located
difference is therefore a kernel fact about the aligned components, not a finding that the
candidate is unfaithful there. Faithfulness to the informal statement needs evidence beyond
both, and we do not supply it automatically: the labels say whether a pair is unfaithful, not
where, and only a reader can confirm that a location is the reason. We keep mechanical
locations and confirmed unfaithfulness apart throughout. The evidence is for candidate
evaluation and for auditing the references themselves, and anyone can re-check it: code, data
and the Lean developments are at \url{https://anonymous.4open.science/r/4444mmmm-FC0B/}, where
\texttt{paper\_check.py} recomputes every main table of this paper from the stored results
(Appendix~\ref{app:repro} lists the units, denominators and label versions of every count).

\paragraph{Contributions.} One claim per question, each with the evidence that tests it.
\begin{enumerate}
\item[A.] \textbf{The failure, measured} (RQ1, \S\ref{sec:measure}). Closed false accepts grow
with proving power: on all 1,140 labelled ProofNetVerif pairs~\cite{proofnetverif}, closed false
accepts rise from 1.7\% with tactics to 39\% (95\% interval 36--43\%) with Opus across models,
and from 12 to 305 (44\%) within one model as its budget grows (Table~\ref{tab:budget}). The
budget sweep nests the provers; across models, which are not nested, the growth is an observed
association that Proposition~\ref{thm:mono} explains, not a law.
\item[B.] \textbf{Certified component comparison} (RQ2, \S\ref{sec:framework}). Comparison goals
over the hypotheses and conclusions of canonical Lean telescopes are fixed before proof search,
so the prover may be unrestricted; every decided verdict, rejects included, is a kernel proof
about the exhibited telescopes, and regrouping certificates tie the telescopes to the original
statements. We state what each verdict certifies, what it is relative to (the canonical split
and alignment, Proposition~\ref{prop:split}) and what it does not show: a certificate
establishes a formal difference, not unfaithfulness (Tables~\ref{tab:layers} and
\ref{tab:scope}).
\item[C.] \textbf{What it gains, what it costs, and where it fails} (RQ3,
\S\ref{sec:measure}--\ref{sec:audit}). Against those closed false accepts, with the same prover
ours stay at 4--15, at fewer model calls and tokens than the closed check spends
(Table~\ref{tab:cost}). The cost is coverage, reported outcome by outcome
(Table~\ref{tab:dist}): 137 faithful pairs are accepted only by closed equivalence, 95 of them
returned as \V{residual} because the alignment ends early, and most undecided pairs fail in
canonicalization and alignment, not in proof search. What the cost buys is a mechanical location:
244 of the 328 rejects of unfaithful pairs name the aligned component that differs, which a
closed verdict never does; whether it is where the candidate is unfaithful is for a reader to
confirm. Independent evidence bounds the claim: an external adjudication of the pairs labelled
equivalent~\cite{semanteq}, selection-dependent because most of its reviews were prompted by
another method's failures, relabels 46 of our 51 rejects there as non-equivalent, and in a blind
two-annotator audit (\S\ref{sec:blind}), the 43 component-goal rejects among them rest on 30
distinct statements; of the 26 distinct statements on which the annotators agree, 16 are
unfaithful candidates, 7 faithful candidates whose reference says something else and 3
canonicalization artefacts, and a reader confirms 9 of our 14 false accepts (6 of 11 distinct
statements). The failures are measured too: accepts still collapse on constant conclusions and
absorption (\S\ref{sec:survive-exp}), and rules written to remove the artefacts the audit found
create others of the same kind (Appendix~\ref{app:k19}). Two applications use the certificates:
combined with SemantEq's accepts, our certified rejects give a two-sided evaluator that decides
54\% of the shared test pairs with no false accept observed under either label set, though with
false rejects (\S\ref{sec:measure}); and applied to benchmark references, the same goals give
kernel-checked evidence for 34 errors that manual curation left in ProofNet\#, miniF2F and
PutnamBench (\S\ref{sec:audit}).
\end{enumerate}
Supporting experiments test known failure modes: candidates constructed or generated by a
language model to collapse (\S\ref{sec:survive-exp}), which closed equivalence rewards and our
accepts partly still do, and best-of-$n$ selection (\S\ref{sec:select}), where certified rejects
veto LLM-judge accepts without trusting either side.
We do not claim a better binary accept than BEq+, which is as precise and sometimes recalls more
where candidates are close to the reference in form (Table~\ref{tab:g4}), or than SemantEq, whose
model-free search reports 100\% precision on the ProofNetVerif test split with 51.4\% recall
under the published labels and 65.6\% under its adjudicated ones~\cite{semanteq} (ours: 94.7\%
and 55.6\% under the published labels, on a partly different set of pairs). Nor do we claim an
automatic faithfulness verdict: an accept is a necessary condition, and a reject is a certified
formal difference whose relevance needs the informal statement. What the criterion adds is a
certificate for every decided verdict, rejects included, and a precise account of the accepts
that remain relative to the theory.

%======================================================================
\section{A certified component criterion}\label{sec:framework}

\begin{definition}[Telescope]
A statement is a triple $S=(\Gamma,H,C)$ of a context $\Gamma$ of typed variables $\bar x$,
a hypothesis $H(\bar x)$ (the conjunction of the Prop-typed binders) and a conclusion
$C(\bar x)$. Its closure is $\Cl S:=\forall\bar x{:}\Gamma.\,H\to C$.
\end{definition}

$\Gamma$ is the ordered, dependent context of the elaborated Lean term: later types may mention
earlier variables, and it includes implicit arguments and type-class instances, which are data.
$H$ collects the Prop-valued binders in order, and is a conjunction only where no later type
mentions their proofs: the first Prop binder whose proof is used later ends the prefix, and it
and everything after it stay inside $C$, apart from hoisted hypotheses (below). $C$ is the remaining
body after the normalizations below, so it may differ syntactically from the source conclusion;
the regrouping certificates below prove that the closure is unchanged. For example,
\texttt{theorem (G : Type) [Group G] [Fintype G] (h : Nat.card G = 2) (a : G) : a * a = 1}
has $\Gamma=(G,\ \mathit{inst}_{\mathrm{Group}},\ \mathit{inst}_{\mathrm{Fintype}},\ a)$,
$H=(\mathrm{card}\,G=2)$ and $C=(a\cdot a=1)$.

The decomposition is computed by Lean (\texttt{CanonV2}): we take the maximal telescope of the
elaborated term, read $\neg\exists x.P$ as $\forall x.\neg P$, keep a final $\neg Q$ or $a\ne b$
in the conclusion, eliminate one-point binders $x=t$, and hoist hypotheses that do not depend on
later binders. Binders of $R$ and $S$ are aligned by type and role, jointly for the shared prefix.
The rules follow three principles. Each is a rewrite whose effect on the closure the kernel
re-checks (the regrouping certificates below), so a rule can cost coverage but not soundness.
Each is the same for $R$ and $S$ and is applied before any prover runs. And each removes a
presentation choice that textbooks and benchmarks make freely (pushing a negation through an
existential, naming a value with $x=t$, the order of independent hypotheses) while keeping in the
conclusion what is asserted there. We did not take the finest split (every Prop binder a
hypothesis) or the coarsest (everything in the conclusion): the coarsest is closed equivalence,
and the finest moves a final implication $P\to Q$ of the conclusion into the hypotheses, which
changes vacuity-only verdicts (Proposition~\ref{prop:split}). Appendix~\ref{app:sens} switches
each rule off, and Appendix~\ref{app:more} tests stability under 299 surface re-presentations
(binder order, conjoined or curried hypotheses, negation forms). Renaming is absorbed by
$\alpha$-equivalence of Lean terms; implicit arguments and coercions are not varied
systematically, and an inserted instance binder is a known failure of the positional
alignment (\S\ref{sec:conclusion}).
For aligned statements we compare
\begin{align*}
R\Ecl S &:\iff \Cl R\leftrightarrow\Cl S,\\
R\Econd S &:\iff \forall\bar x.\ (H_R\leftrightarrow H_S)\wedge\bigl(H_R\to(C_R\leftrightarrow C_S)\bigr),\\
R\Ecomp S &:\iff \forall\bar x.\ (H_R\leftrightarrow H_S)\wedge(C_R\leftrightarrow C_S),
\end{align*}
with $\Ecomp\Rightarrow\Econd\Rightarrow\Ecl$.

\begin{definition}[Verdicts]
A checker searching for proofs in a theory $\T$ (Lean + Mathlib) returns
\V{accept} if $\T\vdash R\Ecomp S$;
\V{reject} if $\T\vdash\neg(R\Econd S)$;
\V{off-domain} if $\T\vdash R\Econd S$ and $\T\vdash\neg(R\Ecomp S)$;
and \V{unknown} otherwise. When the alignment of the two contexts ends early
(\S\ref{sec:survive}), a proof of $R\Ecomp S$ is reported as \V{residual} instead of \V{accept}.
\end{definition}

Every verdict other than \V{unknown} is a Lean proof of one of five goals generated from
$(R,S)$ by a Lean metaprogram before any prover is called: \texttt{acc}, \texttt{cond},
\texttt{nhyp} (hypotheses differ somewhere), \texttt{ncin} (conclusions differ inside $H$) and
\texttt{ncomp}. Because the goals are fixed in advance, the meaning of a certificate does not
depend on the prover, and a stronger prover, one that proves a superset of what another proves,
can only turn \V{unknown} into a decided verdict; different models need not be stronger in this
sense, so across models this is a tendency, not a guarantee.
\V{Unknown} is epistemic: the budget ran out. \V{Off-domain} is semantic: the statements
provably agree on every input satisfying the hypotheses and provably differ elsewhere, as for
the extreme value theorem against its quantifier swap ($\exists x\in K\,\forall y\in K$ vs.\
$\forall y\in K\,\exists x\in K$), which differ only at $K=\emptyset$, or for a junk value of total
division. Whether such a difference matters cannot be decided from the formal statements: no
relation that is sound, closed under restriction and sensitive to the raw conclusion can ignore
differences outside the hypotheses (Proposition~\ref{thm:nogo}). The criterion therefore reports
it instead of folding it into \V{accept} or \V{reject}.

\begin{example}[Running example]\label{ex:run}
Let $R$ be $\forall n{:}\mathbb N,\ 1\le n\to 2\le 2^n$. Every candidate below is a true theorem,
so closed equivalence accepts all of them.
\begin{center}\small
\begin{tabular}{lll}
\toprule
candidate $S$ & verdict & certificate\\
\midrule
$S_1$: $\forall n,\ 0<n\to 2\le 2^n$ & \V{accept} & hypotheses and conclusions agree at every $n$\\
$S_2$: $\forall n,\ n\le n^2$ & \V{reject} & hypotheses differ at $n=0$\\
$S_3$: $\forall n,\ 2\le n\to 2\le 2^n$ & \V{reject} & hypotheses differ at $n=1$ (candidate weaker)\\
$S_4$: $\forall n,\ 1\le n\to 2\le 2^n\vee n=0$ & \V{off-domain} & agree where $1\le n$; conclusions differ at $n=0$\\
$S_5$: $\forall n,\ 1\le n\wedge 2\le 2^n\to 2\le 2^n$ & \V{accept} & absorption; flagged trivial (\S\ref{sec:survive})\\
\bottomrule
\end{tabular}
\end{center}
\end{example}

\paragraph{Certificates.} The components are computed by code, so the checker also proves two
regrouping goals against the original elaborated statements,
$\texttt{split\_r}:\Cl R\leftrightarrow\forall\bar x.\,H_R\to C_R$ and its counterpart
\texttt{split\_s} for $S$.

\begin{proposition}[Certificate semantics]\label{prop:cert}
Suppose that for $(R,S)$ the checker exhibits telescopes $T_R=(\bar x,H_R,C_R)$,
$T_S=(\bar x,H_S,C_S)$, kernel proofs of \texttt{split\_r} and \texttt{split\_s}, and a proof of one
of the five goals. Then \V{accept} and \V{residual} give $T_R\Ecomp T_S$ and hence $\Cl R\leftrightarrow\Cl S$;
\V{reject} gives $\neg(T_R\Econd T_S)$; \V{off-domain} gives $T_R\Econd T_S$ and
$\neg(T_R\Ecomp T_S)$. These facts hold whatever code produced $T_R$ and $T_S$. A closed
refutation $\T\vdash\neg(\Cl R\leftrightarrow\Cl S)$ together with both regrouping certificates
also gives $\neg(T_R\Econd T_S)$.
\end{proposition}

The proof is composition of certificates in the kernel. Every decided verdict we report carries
both certificates; a pair whose goal is proved but whose \texttt{split\_r} or \texttt{split\_s}
is not counts as \V{unknown}. A canonicalizer bug that changes a closure therefore costs coverage,
not correctness. The last sentence is a second route to \V{reject}, and every table uses it: a
pair whose component goals stay open but whose closed equivalence is refuted, with both
certificates, is a reject that certifies a difference without locating it. So is a component
reject when neither telescope has a binder or a hypothesis in its prefix, since its conclusion
goal then compares the closed statements.

\paragraph{Decision procedure.} For a pair $(R,S)$ the checker
\begin{enumerate}\setlength{\itemsep}{0pt}
\item elaborates both statements, computes their canonical telescopes and aligns the binders; if
canonicalization fails, the verdict is \V{unknown};
\item generates, before any prover runs, the two regrouping goals, the five component goals and
the closed goals: both directions of $\Cl R\leftrightarrow\Cl S$ and their negations;
\item runs the prover budget on every goal (fixed tactics, then model proposals and repair
rounds) and keeps a goal as proved only if the kernel checks its proof;
\item returns \V{unknown} unless both regrouping goals are proved, and otherwise \V{accept} if
\texttt{acc} is proved (\V{residual} if the alignment ended early), \V{reject} if \texttt{nhyp} or
\texttt{ncin} is, \V{off-domain} if \texttt{cond} and \texttt{ncomp} are, \V{reject} if only the
closed equivalence is refuted, and \V{unknown} otherwise;
\item reports a reject as located when it rests on \texttt{nhyp} or \texttt{ncin} and a telescope
has a prefix binder or a hoisted hypothesis, and attaches to an accept the one-sided triviality
and constancy certificates of \S\ref{sec:survive} as flags, which do not change the verdict.
\end{enumerate}
In a consistent theory no two of the cases in step~4 can both be proved, so their order only
fixes how a run is read.

Table~\ref{tab:layers} separates what is guaranteed by the kernel from what depends on the rules
and from what needs the informal statement. Only the first layer is unconditional; the second is
measured in \S\ref{sec:measure} and Appendix~\ref{app:sens}, and the third is outside the
criterion.

\begin{table}[t]
\centering\small
\begin{tabular}{p{0.17\linewidth}p{0.4\linewidth}p{0.35\linewidth}}
\toprule
layer & what it guarantees & where it fails, and where we measure it\\
\midrule
kernel & every decided goal, and both regrouping certificates, hold in Lean with Mathlib & never,
given the trusted kernel; resampling the prover never flips a verdict (\S\ref{sec:measure})\\
rules & which components the goals compare: the canonical telescope and the aligned binders; the
certificates guarantee only that the components recompose to the original statements & a partial
alignment can pair binders of different meaning (22 of the 51 rejects of labelled-equivalent pairs locate a difference on one);
canonicalization artefacts (Appendix~\ref{app:rejects}); \V{residual} where the alignment ends early\\
informal & that a certified difference matters and lies where it is located, or that an accept is faithful & labels, an external adjudication (\S\ref{sec:measure}) and the
blind audit (\S\ref{sec:blind}); accepts that collapse (\S\ref{sec:survive})\\
\bottomrule
\end{tabular}
\caption{Three layers of guarantee. A verdict is a kernel fact about the components that the
rules exhibit; that these are the right components, and that the difference matters for the
informal statement, are separate claims.}
\label{tab:layers}
\end{table}

\paragraph{What the verdicts are relative to.} The regrouping certificates bind the closure of
each telescope to its statement; they do not bind the split of that closure into hypotheses and
conclusion, nor the alignment of the two contexts. For $R$ of Example~\ref{ex:run}, the telescope
$(n,\top,\,1\le n\to 2\le 2^n)$ has the same closure and passes \texttt{split\_r}, but compared
with $S_1$ it yields a \V{reject} at $n=0$. A component verdict is therefore a certified fact about
the exhibited telescopes and alignment, and the canonicalization rules above are part of the
definition of the criterion, not an implementation detail. This dependence cannot be removed
without losing what the criterion is for. Folding the context away collapses every sound relation
to $\Ecl$ (Proposition~\ref{thm:repr}), and already allowing a Prop to move between hypotheses and
conclusion, with the alignment fixed, collapses every sound relation closed under restriction to
$\Epi$, which relates any two true statements (Proposition~\ref{prop:split}). The split and the
alignment are the information that lets the criterion separate true statements.

The dependence is also confined. Two true statements are $\Epi$-related, so a \V{reject} outside
$\Epi$, certified by the goal \texttt{npimp} (Appendix~\ref{app:theory}), certifies that one of the
two statements is false, survives every regrouping of either side, and, together with a direction
goal, says which side is false. Only \emph{vacuity-only} rejects, where the pointwise claims
$H\to C$ agree and the hypotheses differ, depend on the split, and they become accepts once both
sides are fully curried. We make the remaining dependence explicit rather than removing it: the
rules are deterministic and the same for $R$ and $S$, every verdict exhibits both telescopes,
Proposition~\ref{prop:canon} lists the re-presentations under which the canonical form is
invariant, 292 of 299 surface re-presentations are proved component-equivalent to the original
(Appendix~\ref{app:more}), and switching each rule off changes decided verdicts only on the
split, never a reject certified by \texttt{npimp} (Appendix~\ref{app:sens}). Where the rules split two presentations of the same content differently, the verdict
changes: 2 of the 14 rejects on labelled-equivalent pairs of the 300-pair subset are such
canonicalization artefacts, both vacuity-only as the proposition predicts.

\paragraph{Why components, and what a reject means.} The three relations are the greatest ones
satisfying nested lists of four closure properties (Proposition~\ref{thm:char}): (S)~soundness,
$R\sim S\Rightarrow\Cl R\leftrightarrow\Cl S$, alone gives $\Ecl$; adding (Res)~closure under
restriction $(\Gamma,H,C)\mapsto(\Gamma,H\wedge P,C)$ gives pointwise implication
$R\Epi S:\iff\forall\bar x.\,(H_R\to C_R)\leftrightarrow(H_S\to C_S)$, which still relates any two true theorems;
adding (Vac)~that related statements have hypotheses both satisfiable or both unsatisfiable gives
$\Econd$; and adding (Weak)~that the raw claims $(\Gamma,\top,C)$ are related as well gives $\Ecomp$. Between two true statements, vacuity respect is the first property
that separates anything, and it separates exactly at a point where the hypotheses differ. A
\V{reject} is therefore a certificate that hypotheses or conclusions differ at some point (the
goals are negated universal statements, so the point exists but the proof need not name it),
not by itself a verdict of unfaithfulness: a decorative hypothesis ($1\le m$ where $m=0$ is
degenerate) is separated as surely as a substantive one, and telling them apart needs the
informal statement. Four further goals over the same components classify each reject as
\emph{vacuity-only} (the pointwise claims $H\to C$ agree), \emph{weaker}, \emph{stronger} or
incomparable (Appendix~\ref{app:theory}); a reject between two true statements is always
vacuity-only, and one certified not to be (\texttt{npimp}) certifies a false statement.

\paragraph{Prover strength.} Model a prover by the set of propositions it proves within its
budget, and compare a prover with a stronger one that proves a superset
(Proposition~\ref{thm:mono}). (i)~Every decided verdict set grows and \V{unknown} shrinks.
(ii)~Closed equivalence accepts when it proves $\Cl R\leftrightarrow\Cl S$, so its false-accept
set grows and, for a prover that proves every theorem of $\T$, contains every unfaithful pair
whose two statements are both provable or both refutable in $\T$, in particular every provable
unfaithful candidate of a provable reference. (iii)~The component false-accept set also grows, but it is contained in
$\{(R,S)\text{ unfaithful}: R\Ecomp S\}$, a set fixed by the alignment and the theory before any
prover runs. The question is what that set contains. The provers of \S\ref{sec:measure} are
different models, not nested in this sense, so there the proposition explains the mechanism and the
tables measure its size.

\subsection{What an accept certifies}\label{sec:survive}

\paragraph{Accept and residual.} The alignment can stop early: a data binder of one
side may have no counterpart (an extra \texttt{Subgroup} binder, an inserted instance), or the
statement may have no binders at all. Everything after the cut is then quantified inside the
conclusion and compared as a closed proposition, so the collapse returns one level down: if both
remainders are provable, $R\Ecomp S$ reduces to $\forall\bar x.\,H_R\leftrightarrow H_S$, whatever the
remainders say (Proposition~\ref{prop:residue}). A proof of $R\Ecomp S$ is therefore an
\V{accept} only when the alignment is total and the prefix nonempty; otherwise it is
\V{residual}. The distinction is computed from the canonical telescopes alone, before any prover
runs. An \V{accept} certifies $H_R(\bar x)\leftrightarrow H_S(\bar x)$ and
$C_R(\bar x)\leftrightarrow C_S(\bar x)$ at every aligned input; a \V{residual} certifies this
only for the binders that were aligned, and for the rest inherits the collapse of closed
equivalence. On the full set its precision is that of closed equivalence (95 faithful against 100
unfaithful pairs with Opus, against 384 and 273 for closed equivalence), so it is not evidence of
faithfulness. The tables report it in a column of its own.

\paragraph{Two collapses inside an accept.} Pointwise agreement is still agreement relative
to $\T$, and it collapses whenever the content of the problem is itself provable pointwise.
\emph{Constant conclusions}: if $\forall\bar x.\,C_R$ and $\forall\bar x.\,C_S$ are both theorems, an
accept certifies only the hypotheses. \emph{Absorption}: for $R=(\Gamma,H,C)$ let
$R^{\mathrm{abs}}:=(\Gamma,H\wedge C,C)$. Then $R\Ecomp R^{\mathrm{abs}}$ holds exactly when
$\forall\bar x.\,H\to C$, i.e.\ when $R$ is true, although $R^{\mathrm{abs}}$ is a theorem by
projection and has lost the problem. The same holds for every relation under which statements with
provably equivalent hypotheses and identical conclusions are related, which includes every check
that accepts any proof of the component equivalences: such a check cannot tell a true statement
from its absorption once the prover proves the statement. So faithfulness is not a property of
the extensions of the components, and a remedy has to restrict how a statement may be written or
which proofs count, as BEq+~\cite{beq}, shadow benchmarks~\cite{shadowbench}, theory-level
checks~\cite{tla} and SemantEq's allowed steps and proof-term audits~\cite{semanteq} do.

We use two one-sided certificates that depend on one statement at a time. They annotate an
accept without changing it, and a rule that uses them can only withhold accepts. \emph{Constancy}: the goals $\forall\bar x.\,C$ and $\forall\bar x.\,\neg C$ on the aligned
context; an accept whose conclusion is certified constant is a \emph{collapse}, one whose two
constancy goals are both refuted is \emph{informative} (Appendix~\ref{app:theory}).
\emph{Triviality}: $S$ closes by \texttt{intros; first | assumption | tauto} within 20{,}000
heartbeats while $R$ does not. Both certify that the candidate is degenerate, not that it is
unfaithful, and \S\ref{sec:survive-exp} measures how much of the false-accept set they remove and
what they cost on faithful pairs.

\paragraph{Embeddings and canonical forms.} For type drift along an embedding $e$
($\mathbb N\hookrightarrow\mathbb Z$, $\mathbb N\to\mathsf{ZMod}\,n$) we compare $S$ with $R\circ e$:
agreement on the image of $e$ implies $\Cl R\to\Cl S$, and the checker reports \V{off-image} when it
also certifies that $e$ is not surjective. The canonical form is stable under swapping independent
binders, currying adjacent hypotheses in the shared prefix, and writing $\neg P$ as
$P\to\mathrm{False}$ (Proposition~\ref{prop:canon}); Appendix~\ref{app:more} tests this on 299
re-presentations. All propositions of this section and Appendix~\ref{app:theory} have short
proofs; the non-metamathematical ones except Proposition~\ref{prop:canon} are also checked in
Lean~4 (\texttt{T1}--\texttt{T9}, listed in Table~\ref{tab:theory}).
Table~\ref{tab:terms} fixes the vocabulary used in the rest of the paper.

\begin{table}[t]
\centering\small
\begin{tabular}{lp{0.68\linewidth}}
\toprule
term & meaning\\
\midrule
\multicolumn{2}{l}{\emph{Verdicts: kernel proofs of goals fixed before proving}}\\
\V{accept} & $T_R\Ecomp T_S$ with total alignment and nonempty prefix: hypotheses and conclusions agree at every input\\
\V{residual} & $T_R\Ecomp T_S$ with an alignment that ends early: the unaligned remainder is compared closed and inherits its collapse; not evidence of faithfulness\\
\V{reject} & $\neg(T_R\Econd T_S)$: they differ at some input inside the hypotheses\\
\V{off-domain} & $T_R\Econd T_S$ and $\neg(T_R\Ecomp T_S)$: differences only outside the hypotheses\\
\V{off-image} & agreement along an embedding that is certified not surjective\\
\V{unknown} & no goal proved within the budget, or a regrouping certificate missing\\
\midrule
\multicolumn{2}{l}{\emph{Annotations; they never change a verdict}}\\
constancy & $\forall\bar x.\,C$ or $\forall\bar x.\,\neg C$ proved: \emph{collapse}; both refuted: \emph{informative}\\
triviality & the candidate closes by propositional logic and the reference does not (a heuristic)\\
direction & a reject is \emph{vacuity-only} (depends on the split) or, certified by \texttt{npimp}, \emph{weaker}, \emph{stronger} or incomparable (one side is false)\\
\midrule
\multicolumn{2}{l}{\emph{Phenomena}}\\
collapse & an accept that holds because both sides are provable, not because they say the same\\
absorption & $(\Gamma,H\wedge C,C)$, component-equivalent to a true $(\Gamma,H,C)$\\
artefact & a verdict caused by how the canonical rules split a closure\\
\bottomrule
\end{tabular}
\caption{Terms. Only the six verdicts are outputs of the criterion; everything else annotates them.}
\label{tab:terms}
\end{table}

%======================================================================
\section{Measuring the collapse}\label{sec:measure}

\paragraph{Data and provers.} ProofNetVerif~\cite{proofnetverif} pairs candidate formalizations of
ProofNet test problems with human equivalence labels; 1,140 labelled pairs (446 equivalent, 694
inequivalent) have statements that elaborate in our Mathlib, and a stratified subset of 300 (150
each) is used where an experiment is expensive. ProofNet-Verified~\cite{pnv} pairs 317 ProofNet\#
statements with manually corrected versions and records the relation of each original to its
informal statement, so its labels are not ours. Each goal is attacked by fixed tactics, then by
one proposal and one repair round of a language model whose output is checked by the kernel:
Qwen3.8-27B without thinking, GPT-4.1-mini, DeepSeek-V4.1-Flash, Gemini~4 Argon, GPT-6 Astra and
Claude Opus~5.5. The closed baseline proves $\Cl R\leftrightarrow\Cl S$ with the same tactics, model
and rounds; Table~\ref{tab:cost} compares what the two sides spend. BEq+ is our port of the LeanInteract implementation (\texttt{beq\_plus}, including
\texttt{exact?} reference checks and \texttt{convert}-based saturation). Model proofs may not
contain \texttt{sorry}, \texttt{admit}, \texttt{native\_decide} or new axioms; accepted model
proofs are recompiled alone in a fresh file with \texttt{\#print axioms}, and every one rechecked
uses only \texttt{propext}, \texttt{Classical.choice} and \texttt{Quot.sound}
(Appendix~\ref{app:repro} states which proofs this covers).

\begin{table}[t]
\centering\small
\begin{tabular}{lrrrrrrr}
\toprule
& \multicolumn{2}{c}{closed $\Ecl$} & \multicolumn{2}{c}{ours, \V{accept}} & \multicolumn{2}{c}{ours, \V{residual}} & certified\\
\cmidrule(lr){2-3}\cmidrule(lr){4-5}\cmidrule(lr){6-7}
prover & FA/694 & acc/446 & FA/694 & acc/446 & FA/694 & acc/446 & rejects (U / F)\\
\midrule
tactics & 12 & 215 & 4 & 188 & 1 & 38 & 5 / 0\\
Qwen3.8-27B & 18 & 232 & 5 & 203 & 4 & 42 & 17 / 2\\
GPT-4.1-mini & 12 & 238 & 4 & 193 & 2 & 39 & 5 / 0\\
DeepSeek-V4.1-Flash & 108 & 328 & 9 & 249 & 41 & 82 & 168 / 14\\
Claude Opus 5.5 & \textbf{273} & 384 & 14 & 248 & 100 & 95 & 328 / 51\\
\midrule
BEq+ & \multicolumn{2}{c}{FA 2, acc 207} & & & & &\\
\bottomrule
\end{tabular}
\caption{All 1,140 labelled ProofNetVerif pairs (694 inequivalent, 446 equivalent). FA: accepts of
inequivalent pairs; acc: accepts of equivalent pairs; certified rejects on inequivalent (U) and
equivalent (F) pairs. Gemini~4 Argon and GPT-6 Astra were run on the subsets only
(Table~\ref{tab:sweep}). \V{Residual}: a proof of $\Ecomp$ whose alignment ends early
(\S\ref{sec:survive}). Wilson 95\% intervals for the FA rate with Opus: closed 35.8--43.0\%, ours
1.2--3.4\%.}
\label{tab:full}
\end{table}

\begin{table}[t]
\centering\small\setlength{\tabcolsep}{4pt}
\begin{tabular}{lrrrrrrrr}
\toprule
& \multicolumn{4}{c}{inequivalent (694)} & \multicolumn{4}{c}{equivalent (446)}\\
\cmidrule(lr){2-5}\cmidrule(lr){6-9}
outcome & $n$ & \% & 95\% CI & $\Ecl$ & $n$ & \% & 95\% CI & $\Ecl$\\
\midrule
\multicolumn{9}{l}{\emph{component certificate, total alignment}}\\
\V{accept} & 14 & 2.0 & 1.2--3.4 & 14 & 248 & 55.6 & 51.0--60.1 & 247\\
\V{reject}, located & 140 & 20.2 & 17.4--23.3 & 54 & 17 & 3.8 & 2.4--6.0 & 8\\
\multicolumn{9}{l}{\emph{component certificate, alignment not total}}\\
\V{residual} & 100 & 14.4 & 12.0--17.2 & 88 & 95 & 21.3 & 17.8--25.3 & 95\\
\V{reject}, located & 104 & 15.0 & 12.5--17.8 & 34 & 22 & 4.9 & 3.3--7.4 & 5\\
\multicolumn{9}{l}{\emph{certificate about the closed statements only}}\\
\V{reject}, not located & 84 & 12.1 & 9.9--14.7 & 0 & 12 & 2.7 & 1.5--4.6 & 0\\
\quad closed equivalence refuted & 54 & 7.8 & 6.0--10.0 & 0 & 8 & 1.8 & 0.9--3.5 & 0\\
\quad component goal, both prefixes empty & 30 & 4.3 & 3.0--6.1 & 0 & 4 & 0.9 & 0.3--2.3 & 0\\
\multicolumn{9}{l}{\emph{undecided}}\\
\V{unknown} & 252 & 36.3 & 32.8--40.0 & 83 & 52 & 11.7 & 9.0--15.0 & 29\\
\quad canonicalization failed & 3 & 0.4 & 0.1--1.3 & 2 & 2 & 0.4 & 0.1--1.6 & 0\\
\quad regrouping certificate missing & 85 & 12.2 & 10.0--14.9 & 12 & 12 & 2.7 & 1.5--4.6 & 5\\
\quad no binder aligned & 59 & 8.5 & 6.6--10.8 & 12 & 4 & 0.9 & 0.3--2.3 & 3\\
\quad partial alignment & 30 & 4.3 & 3.0--6.1 & 12 & 8 & 1.8 & 0.9--3.5 & 0\\
\quad full alignment, goals open & 75 & 10.8 & 8.7--13.3 & 45 & 26 & 5.8 & 4.0--8.4 & 21\\
\midrule
all pairs & 694 & & & 273 & 446 & & & 384\\
\bottomrule
\end{tabular}
\caption{Every outcome with Opus on all 1,140 pairs (the configuration of Table~\ref{tab:full}),
grouped by the evidence behind it, with Wilson intervals; each pair has exactly one outcome, and
no pair is \V{off-domain}. $\Ecl$: how many of the pairs closed equivalence accepts with the
same prover, so the column shows where its 273 false and 384 correct accepts go. Located: a
component goal says whether the aligned hypotheses or conclusions differ, a mechanical location
rather than a reader's judgment; the three reject rows sum to 328 and 51. Not located: the reject certifies only that the closed statements differ,
through a refuted closed equivalence with both certificates (Proposition~\ref{prop:cert}) or
because neither telescope has a binder or a hypothesis in its prefix, so that the conclusion
goal compares the closed statements. \V{Unknown} is split by where it arises, and only its last
line is a budget failure; canonicalization failures and missing certificates are
\V{unknown}, never \V{reject}.}
\label{tab:dist}
\end{table}

\begin{table}[t]
\centering\small
\begin{tabular}{llrrrrrr}
\toprule
& & \multicolumn{2}{c}{closed $\Ecl$} & \multicolumn{2}{c}{ours, \V{accept}} & \multicolumn{2}{c}{ours, \V{residual}}\\
\cmidrule(lr){3-4}\cmidrule(lr){5-6}\cmidrule(lr){7-8}
set & prover & FA & acc. & FA & acc. & FA & acc.\\
\midrule
ProofNetVerif (150 / 150)
  & tactics & 1 & 73 & 0 & 54 & 0 & 18\\
  & Qwen3.8-27B & 3 & 77 & 0 & 62 & 1 & 20\\
  & GPT-4.1-mini & 1 & 81 & 0 & 58 & 0 & 19\\
  & DeepSeek-V4.1-Flash & 22 & 108 & 0 & 74 & 11 & 34\\
  & Gemini 4 Argon & 34 & 120 & 0 & 72 & 20 & 37\\
  & GPT-6 Astra & 51 & 128 & 0 & 75 & 23 & 40\\
  & Claude Opus 5.5 & \textbf{61} & 130 & 0 & 73 & 19 & 39\\
ProofNet-Verified (75 / 242)
  & tactics & 1 & 235 & 1 & 174 & 0 & 34\\
  & Qwen3.8-27B & 1 & 235 & 1 & 185 & 0 & 40\\
  & GPT-4.1-mini & 1 & 235 & 1 & 185 & 0 & 40\\
  & DeepSeek-V4.1-Flash & 7 & 236 & 2 & 195 & 2 & 40\\
  & Gemini 4 Argon & 12 & 237 & 2 & 196 & 2 & 40\\
  & GPT-6 Astra & 17 & 237 & 4 & 196 & 3 & 41\\
  & Claude Opus 5.5 & \textbf{25} & 239 & 4 & 197 & 4 & 41\\
\bottomrule
\end{tabular}
\caption{Prover configurations on the 300-pair ProofNetVerif subset and on ProofNet-Verified, ordered
by closed proofs found; different models, not budgets of one model. FA: false accepts on unfaithful
pairs; acc.: accepts on faithful pairs. Tactics: no language
model, for the comparison goals and the regrouping certificates alike; each model adds one proposal and one repair round (Qwen3.8-27B without thinking). The four
false accepts on ProofNet-Verified are two kernel-proved label errors and two pointwise
collapses (\S\ref{sec:exp}).}
\label{tab:sweep}
\end{table}

\paragraph{The paradox, measured (Table~\ref{tab:full}).} The measurements match the mechanism of
Proposition~\ref{thm:mono}(ii). Configurations that add few closed proofs to fixed tactics look
excellent: 12 false accepts of 694 with tactics alone and with GPT-4.1-mini, 18 with Qwen. With
the configurations that find many closed proofs, false accepts are 108 with DeepSeek and 273
(39\%) with Opus, while the accepted equivalent pairs are 215 with tactics and 384 with Opus: 23
times the false accepts at 1.8 times the recall. The 300-pair subset and the independently
labelled ProofNet-Verified show the same ordering across seven prover configurations
(Table~\ref{tab:sweep}): with DeepSeek, Gemini, Astra and Opus, closed false accepts are 22, 34,
51 and 61 of 150, and 7, 12, 17 and 25 of 75. Astra's numbers are a lower bound on its proving
power, since the provider's content filter refused 37 of its benign Lean proof requests, counted
as failed proofs. On that subset we have no false accept under any configuration; on
ProofNet-Verified ours are two label errors and two pointwise collapses (\S\ref{sec:exp}). The
configurations differ in model, not only in strength, so the tables show an association with
measured proving power; Proposition~\ref{thm:mono}(ii) supplies the mechanism, and
Table~\ref{tab:budget} varies the budget of one model. Our false accepts also move with the
prover, since a stronger prover certifies more accepts, but they stay at 4--14 of 694
(0.6--2.0\%); \S\ref{sec:survive-exp} dissects them. \V{Residual} tracks closed equivalence one
level down: 1 to 100 false against 38 to 95 faithful. The verdicts are paired, so we compare
them on the same pairs: with Opus, 259 inequivalent pairs are accepted only by closed equivalence
and none only by our criterion (exact McNemar $p<10^{-30}$), at a cost of 137 against 1
equivalent pairs ($p<10^{-30}$). Table~\ref{tab:dist} gives every outcome of this
configuration, grouped by the evidence behind it, and its $\Ecl$ columns show what becomes of the
closed accepts, so that the cost is visible next to what it buys. Of the 137 equivalent pairs
that only closed equivalence accepts, 95 are \V{residual}, 29 \V{unknown} and 13 rejected with a
certificate; of the 259 closed false accepts that we do not accept, 88 are \V{residual}, 83
\V{unknown} and 88 rejects that locate the difference, 54 of them on a total alignment. In
coverage, we decide 64\% of the inequivalent and 88\% of the equivalent pairs, counting
\V{residual} as decided, and 49\% and 67\% without it. Of the 252 \V{unknown} inequivalent pairs only 75 have a full alignment with every
comparison goal open; the other 177 fail earlier, at canonicalization, a regrouping certificate
or the alignment, so most of the coverage gap is in the front end, not in proof search. A stronger prover mostly turns \V{unknown} into certified rejects, from 5 to 328 on
inequivalent pairs. The certificates also say where in the aligned telescopes the difference
is, which no closed verdict does: of these 328, a certificate shows that the hypotheses differ in 146, the conclusions
inside the hypotheses in 59, and both in 39. The other 84 certify a difference of the closed
statements without locating it: for 54 no component goal is decided, and the reject is a refuted
closed equivalence together with the two regrouping certificates (Proposition~\ref{prop:cert});
for 30 neither telescope has a binder or a hypothesis in its prefix, so the conclusion goal
compares the closed statements. The location rests on the alignment layer of
Table~\ref{tab:layers}: 104 of the 244 located rejects have an alignment that is not total, and
in the 82 of them where closed equivalence is not refuted, the component that differs is only
as meaningful as that alignment. Even on a total alignment a location is mechanical: the labels
say that these pairs are inequivalent, not where, and no reader has checked the 244 locations.
The blind audit of \S\ref{sec:blind} checks the located rejects of labelled-equivalent pairs, and
where an annotator names the kind of difference, the located component matches it in 81--91\% of
the cases. Closed equivalence proves 65 of the 146 (45\%) and 16 of the
39, but only 7 of the 59 (12\%), and it refutes 27 of the 59. This is the pattern that Proposition~\ref{thm:char}
predicts: a candidate that changes only its hypotheses often remains a true theorem, and closed
equivalence relates any two true theorems, while a candidate that changes its conclusion inside
the hypotheses is often false, which a closed refutation can detect. With DeepSeek the closed
proofs are 19 of 104, 2 of 8 and 0 of 31. On equivalent pairs the strongest prover rejects 51 of 446, each with a kernel
certificate of a difference. Two true statements are pointwise equivalent, so the direction goals
sort these rejects by what they certify (Proposition~\ref{prop:split}), into $30+8+13$. In 30
of them one of the two statements is certified false: in $15+13$ of these also which one, in
the other 2 not. For 15, closed equivalence is refuted and the candidate is the
stronger, false side, so the label is wrong; 8 of them are rejects only through this
refutation: the candidate quantifies over every group (Dummit--Foote 4.4.6b, 3 pairs), admits
$n=0$ (Ireland--Rosen 3.10, 2 pairs), writes the union of subspaces as their sum (Artin 3.7.2),
divides in $\mathbb N$ (Ireland--Rosen 2.27a) or drops the order topology (Munkres 18.8a). For 13, \texttt{npimp} and \texttt{rs} certify that
the \emph{reference} is false. They come from two ProofNet references, both unchanged in
ProofNet\#: Dummit--Foote 2.1.13 (8 pairs) states the hypothesis for a single $x$, so the subgroup
$\mathbb Z\le\mathbb Q$ refutes it, and Pugh 2.92 (5 pairs) drops the Hausdorff assumption under
which nested nonempty compact sets meet. For the other 2, \texttt{npimp} certifies that one side
is false but no direction is proved, so which side is open. 8 rejects are vacuity-only, among them the two canonicalization artefacts of
Appendix~\ref{app:rejects} and a third that the blind audit finds (Putnam 2018 B4,
\S\ref{sec:blind}); these are the only rejects that depend on the split. The remaining 13
have a proved direction but no \texttt{npimp}, so whether they are vacuity-only is open. The 51
come from only 22 problems, and 22 of them locate the difference on an alignment that is not
total, where the positional prefix can pair binders of different meaning. Deciding which of the open ones a reader would waive needs
a reader; the 43 decided by component goals are all in the blind audit, which judges 16 of the 26
statements it agrees on unfaithful and finds three canonicalization artefacts (\S\ref{sec:blind}).

\paragraph{An external adjudication.} The authors of SemantEq re-adjudicated the
ProofNetVerif test pairs labelled equivalent, by hand after a language-model screen, under the
criterion that the candidate makes the same claim as the reference~\cite{semanteq}; 444 of our
446 equivalent pairs are among them. They relabel 46 of our 51 rejects as non-equivalent: all 8
rejects through a refuted closed equivalence and 38 of the 43 that rest on component goals. Of
the 5 they keep equivalent, 4 are vacuity-only, among them both canonicalization artefacts of
Appendix~\ref{app:rejects} (1135 and 358; their proof audit certifies 358 in both directions, and
none of our other rejects), and 1 has a proved direction only. Of the 35 component rejects that
are not vacuity-only, 34 are relabelled, so the disagreements fall where
Proposition~\ref{prop:split} says that a reject depends on the split. Of our 248 accepts they
relabel 3, all Artin 2.8.6, which they attribute to their criterion (an equality of centres
against an isomorphism). The other way round, of the 96 pairs they relabel that we cover, we
reject 46, accept 3 and leave 12 \V{residual} and 35 \V{unknown}. The adjudication is
independent of our verdicts but not of theirs, since most of its reviews were prompted by pairs
that their own method did not accept, and it covers only pairs labelled equivalent, so it says
nothing about our false accepts (\texttt{k17\_semanteq.py}).

\paragraph{Combining with SemantEq.} The two criteria give evidence on opposite sides, so they
compose. Table~\ref{tab:common} scores them, BEq+ as SemantEq released it, and two combinations
under ProofNetVerif's labels and under SemantEq's adjudicated ones. On the 1,135 test pairs that both cover, SemantEq-Eval accepts 235 and never a pair
labelled inequivalent, and it decides nothing else; our criterion decides 640 (262 accepts, 14 of
them false, and 378 certified rejects). The two-sided rule takes SemantEq's accepts as accepts and
our certified rejects as rejects, leaves a pair with both undecided and does not use our accepts. It
decides 611 pairs (54\%, against 21\% for SemantEq alone) with no false accept observed under either
label set (at most 1.6\% of its 234 accepts at 95\% confidence); of its 377
rejects, 50 are of pairs labelled equivalent by ProofNetVerif and 4 under SemantEq's adjudicated
labels (98.9\% precision). The two disagree on one pair, the artefact 358. Our accepts are left
out because all our errors lie outside SemantEq's accepts: of its 235 accepts we accept 191, leave 40
\V{residual} and 3 \V{unknown} and reject 1, while SemantEq accepts none of
our 14 false accepts, and our accepts add 71 pairs that SemantEq does not accept, 57 of them
labelled equivalent. Accepting what either accepts would take on all 14; accepting only what both
accept gives up 43 accepts and no false one. The rule follows from what each side certifies, but we
fixed it after seeing SemantEq's verdicts, so this comparison is exploratory (\texttt{k18\_combine.py}).

\begin{table}[t]
\centering\small
\begin{tabular}{lrrrrrrrr}
\toprule
& \multicolumn{4}{c}{accepts} & \multicolumn{3}{c}{rejects} & \\
\cmidrule(lr){2-5}\cmidrule(lr){6-8}
check & $n$ & false & prec. & recall & $n$ & false & prec. & decided\\
\midrule
\multicolumn{9}{l}{\emph{ProofNetVerif labels: 444 equivalent pairs, 691 not}}\\
BEq+ & 233 & 8 & 96.6 & 50.7 & --- & --- & --- & 20.5\\
SemantEq-Eval & 235 & 0 & 100.0 & 52.9 & --- & --- & --- & 20.7\\
ours & 262 & 14 & 94.7 & 55.9 & 378 & 51 & 86.5 & 56.4\\
two-sided & 234 & 0 & 100.0 & 52.7 & 377 & 50 & 86.7 & 53.8\\
both accept & 191 & 0 & 100.0 & 43.0 & 378 & 51 & 86.5 & 50.1\\
\midrule
\multicolumn{9}{l}{\emph{SemantEq's adjudicated labels: 348 equivalent pairs, 787 not}}\\
BEq+ & 233 & 10 & 95.7 & 64.1 & --- & --- & --- & 20.5\\
SemantEq-Eval & 235 & 0 & 100.0 & 67.5 & --- & --- & --- & 20.7\\
ours & 262 & 17 & 93.5 & 70.4 & 378 & 5 & 98.7 & 56.4\\
two-sided & 234 & 0 & 100.0 & 67.2 & 377 & 4 & 98.9 & 53.8\\
both accept & 191 & 0 & 100.0 & 54.9 & 378 & 5 & 98.7 & 50.1\\
\bottomrule
\end{tabular}
\caption{Reference checks on the 1,135 ProofNetVerif test pairs that SemantEq and we both cover,
under two label sets. False: accepts of pairs labelled not equivalent, rejects of pairs labelled
equivalent; precision, recall (share of the equivalent pairs accepted) and decided (share of all
pairs accepted or rejected) in \%. BEq+ and SemantEq-Eval are the verdicts SemantEq released; ours
is Claude Opus~5.5 at the budget of Table~\ref{tab:full}, with certified rejects only. Two-sided:
SemantEq's accepts and our rejects, the one pair with both left undecided; both accept: our rejects,
and accepts only where both accept. Provers and models differ, so this compares outcomes, not
methods at equal cost.}
\label{tab:common}
\end{table}

\begin{table}[t]
\centering\small
\begin{tabular}{llrrrrrrr}
\toprule
& & \multicolumn{2}{c}{closed $\Ecl$} & \multicolumn{2}{c}{ours, \V{accept}} & \multicolumn{2}{c}{ours, \V{residual}} & \\
\cmidrule(lr){3-4}\cmidrule(lr){5-6}\cmidrule(lr){7-8}
model & budget & FA & acc & FA & acc & FA & acc & reject U\\
\midrule
Claude Opus 5.5 & tactics & 12 & 215 & 4 & 190 & 1 & 38 & 7\\
 & + proposal & 218 & 366 & 13 & 241 & 69 & 90 & 272\\
 & + repair & 274 & 384 & 14 & 248 & 100 & 95 & 328\\
 & + 2nd repair & \textbf{305} & 391 & 15 & 253 & 104 & 95 & 357\\
DeepSeek-V4.1-Flash & tactics & 12 & 215 & 4 & 190 & 1 & 38 & 7\\
 & + proposal & 62 & 288 & 6 & 243 & 19 & 65 & 89\\
 & + repair & 108 & 328 & 9 & 249 & 41 & 82 & 168\\
 & + 2nd repair & \textbf{151} & 342 & 10 & 253 & 46 & 85 & 195\\
\bottomrule
\end{tabular}
\caption{One model, nested budgets, all 1,140 pairs. Each level adds to the one above it: one
proposal per open goal, one repair round (the budget of Table~\ref{tab:full}), a second repair
round. Columns as in Table~\ref{tab:full}. The levels vary only the comparison goals; the
regrouping certificates are those of the full run at every level, which is why the tactics level
has 2 more faithful accepts and 2 more rejects here than the tactics row of Table~\ref{tab:full}, where the
certificates must be found by tactics too.}
\label{tab:budget}
\end{table}

\paragraph{One model, growing budget (Table~\ref{tab:budget}).} To separate budget from model we
replay the full-set run of a model level by level with the prompts, goals and caches of
Table~\ref{tab:full}, and add a second repair round with new calls. Each level keeps the proofs of
the previous one, so the provers are nested and Proposition~\ref{thm:mono}(ii) already rules out a
decrease; the table measures the size of the increase. With Opus, closed false accepts go from 12
to 218, 274 and 305 of 694 (44\%). The second repair round adds 31 closed false accepts against 7
closed accepts of equivalent pairs. Our false accepts go from 4 to 15, and \V{residual} ones,
which follow the closed collapse on the unaligned remainder, from 1 to 104. Replaying the repair level reproduces the stored verdicts on all but 4 of the 1,140 Opus
pairs, where the Lean re-check of the same cached proofs came out differently; one of the four
is the extra closed false accept (274 here, 273 in Table~\ref{tab:full}). DeepSeek, which proves
less, shows the same shape: closed false accepts 12, 62, 108 and 151, ours 4, 6, 9 and 10, and
its replay reproduces every stored verdict. GPT-4.1-mini and Qwen
add almost nothing to the tactics at any level (Table~\ref{tab:full}), so the collapse needs a
model that actually proves.

\paragraph{What each level costs (Table~\ref{tab:cost}).} The levels are matched in rounds, not in
cost, so we also count what each side spends: model calls, prompt plus completion tokens read
from the cached responses, and Lean checks, each under the same heartbeat limit (the runner does
not time single checks). The two sides spend differently by construction: the closed check asks
for each of its four goals in a separate call and repairs every failed one, while the component
check sends the open goals of a pair in one call and stops once a verdict is decided. With both
models, the component check at its largest budget costs less than the closed check at its smallest
model budget, in calls, tokens and Lean checks alike: with Opus 1{,}747 calls and 6.5M tokens
against 3{,}914 and 8.1M, with 15 against 218 false accepts; with DeepSeek 2{,}000 calls and 41.5M
tokens against 3{,}930 and 75.4M, with 10 against 62. With no model calls at all both sides run
fixed tactics, and the false accepts are 4 against 12 (Table~\ref{tab:budget}). Since the closed
false accepts can only grow with the budget (Proposition~\ref{thm:mono}), the difference in
precision is not bought with budget. The difference in coverage is not removed by budget either:
the closed check at 8.1M tokens accepts 366 equivalent pairs, ours at 6.5M accepts 253 and returns
\V{residual} on 95 more. The replay reads every proof and every response from the caches of
Table~\ref{tab:budget} and reproduces its verdict counts exactly (\texttt{k15\_cost.py}).

\begin{table}[t]
\centering\small
\begin{tabular}{llrrrrrrrrrr}
\toprule
& & \multicolumn{5}{c}{closed $\Ecl$} & \multicolumn{5}{c}{ours, \V{accept}}\\
\cmidrule(lr){3-7}\cmidrule(lr){8-12}
model & budget & calls & tokens & Lean & FA & acc & calls & tokens & Lean & FA & acc\\
\midrule
Claude Opus 5.5 & + proposal & 3914 & 8.1M & 3651 & 218 & 366 & 1127 & 4.2M & 1142 & 13 & 241\\
 & + repair & 6102 & 12.4M & 5653 & 274 & 384 & 1535 & 5.7M & 1493 & 14 & 248\\
 & + 2nd repair & 7840 & 15.5M & 7007 & 305 & 391 & 1747 & 6.5M & 1702 & 15 & 253\\
DeepSeek-V4.1-Flash & + proposal & 3930 & 75.4M & 3311 & 62 & 288 & 1128 & 23.6M & 838 & 6 & 243\\
 & + repair & 6608 & 101.2M & 4318 & 108 & 328 & 1693 & 34.8M & 1344 & 9 & 249\\
 & + 2nd repair & 7192 & 112.2M & 4822 & 151 & 342 & 2000 & 41.5M & 1588 & 10 & 253\\
\bottomrule
\end{tabular}
\caption{Cumulative cost of the levels of Table~\ref{tab:budget} on all 1,140 pairs: model calls,
prompt plus completion tokens (DeepSeek's include its reasoning tokens), and Lean checks after
the tactic level. FA and acc as in Table~\ref{tab:budget}. Calls that failed or were refused are
not counted.}
\label{tab:cost}
\end{table}

\paragraph{Resampling the prover.} Decided verdicts are kernel proofs, so a second run of the same
prover can change coverage but cannot contradict the first. We reran Opus on the 300-pair subset
with every proposal and repair drawn afresh. Of the 300 verdicts, 59 changed, all between a
decided verdict and \V{unknown} (36 newly decided, 23 no longer decided). No accept became a
reject or the reverse, and no closed proof became a closed refutation. Closed false accepts were
61 and 58 in the two runs (66 in their union), ours 0 in both, \V{residual} ones 19 and 23 (26). Both
checks vary from run to run in what they decide; the meaning of a decided verdict does not.

%======================================================================
\section{Evaluating candidate formalizations}\label{sec:exp}

\begin{table}[t]
\centering\small
\begin{tabular}{lrrrrr}
\toprule
& \multicolumn{2}{c}{inequivalent (150)} & \multicolumn{3}{c}{equivalent (150)}\\
\cmidrule(lr){2-3}\cmidrule(lr){4-6}
method & accept (false) & reject & accept & reject, total & of which false\\
\midrule
closed $\Ecl$, weak prover & 22 (15\%) & -- & 108 (72\%) & -- & --\\
closed $\Ecl$, strong prover & \textbf{61 (41\%)} & -- & 130 (87\%) & -- & --\\
BEq+ & 0 & -- & 70 (47\%) & -- & --\\
ours, weak prover & \textbf{0} & 33 (22\%) & 74 (49\%) & 3 & 2 (1.3\%)\\
ours, strong prover & \textbf{0} & 69 (46\%) & 73 (49\%) & 17 & 2 (1.3\%)\\
\quad \V{residual}, weak & 11 (7\%) & -- & 34 (23\%) & -- & --\\
\quad \V{residual}, strong & 17 (11\%) & -- & 39 (26\%) & -- & --\\
\bottomrule
\end{tabular}
\caption{ProofNetVerif, 300 pairs. Weak prover: DeepSeek-V4.1-Flash; strong: Claude Opus 5.5.
``Of which false'': rejects caused by canonicalization artefacts (Appendix~\ref{app:rejects}).}
\label{tab:g1}
\end{table}

\paragraph{The 300-pair subset (Table~\ref{tab:g1}).} The stratified subset is where we compare
with BEq+, judges and the reject analysis in detail. With the strong prover closed equivalence
false-accepts 61 of 150 inequivalent pairs; BEq+ has none and pays with 47\% recall. We have
no false accept under either prover (at most 2.5\% of the 150 at 95\% confidence) and accept
73--74 equivalent pairs; another 11--17
inequivalent and 34--39 equivalent pairs are \V{residual}. Of the rejects on labelled-equivalent
pairs, 2 under either prover are canonicalization artefacts. The others involve a false statement
(strong: 11, weak: 0; in four of them the false side is the reference, Pugh 2.92 or Dummit--Foote
2.1.13, and three are rejects only through a refuted closed equivalence, so the two statements
differ in truth value and the label is wrong), or are certified hypothesis differences that the
annotator tolerated (strong: 4, weak: 1), so counting those as false too gives 4.0\% (strong) and
2.0\% (weak). The strong run also assigns 2 inequivalent pairs \V{off-image} (accepts before
embedding alignment, counted as accepts in Table~\ref{tab:sweep}) and leaves 62 inequivalent and
21 equivalent pairs \V{unknown}; 57 of these 83 lack a regrouping certificate or a full
alignment, and only 26 are fully aligned with every component goal open
(Appendix~\ref{app:unk}). Our accepts have no false accept on this subset under any of the
seven configurations we ran on it. On the full 1,140-pair set they have 4--14
(Table~\ref{tab:full}): the zero is a property of the subset, not of the criterion, and
\S\ref{sec:survive-exp} explains the full-set cases. With the strong prover, 54 of the 73
accepts on equivalent pairs are informative (Proposition~\ref{prop:informative}), 16 are
pointwise collapses and 3 are unclassified; of the 19 \V{residual} verdicts on inequivalent
pairs (17 plus the 2 later turned \V{off-image}), 17 are collapses of the remainder and 2 have an informative remainder (pairs 796 and
1355 of the test split). Of the strong run's rejects, 21 decide no component goal: their
closed equivalence is refuted and both regrouping certificates hold, which by
Proposition~\ref{prop:cert} refutes $\Econd$ of the telescopes (18 inequivalent pairs and the 3
label errors above). Embedding alignment converts two \V{residual} verdicts
(Fermat's little theorem over $\mathbb Z$ vs.\ $\mathbb N$) into \V{off-image}. An audit of 20
inequivalent pairs left \V{unknown} found 19 labels correct and 1 disputed.


\paragraph{Regrouping certificates.} Each verdict in Table~\ref{tab:g1}
is a proof of a goal about the canonical components, together with the two regrouping
certificates that bind the closures of the telescopes to the original elaborated statements
(\S\ref{sec:framework}). With the fixed tactics alone, 37 regrouping goals of the 194 decided pairs (strong prover)
were missing; a second tactic tier that instantiates the other side directly proved 30 of
them and one proposal round the remaining 7, so all 194 decided verdicts, including all
63 rejects, are certified end to end. On the 286 decided ProofNet-Verified pairs
(Table~\ref{tab:g4}) 66 regrouping goals were missing; the second tactic tier proved 40 and
one proposal round the remaining 26, so all 286 decided verdicts there, including all 40
rejects, are certified end to end as well.
The direction goals of Proposition~\ref{prop:vac} split the 63 strong-prover rejects on
ProofNetVerif into 16 vacuity-only rejects (the pointwise claims agree, so a reader may waive
the difference), 13 weaker and 26 stronger candidates, 4 that differ in both directions and 4
left undecided; the 40 rejects on ProofNet-Verified split into 5, 3, 27, 3 and 2.

\begin{table}[t]
\centering\small
\begin{tabular}{lrrrrrr}
\toprule
& \multicolumn{3}{c}{unfaithful, candidate true (24)} & \multicolumn{2}{c}{unfaithful, other (51)} & faithful (242)\\
\cmidrule(lr){2-4}\cmidrule(lr){5-6}\cmidrule(lr){7-7}
method & accept & reject & unknown & accept & reject & accept\\
\midrule
closed $\Ecl$, strong prover & \textbf{21} & -- & 3 & 4 & 21 & 239\\
ours, strong prover & 3 & 7 & 11 & 1 & 37 & 197\\
\quad \V{residual} & 3 & -- & -- & 1 & -- & 41\\
\bottomrule
\end{tabular}
\caption{ProofNet-Verified: 317 pairs of a ProofNet\# statement and its manually corrected
version, with the relation label and the truth label of the candidate taken from the audit.
``Candidate true'' is the audit's provability label; ``other'' is unfaithful with the candidate
not marked true. 227 of the 242 faithful pairs are textually identical. Our 8 proofs of
$\Ecomp$ on unfaithful pairs (4 accepts, 4 \V{residual}) split by cause: 3 are kernel-proved
equivalences (the labels are wrong), 2 are pointwise collapses (both conclusions hold at every
aligned input) and 3 are collapses of an unaligned remainder. The closed criterion also refutes two
pairs labelled faithful: the candidate is false (Munkres 21.6b and 24.3a).}
\label{tab:g4}
\end{table}

\paragraph{Independent labels (Table~\ref{tab:g4}).} The labels in Table~\ref{tab:g1} are
ProofNetVerif's, but the classification of rejects on labelled-equivalent pairs is ours. ProofNet-Verified~\cite{pnv}
provides a second, independent setting: each ProofNet\# statement is paired with a manually
corrected version, and the audit records the relation of the original to the informal
statement (faithful, weaker, stronger, incomparable) and whether the original is provable.
We run the same pipeline, with the same strong prover, on the 317 pairs whose two statements
elaborate (13 fail to convert, 28 to elaborate). The 24 pairs whose candidate is unfaithful
but true are exactly the pairs Proposition~\ref{thm:mono}(ii) is about: closed equivalence accepts
21 of them and leaves the other 3 \V{unknown} only because the prover failed; it also
accepts 4 of the 51 remaining unfaithful pairs. Our criterion proves $\Ecomp$ for 8 unfaithful
pairs in total, 4 as \V{accept} and 4 as \V{residual}. Three of these are kernel-proved equivalences, so the label is wrong: one pair is
textually identical (Dummit--Foote 2.4.4), one differs by a type equation that the one-point rule
eliminates (Dummit--Foote 8.3.6a), and in one the ``stronger'' candidate \texttt{SeparatedNhds A B}
is provably equivalent to the reference's $A\cap\overline B=\emptyset\wedge\overline A\cap B=\emptyset$
for all sets in a metric space (Rudin 2.19a; metric spaces are completely normal). Two are
pointwise collapses (Proposition~\ref{prop:residue}): the alignment is total, the hypotheses
agree, and both conclusions hold at every aligned input. In Artin 2.8.6 these are
$Z(G\times H)=Z(G)\times Z(H)$ and the existence of an isomorphism between them; in Pugh 4.15a
both are ``$f$ is uniformly continuous iff it has a modulus of continuity'', once with and once
without continuity of the modulus, and both versions are theorems (the prover constructs the
modulus). The remaining three are \V{residual}: a binder without counterpart stops the alignment
and the remainder is compared closed (the reference's subgroup in Dummit--Foote 4.3.26, an inserted
\texttt{[Fintype $\alpha$]} in Herstein 3.2.21, the abstract normed space against
\texttt{EuclideanSpace} in Pugh 2.41). None is a pair separated by a hypothesis or a conclusion that the criterion
compared and got wrong. The two faithful pairs the closed criterion refutes are genuine
errors that the audit missed: Munkres 21.6b and 24.3a, where \texttt{Icc 0 1} is elaborated in
$\mathbb N$ (in 21.6b the sequence $x^n$ is then eventually constant and does converge uniformly;
\texttt{BenchErrors.lean} proves the negation of both ProofNet\# statements). Rejects on
the faithful side: these two, by their refuted closed equivalence (Proposition~\ref{prop:cert});
4 of the rejects of unfaithful pairs also come from a closed refutation alone. Of the 246 proofs of $\Ecomp$, 201 are accepts and cover 197 of the
242 faithful pairs; the four on unfaithful labels are two of the kernel-proved label errors
(Dummit--Foote 2.4.4, Rudin 2.19a) and the two pointwise collapses. The other provers give
185--195 accepts on faithful pairs, with 1--2 on unfaithful labels, all among these four.
Of the 197 faithful accepts, 124 are informative (Proposition~\ref{prop:informative}),
41 pointwise collapses and 32 unclassified; on unfaithful labels the only informative
accept is Rudin 2.19a, a kernel-proved label error. Dummit--Foote 2.4.4 is a collapse as well
as a label error, since its conclusion holds everywhere.
BEq+ does well in this setting, because 227 of the 242 faithful pairs are textually identical:
it accepts 234 faithful pairs (97\%, Wilson 95\% interval 94--99\%) with a single false accept,
the textually identical label error Dummit--Foote 2.4.4. As a binary accept it beats ours here
(197 faithful, 4 unfaithful); counting \V{residual} as well would reach 238 faithful pairs, a
superset of BEq+'s ($p=0.125$, exact McNemar), but with 8 unfaithful. Accepting a pair when we
accept it or BEq+ does gives 237 faithful and the same 4 unfaithful. What BEq+ does not give is a
reason when it fails: its 8 non-accepts on faithful pairs and 74 on unfaithful pairs carry no
certificate, while 40 of our verdicts on this set are certified rejects.

\begin{table}[t]
\centering\small
\begin{tabular}{llrrrrrrr}
\toprule
& & \multicolumn{3}{c}{ProofNetVerif} & \multicolumn{3}{c}{ProofNet-Verified} & \\
\cmidrule(lr){3-5}\cmidrule(lr){6-8}
judge & input & FA/150 & acc/150 & veto & FA/75 & acc/242 & veto & wrong\\
\midrule
Opus 5.5 & NL & 20 & 124 & 14 / 117 & 11 & 240 & 7 / 240 & 11 / 4\\
Opus 5.5 & NL+ref & 16 & 128 & 13 / 120 & 5 & 238 & 3 / 238 & 6 / 3\\
V4.1-Flash & NL & 14 & 122 & 10 / 116 & 15 & 238 & 7 / 238 & 8 / 12\\
V4.1-Flash & NL+ref & 10 & 115 & 8 / 111 & 2 & 237 & 2 / 237 & 8 / 3\\
Gemini 4 Argon & NL & 13 & 121 & 10 / 115 & 13 & 238 & 5 / 238 & 9 / 10\\
Gemini 4 Argon & NL+ref & 12 & 122 & 9 / 117 & 2 & 239 & 1 / 239 & 7 / 2\\
GPT-4.1-mini & NL & 37 & 141 & 28 / 130 & 46 & 239 & 21 / 239 & 17 / 34\\
GPT-4.1-mini & NL+ref & 20 & 129 & 17 / 122 & 16 & 238 & 9 / 238 & 13 / 13\\
\midrule
ours (Opus) & ref & 0 & 73 & & 4 & 197 & & \\
\quad with \V{residual} & ref & 17 & 112 & & 8 & 238 & & \\
\bottomrule
\end{tabular}
\caption{LLM judges (temperature 0, one sample). FA: accepts on unfaithful pairs; acc: accepts on
faithful pairs. Veto: FA / acc after every judge accept that we \V{reject} is overruled. Wrong:
judge verdicts contradicted by both a kernel certificate and the label (ProofNetVerif /
ProofNet-Verified).}
\label{tab:judge}
\end{table}

\paragraph{LLM judges (Table~\ref{tab:judge}).} Judges are the practical alternative, and
with a reference they are competitive on rates: Opus 5.5 accepts 16 of 150 unfaithful
ProofNetVerif pairs at 128 of 150 faithful ones, and on ProofNet-Verified 5 of 75 at 238 of
242. Our accepts trade recall for precision: 0/73 and 4/197 on these two sets, two of the four
being label errors and two pointwise collapses between true conclusions, and 14/248 on the full
set (\S\ref{sec:survive-exp}); counting \V{residual} as accepts would give 17/112 and 8/238, close
to the judges. Three properties separate the two. (i) \emph{Stability.} Across the eight
judge configurations Fleiss' $\kappa$ is 0.81 and 0.77; 84 of 300 and 63 of 317 pairs receive
both verdicts, and adding the reference flips 8--45 verdicts of the same model. (ii)
\emph{Checkable errors.} 39 and 41 pairs receive at least one judge verdict that is
contradicted by a kernel certificate agreeing with the label (an accept of a labelled-unfaithful
pair that we reject or whose closed equivalence is refuted, or a reject of a labelled-faithful
pair that we prove component-equivalent); for the strongest judge 6 and 3.
Our verdicts cannot be contradicted this way, and a disputed one can be re-checked by the
kernel. (iii) \emph{Composition.} Because our \V{reject} is a certificate, it can veto a
judge's accept without trusting either side. On ProofNet-Verified, where we reject no faithful
pair, the veto lowers the false accepts of the eight configurations together from 110 to 55
(by 0 to 25 per configuration) at no cost in recall; on ProofNetVerif it removes 2--9 per
configuration and costs 4--11 faithful accepts, all of them among the 14 rejects of
Appendix~\ref{app:rejects} (label errors, tolerated hypothesis differences and two artefacts).

\paragraph{Further checks (Appendix~\ref{app:more}).} On 386 mutants with injected drift of
known class no verdict contradicts the class; statements in decidable fragments receive no
\V{unknown}; 292 of 299 surface re-presentations are proved component-equivalent to the original;
and on the FormalRx-CC benchmark the goals certify 83\% of relation labels with no language
model, against 10\% Lean-passing witnesses for the original protocol.

%======================================================================
\section{Where accepts still collapse}\label{sec:survive-exp}

\paragraph{Anatomy on the full set.} With Opus, 114 of the 694 inequivalent pairs have a proof
of $\Ecomp$ (Table~\ref{tab:full}). 100 of them are \V{residual}: the alignment ends early and
the remainder is compared closed, and for 91 of these the remainder is certified constant on the
aligned inputs, the situation of Proposition~\ref{prop:residue}. This is why \V{residual} is not
an accept. Of the 14 accepts, 7 have a conclusion that is certified
constant on the aligned inputs; 4 are absorptions in the sense of \S\ref{sec:survive}, all
variants of one exercise (Artin 11.2.13, ``$a\mid b$ in $\mathbb Z[i]$ implies $a\mid b$ in $\mathbb Z$''
against ``$a\mid b$ implies $a\mid b$''); one has hypotheses that are equivalent relative to $\T$ but
not as the textbook intends them (Herstein 2.2.3), which no $\T$-relative check can see; and two
looked to us like label errors (Putnam 2018 A5, Artin 10.6.7). The blind audit of
\S\ref{sec:blind} confirms 9 of the 14 as unfaithful and judges the other 5 faithful: three Pugh
2.26 candidates with a constant conclusion, Artin 10.6.7, and Putnam 2018 A5 (whose two sides
misstate the exercise in the same way, \S\ref{sec:audit}). The one-sided certificates of
\S\ref{sec:survive} target the constant and absorbed ones. The triviality certificate flags 5 of
the 14 false accepts, including all four absorptions, and none of the 248 faithful accepts (nor
any faithful \V{residual}), nor any faithful accept on the 300-pair subset or on
ProofNet-Verified. The constancy certificate costs
more: 56 of the 248 faithful accepts also have a constant conclusion, so treating constancy
as disqualifying removes correct accepts that certify only the hypotheses.

\begin{table}[t]
\centering\small
\begin{tabular}{lrrr}
\toprule
accept rule (Opus, all 1,140 pairs) & FA / 694 & FA, 95\% CI & acc / 446\\
\midrule
closed $\Ecl$ & 273 & 35.8--43.0\% & 384\\
ours, \V{accept} or \V{residual} & 114 & 13.9--19.4\% & 343\\
ours, \V{accept} & 14 & 1.2--3.4\% & 248\\
ours, not trivial & 9 & 0.7--2.5\% & 248\\
ours, not trivial, not constant & 3 & 0.1--1.3\% & 192\\
BEq+ & 2 & 0.1--1.0\% & 207\\
BEq+ or ours, not trivial & 11 & 0.9--2.8\% & 288\\
BEq+ or ours, not trivial, not constant & 5 & 0.3--1.7\% & 268\\
\bottomrule
\end{tabular}
\caption{Accept rules of increasing strictness on the full set (Wilson intervals). ``Not trivial''
and ``not constant'': the triviality and constancy certificates of \S\ref{sec:survive} do not
fire.}
\label{tab:levels}
\end{table}

\paragraph{Accept rules (Table~\ref{tab:levels}).} The two filters are different checks, one on the
candidate's propositional form and one on the conclusion alone, and they do not cost the same.
Against BEq+, our accept recalls 41 more equivalent pairs (248 against 207, 81 accepted only by
us and 40 only by BEq+, exact McNemar $p<10^{-3}$) at 12 more false accepts (14 against 2,
$p=0.004$). Triviality lost no faithful accept on any set we ran and takes the false accepts from
14 to 9; with only 5 discordant pairs this is not significant on its own (exact McNemar $p=0.06$).
Constancy takes them from 9 to 3 ($p=0.03$) and costs 56 faithful accepts (248 to 192,
$p<10^{-16}$), because a constant conclusion is also common in faithful pairs whose content lies
in the hypotheses. Read as a binary classifier, the strictest rule is indistinguishable from BEq+
on these data: 3 against 2 false accepts ($p=1$) and 192 against 207 faithful accepts, 61 of them
accepted only by us and 76 only by BEq+ ($p=0.23$). The two are complementary: accepting when
either accepts gives 268 of 446 equivalent pairs (60\%, 61 more than BEq+, $p<10^{-17}$) at 5
false accepts (0.7\%, $p=0.25$ against BEq+'s 2), and dropping the constancy condition gives 288
(65\%) at 11 (1.6\%). We report all rules, because which one is appropriate depends on whether a false accept or
a lost accept is more costly; every accept in every rule comes with its certificates.

\paragraph{Reward hacking by construction (Table~\ref{tab:hack}).} For 80 ProofNetVerif references
$R=(\Gamma,H,C)$ a Lean metaprogram builds three candidates that a reader rejects at once but that
are closed-equivalent to $R$ whenever $R$ is true: the constant $(\Gamma,H,\top)$, the absorption
$(\Gamma,H\wedge C,C)$ and the vacuous $(\Gamma,H\wedge\bot,C)$. For BEq+ the same edits are applied to
the source statement. Every check scores every candidate, with $R$ itself as the control.
Closed equivalence rewards a hack exactly when its prover proves the reference: 3 of 80 with
tactics or GPT-4.1-mini, 53--56 with Opus. The restriction of BEq+ blocks the constant (2) and the
vacuous candidate (0) but not absorption (54 of 80): applying the candidate leaves the reference
itself as the remaining goal, which \texttt{exact?} then proves from Mathlib, and the proof mentions
the candidate as the rule requires. Our criterion alone rejects the vacuous candidate with a
certificate (73 of 80), but accepts 45 absorptions and 11 constant candidates, as
\S\ref{sec:survive} predicts: the 11 are references whose conclusion holds on the whole aligned
context (8 more, and 15 controls, are \V{residual} because the alignment ends early), and absorption is invisible to every check invariant under equivalent hypotheses. The
\V{unknown} verdicts are expected as well. For a true reference the constant candidate agrees
with it wherever $H$ holds, so its only possible verdicts are \V{accept} and \V{off-domain}, and
the 54 \V{unknown} are mostly references for which the prover found no point outside $H$ where
the conclusion fails; an absorption is decided only when the prover proves or refutes
$H\to C$ pointwise, which it did not for 24. A further 3 absorption accepts, 3 vacuous rejects
and 3 control accepts lack a regrouping certificate and count as \V{unknown}. With GPT-4.1-mini
our criterion accepts 4 absorptions and no other hack.

The triviality flag is a heuristic, not a verdict. It fires on all 240 hacks and withholds their
reward but leaves the verdict, since it shows that the candidate is degenerate, not that it
differs from $R$; ``no hack is rewarded'' therefore means that no hack is accepted unflagged, not
that any hack is rejected. It does not flag the controls, by construction, nor any faithful
accept on real candidates (above). It only recognizes candidates that close by propositional
logic, so we rewrote the constant and the absorbed candidates to need one \texttt{simp} step
(last two rows). The flag then misses all 80 constants and 11 of the 80 absorptions; closed
equivalence rewards 52 and 51 of the rewritten candidates, BEq+ 0 and 9. With Opus we accept 6
of the missed absorptions unflagged, that is, reward them. The 21 rewritten constants that the
prover settles are \V{residual}, not accepts, but only because the new conclusion
$\forall k,\ 0+k=k$ adds a binder without counterpart; a rewrite without a binder would be
accepted unflagged. Our verdicts otherwise hardly move (46 absorption accepts against 45), as
they should, since the rewrites are equivalent. The flag is a floor, not a defence. As a side
effect, 8 of the 80 references are certified false, since a constant or
absorbed candidate can be rejected only where $R$ fails inside its hypotheses; among them are
Rudin~3.5 (Appendix~\ref{app:refs}) and Rudin~1.8, which denies that $\mathbb C$ carries any linear
order where the textbook means an ordered field.

\paragraph{Reward hacking by a language model.} Constructed hacks test known failure modes; a
model told to game the checkers may find others. For 50 of the references we asked Opus for three
Lean statements that do not formalize the problem but that an equivalence checker would accept,
and scored all 150 (they all elaborate). Closed equivalence with Opus rewards 78 of them and BEq+
none. We accept 8, reject 49 with certificates (14 of them rewarded by closed equivalence), and
return \V{residual} for 29 and \V{unknown} for 64. All 8 accepts are the collapse of
\S\ref{sec:survive} that the model rediscovered: the conclusion is replaced by another statement
that holds at every aligned input (a reflexive equation, one inclusion of an equality of
subgroups, $f(1)=1$ for a ring hom $f$, the converse of $a\cdot v=0\to a=0\vee v=0$, which is
also true). The triviality flag catches 3 of the 8. This is a stress test with 150 generated
attacks, not an estimate of how often such candidates occur, and whether each candidate is
unfaithful rests on our reading, which agrees with the model's stated intent in all 8 cases.

\begin{table}[t]
\centering\small
\begin{tabular}{lrrrrrrrrrr}
\toprule
& \multicolumn{3}{c}{closed $\Ecl$} & & \multicolumn{5}{c}{ours (Opus), all verdicts} & ours + triv.\\
\cmidrule(lr){2-4}\cmidrule(lr){6-10}
candidate (80 each) & tactics & mini & Opus & BEq+ & acc. & resid. & rej. & off-dom. & unk. & rewarded\\
\midrule
$R$ (control) & 80 & 80 & 80 & 80 & 62 & 15 & 0 & 0 & 3 & 62\\
$(\Gamma,H,\top)$ & 3 & 3 & \textbf{54} & 2 & 11 & 8 & 4 & 3 & 54 & 0\\
$(\Gamma,H\wedge C,C)$ & 3 & 3 & \textbf{56} & \textbf{54} & \textbf{45} & 0 & 8 & 0 & 27 & 0\\
$(\Gamma,H\wedge\bot,C)$ & 3 & 3 & \textbf{53} & 0 & 0 & 0 & 73 & 0 & 7 & 0\\
\midrule
$(\Gamma,H,\forall k,\ 0+k=k)$ & 3 & 4 & \textbf{52} & 0 & 0 & 21 & 3 & 3 & 53 & 0\\
$(\Gamma,H\wedge(\forall k,\ k=k\to C),C)$ & 3 & 3 & \textbf{51} & 9 & \textbf{46} & 0 & 8 & 0 & 26 & \textbf{6}\\
\bottomrule
\end{tabular}
\caption{Candidates constructed to collapse, on 80 ProofNetVerif references. Closed, BEq+ and the
last column: rewarded (accepted) candidates. Middle columns: every verdict of our criterion with
Opus (rows sum to 80); \V{residual} is not rewarded. Last column: accepts not flagged by the
triviality heuristic. Below the
rule: the constant and absorbed candidates rewritten so that one \texttt{simp} step, but not
propositional logic, closes them.}
\label{tab:hack}
\end{table}

%======================================================================
\section{Selecting autoformalizations}\label{sec:select}

A faithfulness criterion is also used to choose among sampled formalizations and, as a reward,
to train the sampler. Section~\ref{sec:survive-exp} measures what each check rewards on candidates
constructed to collapse; this section measures what it picks among sampled candidates.

\paragraph{Setup.} For the 317 ProofNet-Verified problems of Table~\ref{tab:g4}, DeepSeek-V4.1-Flash,
GPT-4.1-mini and Claude Opus~5.5 each draw 8 formalizations at temperature 0.8 from the informal
statement (7,608 draws). 5,444 draws convert to a proposition and elaborate together with the
corrected reference; they contain 2,879 distinct candidates. Every candidate is compared with
the reference by closed equivalence and by our criterion (same prover and budget as in
\S\ref{sec:exp}) and judged by the four judge configurations of Table~\ref{tab:judge}.

\paragraph{Ground truth from certificates.} We classify each candidate only by kernel facts
about it and the faithful reference, independently of any selector. It is \emph{certified
different} if its closed equivalence is refuted, or if we reject it with both regrouping
certificates and a total alignment (277 draws); \emph{certified equivalent} if it has an
accept with both regrouping certificates (3,126 draws); and \emph{open} otherwise (2,041 draws).
The regrouping certificates were completed as in \S\ref{sec:exp}: 2,490 of the 2,492 decided
candidates carry both. Rejects with a partial alignment are excluded because the positional prefix
can pair binders of different meaning; counting every reject nearly doubles the certified-different
class (524 draws) and raises the best-of-24 rates below by a factor of 2.5 to 5. Neither certified class is independent of our criterion:
only 101 draws have a refuted closed equivalence, and by Appendix~\ref{sec:axioms} a strict
reject is a certified difference of hypotheses or conclusions that a reader may waive. The rows
of our own selectors are therefore consistency checks, not measurements; the comparison among
the other selectors is fair, since none of them sees these certificates.

\begin{table}[t]
\centering\small
\begin{tabular}{lrrrrrrrr}
\toprule
& \multicolumn{2}{c}{per draw (of rewarded)} & \multicolumn{3}{c}{best of 1} & \multicolumn{3}{c}{best of 24}\\
\cmidrule(lr){2-3}\cmidrule(lr){4-6}\cmidrule(lr){7-9}
selector & diff. & robust & diff. & equiv. & none & diff. & equiv. & none\\
\midrule
closed $\Ecl$ & 129 & 0 / 5047 & 1.7 (0.0) & 41.2 & 33.2 & \textbf{2.8} (0.0) & 57.9 & 0.3\\
judge Opus, NL & 159 & 21 / 5201 & 2.1 (0.3) & 41.2 & 31.3 & \textbf{3.3} (0.5) & 57.2 & 0.0\\
judge Opus, NL+ref & 115 & 0 / 5135 & 1.5 (0.0) & 41.2 & 32.2 & \textbf{2.5} (0.0) & 57.7 & 0.0\\
judge V4.1-Flash, NL+ref & 56 & 0 / 4877 & 0.8 (0.0) & 41.1 & 35.5 & \textbf{1.2} (0.0) & 60.0 & 0.9\\
judge GPT-4.1-mini, NL+ref & 115 & 34 / 4996 & 1.5 (0.5) & 40.3 & 34.1 & \textbf{2.6} (0.9) & 58.4 & 0.0\\
\midrule
judge Opus, NL+ref, vetoed by ours & 0 & 0 / 4805 & 0.0 & 41.2 & 36.6 & 0.0 & 61.4 & 0.9\\
ours, \V{accept} & 0 & 0 / 3126 & 0.0 & 41.3 & 58.7 & 0.0 & 75.4 & 24.6\\
ours, \V{accept} or \V{residual} & 0 & 0 / 4383 & 0.0 & 41.3 & 42.2 & 0.0 & 66.8 & 4.7\\
\bottomrule
\end{tabular}
\caption{Rows below the line use our certificates, which also define the classes, and are
consistency checks, not independent measurements. Best-of-$n$ selection on 317 ProofNet-Verified
problems. Per draw: rewarded draws certified different from the faithful reference; ``robust'':
only differences that survive every regrouping (refuted closed equivalence, or a reject with
\texttt{npimp}), which certify the candidate false or not closed-equivalent. Best of $n$:
first rewarded draw among $n$ random draws (50 orders per problem); percentage of problems whose
pick is certified different (robust in parentheses), certified equivalent, or for which nothing
is rewarded; the rest are open picks.}
\label{tab:select}
\end{table}

\paragraph{Selector preferences (Table~\ref{tab:select}).} The table measures which certified
classes each selector picks, not how often a pick is unfaithful; the latter needs the audit below.
Every selector that does not use certificates
rewards certified-different candidates, 1--3\% of the draws it rewards, and under best-of-$n$
selection the rate at which such a candidate is \emph{picked} grows by a factor of 1.5--1.7 from
$n=1$ to $n=24$ (closed 1.7\% to 2.8\%, Opus judge with reference 1.5\% to 2.5\%): more samples
give more chances that a formally different statement is rewarded first. The blind audit of
\S\ref{sec:blind} finds 13 of the 23 rewarded strict rejects it can judge unfaithful, so about
half of this effect is unfaithfulness and the rest differences a reader waives. The effect is small
in absolute terms and rests on our rejects: closed equivalence is never rewarded on a refuted
pair, and the judges reward 0--32 of the 101 refuted draws. For closed equivalence the 129
rewarded draws are exactly the pairs of Proposition~\ref{thm:mono}(ii): the hypotheses or the
conclusions certifiably differ at some point, but closed equivalence is proved.
Proposition~\ref{prop:split} says what this implies. Two true statements are pointwise
equivalent, so every reject between them is vacuity-only and depends on how the statement was
split into hypotheses and conclusion. A reject that survives every regrouping, one with
\texttt{npimp}, certifies that one side is false; 61 of the 69 such rejects on these candidates
also prove \texttt{sr} and hence that the candidate itself is false. Closed equivalence with a
true reference can therefore never reward a robust difference, and indeed it rewards none
(column ``robust''). The judges without the reference do: Opus reasoning from the informal
statement alone rewards 21 such draws and GPT-4.1-mini with the reference 34, so at $n=24$ they
pick a candidate certified false for 0.5\% and 0.9\% of the problems. Under the robust reading the
growth with $n$ is confined to these two judges; the rest of the effect above consists of
differences that do not survive every regrouping, about half of which a reader waives
(\S\ref{sec:blind}). Vetoing the judge's accept with our reject removes
all of them (0 of 4,805 rewarded draws) and keeps
the judge's coverage: at $n=24$ the vetoed judge picks a candidate for 99.1\% of the problems, and
the share of certified-equivalent picks rises from 57.7\% to 61.4\%. Our criterion alone covers
fewer problems (75.4\% with accepts, 95.3\% when \V{residual} is rewarded too), but every
accept it picks is certified equivalent, and every \V{residual} pick is open.

\paragraph{What this does and does not show.} The certified classes cover 63\% of the usable
draws. Two questions need a reader: how many of the strict rejects that closed equivalence or a
judge rewards are unfaithful (\S\ref{sec:framework} says some will be waivable), and how
many open picks are (38--40\% of problems at $n=24$ for the closed and judge selectors). The
blind audit of \S\ref{sec:blind} answers both, on every problem where a rewarded strict reject
occurs and on 30 open picks per selector. Of the 34 rewarded strict rejects, the annotators agree
on 23 and judge 13 unfaithful (57\%, 95\% interval 37--74\%); the other 10 are faithful candidates
whose certified difference from the corrected reference a reader waives. Of the open picks, 4 of
the 28 agreed picks of closed equivalence are unfaithful (14\%, 95\% interval 6--32\%), 2 of the
22 of the Opus judge with the reference, and 1 of the 25 of our accept (4\%, 1--20\%). Somewhat
more than half of the certified-different picks above are therefore unfaithful: at $n=24$, closed
equivalence picks a certified-different candidate for 2.8\% of the problems, and by this rate an
unfaithful one for about 1.6\%. The growth with $n$ is thus in part a Goodhart effect on
faithfulness, small in absolute terms, and the rest consists of differences a reader waives.

%======================================================================
\section{A blind two-annotator audit}\label{sec:blind}

\IfFileExists{h2_results.tex}{% generated by scripts/real/h2_tex.py from data/real/h2_scores.json; do not edit
\newcommand{\HtwoKappa}{0.66}
\newcommand{\HtwoKappaN}{226}
\newcommand{\HtwoAgree}{188}
\newcommand{\HtwoKappaFirstA}{0.46}
\newcommand{\HtwoKappaFirstB}{0.25}
\newcommand{\HtwoFullFaithfulRejectRate}{38}
\newcommand{\HtwoFullFaithfulRejectN}{43}
\newcommand{\HtwoFullUnfaithfulLocalAcceptRate}{45}
\newcommand{\HtwoFullUnfaithfulLocalAcceptN}{14}
\newcommand{\HtwoFullFaithfulLocalAcceptRate}{100}
\newcommand{\HtwoFullFaithfulLocalAcceptN}{10}
\newcommand{\HtwoFullUnfaithfulRejectRate}{0}
\newcommand{\HtwoFullUnfaithfulRejectN}{10}
\newcommand{\HtwoStrictRejectsRewardedAllRate}{57}
\newcommand{\HtwoStrictRejectsRewardedAllN}{34}
\newcommand{\HtwoStrictRejectsRewardedClosedRate}{57}
\newcommand{\HtwoStrictRejectsRewardedClosedN}{31}
\newcommand{\HtwoStrictRejectsRewardedJudgeRate}{47}
\newcommand{\HtwoStrictRejectsRewardedJudgeN}{25}
\newcommand{\HtwoOpenPicksClosedRate}{14}
\newcommand{\HtwoOpenPicksClosedN}{30}
\newcommand{\HtwoOpenPicksJudgeRate}{9}
\newcommand{\HtwoOpenPicksJudgeN}{30}
\newcommand{\HtwoOpenPicksOursRate}{4}
\newcommand{\HtwoOpenPicksOursN}{30}
\begin{table}[t]
\centering\small\setlength{\tabcolsep}{4pt}
\begin{tabular}{lrrrrrrl}
\toprule
set & items & statements & yes & no & unsure & disagree & rate [95\% CI]\\
\midrule
rejects of labelled-equivalent pairs & 43 & 30 & 10 & 16 & 0 & 4 & 38\% faithful [22--57]\\
accepts of labelled-inequivalent pairs & 14 & 11 & 5 & 6 & 0 & 0 & 45\% faithful [21--72]\\
\quad control: accepts of equivalent pairs & 10 & 10 & 10 & 0 & 0 & 0 & 100\% faithful [72--100]\\
\quad control: rejects of inequivalent pairs & 10 & 10 & 0 & 9 & 0 & 1 & 0\% faithful [0--30]\\
strict rejects rewarded by a selector & 34 & 34 & 10 & 13 & 0 & 11 & 57\% unfaithful [37--74]\\
\quad rewarded by closed equivalence & 31 & 31 & 9 & 12 & 0 & 10 & 57\% unfaithful [36--76]\\
\quad rewarded by the judge & 25 & 25 & 9 & 8 & 0 & 8 & 47\% unfaithful [26--69]\\
open picks of closed equivalence & 30 & 30 & 24 & 4 & 0 & 2 & 14\% unfaithful [6--32]\\
open picks of the judge & 30 & 30 & 20 & 2 & 0 & 8 & 9\% unfaithful [2--28]\\
open picks of our accept & 30 & 30 & 24 & 1 & 0 & 5 & 4\% unfaithful [1--20]\\
\bottomrule
\end{tabular}
\caption{Blind two-annotator audit. ``Statements'': distinct propositions (items that state the same proposition share one answer); the answer columns and rates count statements. A verdict counts only when both annotators agree; inter-annotator $\kappa=0.66$ on 226 statements.}
\label{tab:h2}
\end{table}

\newcommand{\HtwoLocKindNA}{32}
\newcommand{\HtwoLocKindAgreeA}{29}
\newcommand{\HtwoLocKindNB}{43}
\newcommand{\HtwoLocKindAgreeB}{35}
\newcommand{\HtwoMapAlignTotalN}{17}
\newcommand{\HtwoMapAlignTotalRows}{14}
\newcommand{\HtwoMapAlignTotalYes}{1}
\newcommand{\HtwoMapAlignTotalNo}{12}
\newcommand{\HtwoMapAlignEarlyN}{26}
\newcommand{\HtwoMapAlignEarlyRows}{16}
\newcommand{\HtwoMapAlignEarlyYes}{16}
\newcommand{\HtwoMapAlignEarlyNo}{10}
\newcommand{\HtwoMapLocHN}{30}
\newcommand{\HtwoMapLocHRows}{19}
\newcommand{\HtwoMapLocHYes}{12}
\newcommand{\HtwoMapLocHNo}{16}
\newcommand{\HtwoMapLocCN}{7}
\newcommand{\HtwoMapLocCRows}{7}
\newcommand{\HtwoMapLocCYes}{3}
\newcommand{\HtwoMapLocCNo}{2}
\newcommand{\HtwoMapLocHcN}{2}
\newcommand{\HtwoMapLocHcRows}{2}
\newcommand{\HtwoMapLocHcYes}{2}
\newcommand{\HtwoMapLocHcNo}{0}
\newcommand{\HtwoMapLocClosedN}{4}
\newcommand{\HtwoMapLocClosedRows}{2}
\newcommand{\HtwoMapLocClosedYes}{0}
\newcommand{\HtwoMapLocClosedNo}{4}
\newcommand{\HtwoMapDirRefFalseN}{13}
\newcommand{\HtwoMapDirRefFalseRows}{4}
\newcommand{\HtwoMapDirRefFalseYes}{8}
\newcommand{\HtwoMapDirRefFalseNo}{5}
\newcommand{\HtwoMapDirCandFalseN}{7}
\newcommand{\HtwoMapDirCandFalseRows}{7}
\newcommand{\HtwoMapDirCandFalseYes}{0}
\newcommand{\HtwoMapDirCandFalseNo}{6}
\newcommand{\HtwoMapDirOneFalseN}{2}
\newcommand{\HtwoMapDirOneFalseRows}{2}
\newcommand{\HtwoMapDirOneFalseYes}{0}
\newcommand{\HtwoMapDirOneFalseNo}{2}
\newcommand{\HtwoMapDirVacuityN}{8}
\newcommand{\HtwoMapDirVacuityRows}{7}
\newcommand{\HtwoMapDirVacuityYes}{3}
\newcommand{\HtwoMapDirVacuityNo}{3}
\newcommand{\HtwoMapDirDirectionOnlyN}{13}
\newcommand{\HtwoMapDirDirectionOnlyRows}{12}
\newcommand{\HtwoMapDirDirectionOnlyYes}{6}
\newcommand{\HtwoMapDirDirectionOnlyNo}{6}
\newcommand{\HtwoMapFaConstantN}{7}
\newcommand{\HtwoMapFaConstantRows}{7}
\newcommand{\HtwoMapFaConstantYes}{3}
\newcommand{\HtwoMapFaConstantNo}{4}
\newcommand{\HtwoMapFaTrivialN}{4}
\newcommand{\HtwoMapFaTrivialRows}{1}
\newcommand{\HtwoMapFaTrivialYes}{0}
\newcommand{\HtwoMapFaTrivialNo}{4}
\newcommand{\HtwoMapFaNeitherN}{3}
\newcommand{\HtwoMapFaNeitherRows}{3}
\newcommand{\HtwoMapFaNeitherYes}{2}
\newcommand{\HtwoMapFaNeitherNo}{1}
\begin{table}[t]
\centering\small
\begin{tabular}{lrrrrrr}
\toprule
mechanical facts & items & statements & yes & no & unsure & disagree\\
\midrule
\multicolumn{7}{l}{\emph{43 component-goal rejects of labelled-equivalent pairs}}\\
\quad alignment total & 17 & 14 & 1 & 12 & 0 & 4\\
\quad alignment not total & 26 & 16 & 16 & 10 & 0 & 0\\
\quad located in the hypotheses & 30 & 19 & 12 & 16 & 0 & 2\\
\quad located in the conclusion & 7 & 7 & 3 & 2 & 0 & 2\\
\quad located in both & 2 & 2 & 2 & 0 & 0 & 0\\
\quad closed statements only & 4 & 2 & 0 & 4 & 0 & 0\\
\quad reference certified false & 13 & 4 & 8 & 5 & 0 & 0\\
\quad candidate certified false & 7 & 7 & 0 & 6 & 0 & 1\\
\quad one side certified false & 2 & 2 & 0 & 2 & 0 & 0\\
\quad vacuity-only & 8 & 7 & 3 & 3 & 0 & 2\\
\quad direction proved, no side certified false & 13 & 12 & 6 & 6 & 0 & 1\\
\multicolumn{7}{l}{\emph{14 accepts of labelled-inequivalent pairs}}\\
\quad conclusion certified constant & 7 & 7 & 3 & 4 & 0 & 0\\
\quad candidate flagged trivial & 4 & 1 & 0 & 4 & 0 & 0\\
\quad neither & 3 & 3 & 2 & 1 & 0 & 0\\
\bottomrule
\end{tabular}
\caption{The blind audit against the kernel facts behind each reject and false accept of Table~\ref{tab:dist}, by a mapping fixed before the answers returned. The three blocks of the rejects classify the same 43 items; copies of one proposition share its answer. Yes: both annotators judge the candidate faithful. For a reject this means that the located difference does not make the candidate unfaithful: the reference is wrong, the difference is waivable, or the alignment paired the wrong binders; for an accept, that the label is wrong. Where an annotator calls a rejected candidate unfaithful and names a missing, extra or vacuous hypothesis or a different conclusion, the located component agrees in annotator A: 29 of 32; annotator B: 35 of 43.}
\label{tab:h2map}
\end{table}
}{}

Which certified differences matter, and whether an accept is faithful, needs a reader. Two
annotators judged 249 items, shown shuffled as 226 distinct statements (items that state the
same proposition share one answer, and rates count statements): the 43 rejects of
labelled-equivalent pairs that rest on component goals and the 14 false accepts of
\S\ref{sec:measure}, controls from the opposite cells, the selection sample of
\S\ref{sec:select}, and the 47 items of the earlier single-annotator classification. Each item
shows the informal statement, the reference and the candidate, and asks one question: is the
candidate a faithful formalization of the informal statement? The reference is offered as a help
that may be wrong; the sheet shows no verdict, set or selector, and the annotators were asked to
use no Lean, language model or other tool. The mapping of each item to the kernel facts behind it
(Table~\ref{tab:h2map}) was fixed and committed before the answers returned. The annotators
agree on 188 of the 226 statements ($\kappa=0.66$), and a verdict counts only where both agree
(Table~\ref{tab:h2}). The items are chosen by our verdicts, so the audit estimates how often a
verdict of a given kind is right, not how many unfaithful pairs the criterion leaves
\V{residual} or \V{unknown}. The earlier classification asked whether reference and candidate agree,
not whether the candidate is faithful, and where the reference is wrong the two questions have
opposite answers; it agrees with the new answers at $\kappa=0.46$ and $0.25$, and at $0.54$ and
$0.34$ without the pairs whose reference it marked wrong.

One annotator coded unqualified names (\texttt{log}, \texttt{card}, \texttt{ball},
\texttt{IsTopologicalBasis}) as type errors. The sheet shows each statement without the
benchmark's \texttt{open} declarations, under which these names elaborate as intended, so this is
a defect of the sheet, not of the candidates. It concerns 10 statements, on all of which the
other annotator answered ``faithful''; it turns 10 agreements into disagreements and changes no
agreed verdict.

\paragraph{Rejects of labelled-equivalent pairs.} The annotators agree on 26 of the 30 statements
behind the 43 component-goal rejects. 16 candidates are unfaithful, so the label is wrong and the
reject right; in three of them (Pugh 2.92) the reference is certified false as well, and the
candidate adds a compactness assumption that the informal statement lacks, so both sides are
unfaithful. Of the 10 faithful ones, one is Dummit--Foote 2.1.13 (8 pairs), whose reference Lean
proves false; six are faithful candidates that the certificate separates from a reference that
says something else (three Herstein 2.5.52 candidates, whose reference weakens ``more than three
quarters'' to ``at least'', and three Dummit--Foote 1.6.23 candidates, whose reference drops
finiteness); and three are canonicalization artefacts, the two of Appendix~\ref{app:rejects}
(1135 and 358) and Putnam 2018 B4, where the reference assumes $\exists n,\ x_n=0$ and the
candidate takes $n$ as a binder with $x_n=0$, which the rules do not normalize to the same
telescope. All three artefacts are vacuity-only, as Proposition~\ref{prop:split} predicts. One of
the 4 split statements is a Pugh 3.63a candidate stating
\texttt{Summable fun k => 1 / (k * (log k)\^{}p)}, which one annotator accepts and the other
rejects for an unqualified \texttt{log}; Lean proves it false, since $k$ is elaborated as a real
number and the sum ranges over $\mathbb R$, the reading the sheet warns about. By the mapping, 6
of the 7 candidates certified false are judged unfaithful, and the seventh is this one. A location
is not always the reader's: in the three Dummit--Foote 1.6.23 candidates, \texttt{[Fintype G]} is
an instance binder without a counterpart, so the alignment ends there and the difference, a
finiteness hypothesis that the reference lacks, is reported in the conclusion. Two rules written
after the audit remove these artefacts and two of the mislocations, but create three new artefacts
of the same kind (Appendix~\ref{app:k19}). Where an annotator
calls a rejected candidate unfaithful and names a missing, extra or vacuous hypothesis or a
different conclusion, the located component matches it in 29 of 32 cases for the first annotator
and 35 of 43 for the second.

\paragraph{False accepts.} Of the 11 statements behind the 14 false accepts, the annotators agree
on all 11. Six are unfaithful: Artin 11.2.13 (the four absorptions), Herstein 2.2.3, and four of
the seven with a constant conclusion (Herstein 2.1.26, Ireland--Rosen 3.4, Pugh 3.63a, Artin
2.3.2). Five are faithful, so their label is wrong and the accept right: three Pugh 2.26
candidates, whose conclusions are constant on the aligned inputs, Artin 10.6.7, and Putnam 2018
A5, where the annotators read \texttt{ContDiff}~$\mathbb R$~$\top$ as smoothness; in the Mathlib
we use it denotes analyticity, so both sides misstate the exercise (\S\ref{sec:audit}). A reader
thus confirms 9 of the 14 false accepts, and a constant conclusion is not by itself
unfaithfulness: three of the seven constant accepts are faithful.

\paragraph{Controls.} Of 10 sampled accepts of equivalent pairs with an informative conclusion,
all 10 are judged faithful; of 10 sampled rejects of inequivalent pairs, 9 are judged unfaithful
and 1 is split. In the selection sample, 7 of the 10 candidates certified equivalent to the
corrected reference are judged faithful and 3 are split, two of them by unqualified names.

%======================================================================
\section{Auditing references with single-statement probes}\label{sec:audit}

Every experiment above trusts the reference. The same goals also audit it, by comparing a
statement with its own weakenings.

\paragraph{Probes.} For $S=(\Gamma,H_1\wedge\dots\wedge H_m,C)$ let $S_{\setminus k}$ drop $H_k$
and write $\bar x$ for the variables and $H_{\setminus k}$ for the remaining hypotheses. We state
three goals per hypothesis:
\[
\texttt{implied}_k:\ \forall\bar x.\,H_{\setminus k}\to H_k,\qquad
\texttt{nimplied}_k:\ \neg\,\texttt{implied}_k,\qquad
\texttt{idle}_k:\ \forall\bar x.\,H_{\setminus k}\to\neg H_k\to C.
\]
\begin{proposition}\label{prop:probe}
$\texttt{implied}_k$ holds iff $S\Ecomp S_{\setminus k}$. If $\Cl S$ holds, then
$\texttt{idle}_k\wedge\texttt{nimplied}_k$ holds iff $\Cl S_{\setminus k}$ holds and
$\neg(S\Econd S_{\setminus k})$.
\end{proposition}
\begin{proof}
The conclusions coincide, so $S\Ecomp S_{\setminus k}$ and $S\Econd S_{\setminus k}$ both reduce to
$\forall\bar x.\,(H_{\setminus k}\wedge H_k)\leftrightarrow H_{\setminus k}$, which is
$\texttt{implied}_k$. Given $\Cl S$, i.e.\ $\forall\bar x.\,H_{\setminus k}\to H_k\to C$,
$\texttt{idle}_k$ is equivalent to $\forall\bar x.\,H_{\setminus k}\to C$ by cases on $H_k$.
\end{proof}
So an \texttt{implied} certificate says that $H_k$ is redundant, and an \texttt{idle} one says
that the formal statement remains a theorem without $H_k$ although $H_k$ excludes real points:
the hypothesis is decorative. Closed equivalence identifies $S$ with $S_{\setminus k}$ in both
cases (Proposition~\ref{thm:repr}); the component criterion separates the second case, and by
\texttt{true\_true\_reject\_iff} the separation is exactly a hypothesis that does no work.
Neither is an error by itself: competition statements restate positivity conditions, and
junk values make some hypotheses idle. A flag is therefore a kernel-checked pointer that a
human resolves against the informal statement, and recall is limited to errors that surface as
a redundant or decorative hypothesis.

\paragraph{What single-hypothesis probes miss.} Two whole-statement probes complement them:
$\texttt{vacuous}:\ \forall\bar x.\,H\to\bot$ (the hypotheses are jointly unsatisfiable) and
$\texttt{trivial}:\ \forall\bar x.\,C$ (the conclusion holds with every free hypothesis
removed). A vacuous statement is a theorem for every conclusion and a trivial one for every
hypothesis (\texttt{T8.vacuous\_iff}, \texttt{trivial\_closed}). Their relation to the
per-hypothesis probes is asymmetric. If $C$ is trivial then every hypothesis is idle, so
single probes already flag a trivial conclusion as soon as one hypothesis is not implied by
the others (\texttt{T8.trivial\_idle}). Joint vacuity, however, is invisible to them: for
$H_1=(0<x)$, $H_2=(x<0)$ and $C=(x=1)$ over $\mathbb Z$ the hypotheses are individually
satisfiable, neither is implied by the other, neither is idle (the conclusion fails on each
alone), and the statement is vacuous (\texttt{T8.joint\_vacuous\_missed}). The
\texttt{vacuous} probe closes exactly this gap; \S\ref{sec:global} reports what it adds.

\begin{table}[t]
\centering\small
\begin{tabular}{lrrrrl}
\toprule
benchmark & statements & probed & flagged & errors & external check\\
\midrule
ProofNet\# (valid+test) & 358 & 240 & 39 & 13 (+11) & of the 24, 22 confirmed by~\cite{pnv}, 2 missed\\
miniF2F (DeepSeek port) & 488 & 344 & 93 & 8 & each kept by v2~\cite{minif2fv2} or DM~\cite{minif2fdm}\\
PutnamBench (2026-10-06) & 672 & 481 & 87 & 2 & unchanged upstream\\
\bottomrule
\end{tabular}
\caption{Per-hypothesis probes on 1,518 benchmark references. ``probed'': statements with at
least one Prop hypothesis whose goals elaborate; ``errors'': flagged statements that are
unfaithful to the informal statement after manual review (Appendix~\ref{app:refs},
Table~\ref{tab:bench}); +11: errors found by pairwise rejects and not flagged by a probe.}
\label{tab:audit}
\end{table}

\paragraph{ProofNet\#.} Running the probes on all 358 ProofNet\# references
flags 39 statements, each with a kernel certificate. Checking each flagged statement against
its informal statement finds 13 genuine errors, none fixed in ProofNet\#: conclusions such as
\texttt{Nonempty (CommGroup G)} that hold for every group (Herstein 2.5.43, 4.2.5), hypotheses
such as \texttt{IsEmpty (CommGroup G)} that are contradictory and make the statement vacuous
(Herstein 2.5.37, 2.8.12), closure in place of limit points making a statement false (Rudin
2.24), a circle formalized as the two-point sphere of $\mathbb R$, which makes the
statement false (Munkres 24.2), a series written as a sequence (Pugh 3.63a), and six others.
With one more found by the pairwise criterion on ProofNet-Verified (Munkres 21.6b,
Table~\ref{tab:g4}), these are 14 wrong references; ten more found by pairwise rejects on
ProofNetVerif (below) bring the total to 24 of 358 (6.7\%). Appendix~\ref{app:refs} lists the
fourteen with the certified goal and the defect; the remaining 26 probe flags are benign
redundant hypotheses or junk-value degeneracies, so the probe precision is 13 of 39 (33\%).

ProofNet-Verified~\cite{pnv}, an independent manual audit of all 367 ProofNet\# statements,
lets us measure both directions. It confirms 12 of our 14 errors. The other two,
Munkres~24.3a and 21.6b, are marked faithful there but are false: without a type ascription
\texttt{Icc 0 1} is elaborated in $\mathbb N$, a two-point discrete space; swapping its points
is a continuous map without a fixed point, and $x^n$ is eventually constant and converges
uniformly (\texttt{BenchErrors.lean} proves both negations). The
audit also shows the limits of the probe. Of the 96 statements it marks unfaithful, 66 have a
Prop hypothesis and can be probed, and only 17 are flagged; for 5 of these the flagged
hypothesis is benign and the defect lies elsewhere (e.g.\ Rudin~5.7 takes the limit along
the full neighbourhood filter $\mathcal N(x)$ instead of the punctured $\mathcal N_{\ne}(x)$, which
makes it false). The probe is a triage tool with a kernel-checked hit list, not an
auditor: it reduces inspection from 358 statements to 39, of which at least 18 are erroneous,
and its recall is limited to errors that make some hypothesis idle or implied.

\paragraph{Pairwise rejects as an audit.} The probes compare a reference with its own weakenings;
candidates offer other comparisons for free. A reject certified by \texttt{npimp} and the
direction \texttt{rs} proves that the reference is false (Proposition~\ref{prop:split}(iv)). On the
1,140 ProofNetVerif pairs with Opus this happens for 32 pairs and 11 references, all of them in the
ProofNet\# test split and all unchanged there. Two are already in Appendix~\ref{app:refs} or
\S\ref{sec:survive-exp} (Munkres 24.3a, Rudin 1.8). The other nine are new to our list:
\begin{itemize}
\item a hypothesis stated for one $x$ instead of all (Dummit--Foote 2.1.13);
\item a missing assumption: Hausdorff for nested compact sets (Pugh 2.92), nontriviality for
groups without proper subgroups (Herstein 2.3.16), surjectivity for images of normal subgroups
(Herstein 2.7.7), $T_1$ next to \texttt{RegularSpace} (Munkres 31.1);
\item \texttt{Monotone} for ``monotonic'', which excludes $x\mapsto -x$ (Rudin 4.15);
\item a wrong order for a primitive root, $p$ for $p-1$ (Ireland--Rosen 4.6);
\item a summand $X^n$ for $X^i$, so that $P_1=2X$ and $P_2=6X^2$ are not coprime
(Putnam 2014 A5);
\item bounded intervals for intervals, so $\mathbb R$ itself has no decomposition (Rudin 2.29).
\end{itemize}
ProofNet-Verified marks all ten as unfaithful and unprovable, and 24.3a as faithful, which it is
not (above), so on this list the certificates and the independent manual audit agree. For a
reader who does not want to follow the regrouping argument, \texttt{BenchErrors.lean} also
refutes each of the ten statements directly (Table~\ref{tab:refs}). None of
the ten is flagged by a probe: their errors are false conclusions or missing hypotheses, not
idle or implied ones. The certificate needs only a candidate whose comparison with the reference
the prover decides, so every benchmark that collects candidates also audits its references.

\paragraph{Competition benchmarks.} We ran the same probes on the 488 statements of the
DeepSeek-Prover Lean~4 port of miniF2F~\cite{minif2f} and the 672 Lean statements of
PutnamBench~\cite{putnambench} (answer constants made \texttt{opaque}, so a flag holds for every
answer). 93 and 87 statements are flagged. Competition statements carry many redundant
positivity conditions, so precision is lower (Table~\ref{tab:bench}): 8 of the 93 miniF2F flags
and 2 of the 87 PutnamBench flags are genuine errors. For each hard error
\texttt{lean/ExperimentsMathlib/BenchErrors.lean} keeps the benchmark hypotheses verbatim and
derives \texttt{False} or the negation of the statement. Two cases show what the probe adds over
manual review. In \texttt{aime\_1996\_p5} the cubic $x^3+3x^2+4x-11$ is strictly increasing, so
the hypothesis that $a,b,c$ are three distinct real roots is unsatisfiable; the probe flags it
because the hypotheses become mutually implied. In \texttt{putnam\_1991\_a6}, $\to$ binds
weaker than $\wedge$, so the set counted by $A(n)$ contains every tuple that is not a partition
of $n$, is infinite, and contradicts $A(n)\in\mathbb N$. The statement was correct until a
readability refactoring inlined an auxiliary predicate into the set-builder without
parentheses; no proof attempt can notice this, since every prover then succeeds. We compared each miniF2F error with the
two corrected releases. Neither fixes all eight, and each error survives in at least one of
them. DeepMind's fork~\cite{minif2fdm} fixes the three certified ones and keeps the other five.
miniF2F-v2~\cite{minif2fv2}, which reports that all its statements are correct and provable,
fixes the five and keeps the three certified ones:
\texttt{aime\_1996\_p5} unchanged, \texttt{amc12a\_2020\_p13} with natural-number exponents
\texttt{1/c = 0} (false in both v2s and v2c), and a rewrite of \texttt{amc12a\_2003\_p25} whose
domain $\{x\mid 0\le f\,x\}$ compares a curried function pointwise and is therefore all of
$\mathbb R$ (false in v2s, answer set empty in v2c). \texttt{putnam\_1991\_a6} had not been reported;
\ifanonymous it has since been reported to the maintainers with a one-line fix.\else we reported it
upstream together with a one-line fix.\fi{} One more
PutnamBench flag points to an error in the informal text: the docstring of
\texttt{putnam\_2021\_b4} bounds the product by $F_{m-1}$, while the official problem and the
formal statement use $F_m-1$. A reference can also change meaning under a library upgrade,
which no probe sees: ProofNet's Putnam 2018 A5 renders ``infinitely differentiable'' as
\texttt{ContDiff}~$\mathbb R$~$\top$~\texttt{f}, which the Mathlib we use reads as analytic
(smooth is \texttt{ContDiff}~$\mathbb R$~$\infty$~\texttt{f}). The candidate in ProofNet-Verified writes the same, so the pair is
equivalent, but neither side states the exercise.

\paragraph{What the evidence is.} The 34 errors rest on evidence of two strengths. For 20, Lean
proves from the benchmark's own statement that it is false or vacuous: 13 ProofNet\# negations
(Table~\ref{tab:refs}, ``cert.'' and the lower block), the 3 ProofNet\# statements whose
hypotheses the \texttt{vacuous} probe refutes, 3 miniF2F statements and
\texttt{putnam\_1991\_a6}. An exercise or competition problem that is false or vacuous as stated
is wrong whatever its informal text says, so these need no reader. For the other 14 the kernel
certifies a degeneracy (an idle or implied hypothesis, a conclusion that holds everywhere) and
the error is the reading of that degeneracy against the informal statement; there the
certificate locates the defect and a human confirms it. Manual curation is quality control, not
a guarantee, and the certificates show where it failed: ProofNet-Verified, an independent audit
of every ProofNet\# statement, marks Munkres 21.6b and 24.3a faithful although Lean refutes
both; miniF2F-v2, which reports every statement correct and provable, keeps three statements
that Lean proves false or vacuous; and \texttt{putnam\_1991\_a6} had not been reported. Conversely,
22 of our 24 ProofNet\# errors are also in ProofNet-Verified and each miniF2F error is fixed in
one of the two corrected releases, so the probes rediscover what careful audits find and add a
machine-checked reason. Their precision alone (13 of 39, 8 of 93, 2 of 87) makes them a triage
tool, not an auditor. Each of the 34 errors, with its evidence and its status in the independent
records, is a row of \texttt{data/real/bench\_errors.csv}.

\begin{table}[t]
\centering\footnotesize
\begin{tabular}{llp{7.4cm}l}
\toprule
statement & probe & defect & v2 / DM\\
\midrule
\texttt{aime\_1996\_p5} & $H_1,H_2,H_3$ implied & cubic has one real root; hypotheses unsatisfiable (cert.) & kept / fixed\\
\texttt{amc12a\_2003\_p25} & $H_0$ implied & \texttt{Real.sqrt} total, so the domain is $\mathbb R$; $h_2$ never holds (cert.) & still wrong / fixed\\
\texttt{amc12a\_2020\_p13} & $H_2$ implied & \texttt{1/c} in $\mathbb N$ is $0$; $h_2$ reads $1=1$; statement false (cert.) & kept / fixed\\
\texttt{mathd\_algebra\_76} & $H_0$ idle & proves $f(4)=-1$ instead of $f^{5}(4)=1$; odd branch unused & fixed / kept\\
\texttt{amc12a\_2002\_p6} & $H_0$ idle & $\exists p,\ mp\le m+p$ holds with $p=0$; trivially true & fixed / kept\\
\texttt{mathd\_algebra\_144} & $H_0$ idle & count of triangles replaced by the bound $d<10$; $d>0$ unused & fixed / kept\\
\texttt{amc12\_2000\_p6} & $H_1$ idle & proves $pq-(p+q)\ne194$, unrelated to the answer $119$ & fixed / kept\\
\texttt{imo\_1984\_p2} & $H_1$--$H_3$ idle & ``find a pair'' replaced by the bound $19\le a+b$ & v2c only / kept\\
\midrule
\texttt{putnam\_1991\_a6} & $H_0,H_2,H_3$ implied & $\to$ binds weaker than $\wedge$; counted set infinite; vacuous (cert.) & reported\\
\texttt{putnam\_1966\_b2} & $H_0$ idle & adds $n>0$ to ``any ten consecutive integers''; weaker than stated & unchanged\\
\bottomrule
\end{tabular}
\caption{Reference errors among the 93 flagged miniF2F and 87 flagged PutnamBench statements.
``cert.'': \texttt{BenchErrors.lean} derives \texttt{False} or the negation from the benchmark's
own hypotheses (standard axioms only). Last column: status in miniF2F-v2~\cite{minif2fv2} and in
the DeepMind fork~\cite{minif2fdm}, or in the PutnamBench repository.}
\label{tab:bench}
\end{table}

\paragraph{Whole-statement probes.}\label{sec:global} Running \texttt{vacuous} and
\texttt{trivial} on the 335 ProofNet\# statements whose goals elaborate (fixed tactics, then
one proposal and one repair round) flags 14: 3 vacuous and 11 trivial. The three vacuous
statements are Herstein 2.5.37 and 2.8.12 (\texttt{IsEmpty (CommGroup G)}) and
Ireland--Rosen 2.21, all among the fourteen errors, now with a direct certificate
$\forall\bar x.\,H\to\bot$ instead of an implied-hypothesis flag. Twelve of the 14 flags
coincide with single-hypothesis flags, as \texttt{T8.trivial\_idle} predicts for trivial
conclusions. The two new flags are true theorems with a hypothesis that does no work: Axler 3.8
holds without \texttt{FiniteDimensional F V} (a complement of the kernel exists in any
vector space), and Dummit--Foote 2.4.16b holds for $n=0$ as well (the rotations of the
infinite dihedral group have index two). Neither is a formalization error. On ProofNet\#,
then, joint vacuity does not occur without single-hypothesis symptoms, and the
whole-statement probes add certificates rather than recall; the recall ceiling against
ProofNet-Verified remains the 30 of 96 unfaithful statements that have no Prop hypothesis
at all and the errors that leave every hypothesis essential. The same probes on the 486
miniF2F and 654 PutnamBench statements that elaborate flag 15 and 5 respectively
(miniF2F: 5 vacuous, 10 trivial; PutnamBench: 1 vacuous, 4 trivial). Two of the five
miniF2F vacuous flags and one trivial recover errors already in Table~\ref{tab:bench}
(\texttt{aime\_1996\_p5}, \texttt{amc12a\_2003\_p25}, \texttt{amc12a\_2002\_p6}); the other
three vacuous flags (\texttt{mathd\_numbertheory\_35}, with $\forall n,\,n\mid\sqrt{196}$
in place of ``$n$ ranges over the divisors''; \texttt{mathd\_algebra\_437}, whose bounds
$x<n<y$ reverse $\sqrt[3]{-101}<\sqrt[3]{-45}$; \texttt{mathd\_algebra\_275}) are additional
candidates that the single-hypothesis probes missed, and the remaining trivials are idle
positivity or identities that hold without the dropped hypothesis. On PutnamBench the one
vacuous flag is the known \texttt{putnam\_1991\_a6}; of the four trivials, one is the known
weakening \texttt{putnam\_1966\_b2} and the other three (\texttt{putnam\_1971\_b3},
\texttt{putnam\_1997\_b5}, \texttt{putnam\_2000\_b2}) are true with an idle hypothesis. Whole-statement probes therefore
add certificates on ProofNet\# and a small number of extra flags on the competition sets.

\paragraph{Beyond benchmarks.} The same comparison applies whenever two kernel-checked
statements are claimed to say the same thing. Appendix~\ref{app:kaplansky} compares the
challenge statements of a recently published Lean-verified counterexample with the main
theorem of its manuscript: closed equivalence must accept, since both are proved, while the
criterion certifies in one line that the headline challenge assumes less about the group.

%======================================================================
\section{Related work}\label{sec:related}

Table~\ref{tab:related} summarizes the comparison by capability. The combination that the
criterion adds is the second, third and fourth column together: kernel evidence for a
difference, the component it lies in, and fewer false accepts among true, non-equivalent pairs
without restricting the prover. Each of these exists separately in earlier work.

\begin{table}[t]
\centering\footnotesize\setlength{\tabcolsep}{3pt}
\begin{tabular}{p{0.25\linewidth}ccccc}
\toprule
method & \makecell{evidence for\\``equivalent''} & \makecell{evidence for\\``different''} & \makecell{locates the\\difference} & \makecell{reduces false\\accepts among true,\\non-equivalent pairs} & \makecell{undecided\\reported}\\
\midrule
closed equivalence, unrestricted & kernel & kernel & no & no & yes\\
BEq+~\cite{beq} & kernel & no & no & \makecell{by restricted\\search} & no\\
SemantEq~\cite{semanteq} & \makecell{kernel,\\audited} & no & no & \makecell{by allowed-step\\relation} & no\\
symbolic equivalence~\cite{symeq} & prover & no & components & yes & no\\
SA-Pass~\cite{shadowbench} & kernel & no & by shadow & by restriction & no\\
fingerprints~\cite{bpf}, T$^2$~\cite{t2} & kernel & no & no & \makecell{no (complete\\oracle)} & no\\
FormalRx-CC~\cite{formalrxcc} & -- & kernel & given subterms & yes & --\\
LLM judges~\cite{beyondcomp,criticlean} & no & no & in text & yes & no\\
ours & kernel & kernel & \makecell{component,\\or closed only} & \makecell{on aligned\\components} & yes\\
\bottomrule
\end{tabular}
\caption{Capabilities of reference-based checks, as described in the cited papers.
``Reduces false accepts among true, non-equivalent pairs'': avoids accepting as equivalent two
statements that are both true but make different mathematical claims. It does not mean that a
pair the method does not accept is proved different; that is the column ``evidence for
different''. ``By restriction'', ``by restricted search'', ``by allowed-step relation'': only
because the prover, the allowed proofs or the admitted relation are restricted (BEq+ requires each
proof to use the other statement; SemantEq admits only chains of pre-specified steps), so a pair
that is not accepted may still be equivalent. ``Kernel, audited'': the kernel checks the proof
term, and an audit decides whether the proof stays within the allowed relation. No evidence for
``equivalent'' in this table shows faithfulness to the informal statement. ``Undecided
reported'': not accepting is split into a separate undecided verdict whose causes are reported;
a ``no'' does not mean that a paper reports no recall or no count of pairs it does not accept.
``--'': not an aim of the method.
Single-statement checks (vacuity, counterexamples~\cite{faults,meidouglas}) are complementary and
not listed.}
\label{tab:related}
\end{table}

\paragraph{Reference-based equivalence.} Cross-provability between a candidate and a reference
is the most widely used automatic criterion: Euclidean geometry~\cite{eucgeo}, BEq with an LLM
prover~\cite{beqliu} and its deterministic successor BEq+ with the ProofNetVerif labels and
ProofNet\#~\cite{beq}. That all true statements are provably equivalent is folklore (they form
the top element of the Lindenbaum--Tarski algebra), and practitioners know that a strong prover
exploits it: Poiroux et al.\ observe that an unrestricted \texttt{exact?} makes ``unrelated true
theorems'' equivalent and restrict the tactic to the other statement~\cite{beq}, and the case for
theory-level autoformalization states the collapse and proposes restricting equivalence checks to
an admissible fragment of the background theory~\cite{tla}. The recently introduced
SemantEq~\cite{semanteq} restricts which proofs count: it accepts a pair only through chains of
pre-specified transformation steps (logical rearrangement, rewriting with the library's
equations and biconditionals, a change of variables, structural weakening), checked on the proof
term, and that a rewrite preserves meaning rests, as its authors note, on how well the library's
equations reflect mathematical usage. Because the rule is fixed before proving, a stronger prover
raises its recall without changing what an accept means. On the ProofNetVerif test split,
SemantEq-Eval reports 100\% precision under its adjudicated labels, with 65.6\% recall; under two
of four prompts a language model produces proofs in both directions for about two fifths of the
non-equivalent pairs, which is why it audits proof terms rather than relying on compilation. A
pair it does not accept is undecided: no admissible chain passed the audits, which does not show
that the statements differ. Our criterion restricts which propositions are proved instead: it
fixes the component goals and accepts any kernel proof of them, so a reject is a kernel proof
that the aligned components differ, while an accept still collapses on constant conclusions and
absorption, which a restriction on proofs such as SemantEq's is designed to exclude
(\S\ref{sec:survive}). The two give evidence on opposite sides and compose
(\S\ref{sec:measure}). SemantEq's authors also re-adjudicated many of the relevant
ProofNetVerif pairs by hand and relabel 46 of our 51 rejects as non-equivalent; this is external
but selection-dependent evidence, since most reviewed pairs were prompted by SemantEq's own
failures to accept them and the review covers only pairs originally labelled equivalent. We
add a measurement of how fast the error grows with today's provers, across models and across the
budgets of one model (Tables~\ref{tab:full}, \ref{tab:sweep} and~\ref{tab:budget}), and of what
restricting the prover costs in recall (Table~\ref{tab:g1}). Later work reads BEq+ non-equivalence as a
conservative floor~\cite{paraphrase}, and the ProofNet audit of Ground False notes that equivalence
metrics compress model gaps through ``coincidental agreement''~\cite{pnv}.

\paragraph{Faithfulness as a training signal.} Autoformalizers are trained with BEq as a
verifiable reward~\cite{stepfun}, with language-model consistency checks as the
reward~\cite{formarl,reform}, or with a trained critic~\cite{criticlean}; with type-checking alone as
the reward, outputs become vacuous~\cite{onlinerl}. These are the settings in which a check is
optimized against, and \S\ref{sec:survive-exp} measures what each kind of check rewards on
candidates constructed to collapse.

\paragraph{Component comparison.} Comparing hypotheses with hypotheses and conclusions with
conclusions has precedents. Symbolic equivalence~\cite{symeq} calls $P_1\to Q_1$ and $P_2\to Q_2$
equivalent when $P_1\equiv P_2$ and $Q_1\equiv Q_2$ and tests premises for contradiction, to
select among candidates in Isabelle without a reference; this is $\Ecomp$ with a vacuity check.
VERINA~\cite{verina} judges generated program specifications by preconditions and by
postconditions under the reference precondition, our $\Econd$, with an \V{unknown} outcome
reported as bounds; CLEVER~\cite{clever} compares specifications quantified over the
implementation, an open comparison that does not collapse. LeanScorer~\cite{mathesis} matches
premises and conclusions with an LLM, and BEq+ proves the two directions of an entailment.
Premise checks, conclusion checks and directional checks are therefore all in use, and the split
does not by itself distinguish our criterion from them. What does is how the evidence is
produced and what it is tied to. The goals are fixed by a metaprogram before any prover runs, as
SemantEq fixes its allowed steps for accepts, whereas FormalRx-CC (below) lets a model write the
witness of a difference, and 31\% of the
witnesses that compile in its archived runs still compile once the two objects are made
identical (Appendix~\ref{app:more}); each verdict is a kernel proof, where LeanScorer's matching
is an LLM judgment; the regrouping certificates tie
the compared components to the original statements (Proposition~\ref{prop:cert}), which a
decomposition computed by code otherwise leaves unchecked; and the outcomes separate a located
difference from one certified only between the closed statements, an alignment that ends early,
an undecided pair and the accepts that still collapse (\S\ref{sec:survive}), each counted on the
full set (Table~\ref{tab:dist}). The closest
conceptual antecedent is
vacuity detection in model checking~\cite{beer,kupfermanvardi,gurfinkel,namjoshi}, where a
property that holds because a subformula is irrelevant is flagged as vacuous, multi-valued
semantics separate the cases, and proof-based vacuity reads irrelevance off a certificate;
our (Vac) axiom, the third value and the per-hypothesis probes of \S\ref{sec:audit} are the
counterparts for comparing two statements. For a true telescope the probes are exactly these
tests: vacuity in $H_k$ is $\texttt{idle}_k$, vacuity in $C$ is \texttt{vacuous}, and mutual
vacuity in all hypotheses~\cite{gurfinkel} is \texttt{trivial} (\texttt{T8}). A hypothesis implied
by the others is vacuous in this sense but harmless, which is why a flag also requires
$\texttt{nimplied}_k$, and each probe verdict is a kernel proof, a proof-based vacuity check in
the sense of~\cite{namjoshi}.

\paragraph{Probe- and test-based criteria.} Provability fingerprints~\cite{bpf} compare forward
and backward consequence neighbourhoods; SA-Pass~\cite{shadowbench} checks a candidate against
a complete set of shadow statements; T$^2$~\cite{t2} replaces a theorem in its repository and
recompiles its successors. With a complete, unrestricted oracle all three are $\Ecl$-invariant
and collapse on true statements (Proposition~\ref{thm:repr}); SA-Pass avoids this by allowing only
direct projections of the candidate, and its authors note that shadow sets must be chosen to avoid
collapse. T$^2$ is moreover forward-only, so a strengthened or
vacuous candidate passes whenever it is accepted with \texttt{sorry}. FormalRx-CC~\cite{formalrxcc}
is closest in spirit: it notes that two true theorems may be inseparable and asks for local
witnesses on benchmark-supplied subterms. Our framework fixes the meaning of such witnesses in
advance, which is what its archived runs lack (\S\ref{sec:exp}). Reference-free mechanical
checks~\cite{meidouglas} evaluate a statement numerically and replace its goal by an opaque
proposition to detect proofs that ignore it; they report 14 accepted errors among 153 ProofNet
items and are complementary to a reference-based comparison.

\paragraph{Similarity and judges.} Tree edit distance on operator trees (GTED~\cite{gted},
ASSESS~\cite{assess}), learned alignment scores (FormalAlign~\cite{formalalign}) and LLM judges,
alone or with human calibration~\cite{beyondcomp,atf}, give graded or binary scores without
certificates. They are not closed-equivalence invariant and so escape Proposition~\ref{thm:repr}, at
the price of violating soundness (S) or using information beyond the formal statement; none can
tell an \V{off-domain} difference from a real one.

\paragraph{Benchmark audits.} ProofNet\# corrects 31.8\% of the ProofNet Lean~4
port~\cite{beq,proofnet}, and Ground False~\cite{pnv} re-audits ProofNet\# by hand (released as
ProofNet-Verified) and finds 97 of 367 statements unfaithful or ambiguous, confirming errors by proving negations or
trivialized proofs; miniF2F-v2~\cite{minif2fv2} and DeepMind's fork~\cite{minif2fdm} correct
miniF2F~\cite{minif2f}; PutnamBench~\cite{putnambench} fixes misformalizations through its issue
tracker. ``Faults in our formal benchmarking''~\cite{faults} audits five benchmarks with Lean
static checkers (vacuous hypotheses by proving \texttt{False}, counterexamples, junk-value guards)
plus LLM review, and reports 398 certified issues. Our per-hypothesis probes generalize their
vacuity check: \emph{implied} flags redundant or contradictory hypotheses, and \emph{idle} flags
conclusions that hold without a hypothesis, which is how weakened goals such as
\texttt{mathd\_algebra\_76} and \texttt{Nonempty (CommGroup G)} surface. That every miniF2F error we
certify survives in at least one corrected release suggests that manual correction and static
checks are complementary to certified probes.

%======================================================================
\section{Limitations and conclusion}\label{sec:conclusion}

\paragraph{Limitations.} Four limits bound the claims. First, component verdicts are relative to
the exhibited canonical telescopes (\S\ref{sec:framework}). The regrouping certificates fix the
closures, not the split into hypotheses and conclusion, so the canonicalization rules are part of
the criterion, and where they split two presentations of the same content differently the
verdict changes. On the 300-pair subset this happens twice, and both cases are rejects of
labelled-equivalent pairs: in pair 1135 \texttt{[Fintype G]} is data and \texttt{[Finite G]} a
proposition, so the same assumption lands in the context on one side and in the hypotheses on the
other; in pair 358 the statements differ by a de Morgan step plus currying, which the rules do not
normalize to the same telescope. On the full set, 22 of the 43 rejects of labelled-equivalent
pairs that rest on component goals locate the difference on an alignment that is not total and
may include more such cases; the blind audit finds one, Putnam 2018 B4, where an existential
hypothesis and a binder are split differently (\S\ref{sec:blind}). A uniqueness certificate for the canonical form would
close this gap; local rules do not. The two rules of Appendix~\ref{app:k19}, one of which keeps the
alignment going past a binder without a counterpart by reading it as a \texttt{Nonempty}
hypothesis, remove the artefacts they target and create three of the same kind. The alignment
itself is positional, by type and role; other strategies, such as matching binders by the
hypotheses that mention them, are untested. Second, the cross-model tables compare different models ordered by measured
proving power, and the controlled sweep of Table~\ref{tab:budget} varies only the number of
proposal and repair rounds of two models; neither is a law in budget or search effort. The cost
comparison of Table~\ref{tab:cost} counts Lean checks rather than Lean time, and it compares the
two checks as implemented, not their cheapest implementations; it shows that our precision is
not bought with budget, not that the decomposition decides more pairs at equal cost. Third, the criterion certifies formal relations, not faithfulness: an accept is
necessary but, as absorption shows, not sufficient, and a reject is a formal difference that a
reader may waive (a decorative hypothesis, a contrapositive reading). Fourth, which differences
matter rests on human judgment. The blind audit of \S\ref{sec:blind} has two annotators who agree
on 83\% of the statements ($\kappa=0.66$), so a verdict on the rest is open; it covers samples,
not every reject; and its sheet showed the statements without the benchmark's \texttt{open}
declarations, which one annotator read as type errors. The classification of the subset's rejects in
Appendix~\ref{app:rejects} is by one annotator. The audit's analysis was fixed in advance: each of
the 43 rejects and 14 false accepts is mapped to the kernel facts behind it (alignment, located
component, direction certificate, constancy, triviality), and the kind of difference an
annotator names is compared with the located component, by a rule chosen before the answers
returned. The
adjudication of \S\ref{sec:measure} is independent of our verdicts and agrees with most rejects
of labelled-equivalent pairs, but it uses another criterion (the same claim as the reference),
was prompted mostly by another method's failures, and does not cover the false accepts.

Alignment beyond permutations and standard embeddings (e.g.\ \texttt{List} against
\texttt{Finset}) is open, and positional alignment is broken by an inserted instance binder
(two of the four \V{residual} verdicts on unfaithful pairs in Table~\ref{tab:g4}); $\Ecomp$ is strict on junk values,
which guarded equivalence repairs only partially; the transfer theorem gives decidability
only for the arithmetic fragments, which cover 3\% of ProofNetVerif but a third of
miniF2F (Table~\ref{tab:frag}), while for structures and type classes we report only the
empirical \V{unknown} rates; and the labels of ProofNetVerif carry noise, as do its
references, which come from ProofNet\# and are unfaithful in 26.4\% of cases~\cite{pnv}.
A wrong reference turns a faithful candidate into a correct reject and an equally wrong one into
an accept, so Table~\ref{tab:g1} measures agreement with the reference; Table~\ref{tab:g4} uses
corrected references and avoids this. Every reject is a
kernel certificate and therefore holds regardless of the labels, but the classification of
rejects into label errors, tolerated hypothesis differences and artefacts rests on a single
annotator (Appendix~\ref{app:rejects}); the ProofNet-Verified experiment (Table~\ref{tab:g4})
uses labels we did not produce and reaches the same conclusion. Seven of 299 re-presentations in
\S\ref{sec:exp} remain open, all predicted by Proposition~\ref{prop:canon}. The probes of
\S\ref{sec:audit} have low recall against a full manual audit (12 of its 96 unfaithful
statements are flagged at the defect) and are meant to complement it.

\paragraph{Conclusion.} Comparing closed statements up to provable equivalence cannot separate
true theorems, and with the provers in use today the resulting false accepts reach 39\% of
unfaithful pairs. Comparing canonical components with goals fixed in advance makes every
decided verdict a kernel certificate about the exhibited telescopes whose meaning does not depend
on which prover found it, confines what a stronger prover can wrongly accept to pairs that are
component-equivalent relative to the theory and the canonical form, and reports separately the
differences no formal criterion can resolve, those outside the hypotheses, and the pairs whose
alignment ends early, where the closed collapse returns. It pays for this in coverage, which it
reports outcome by outcome, and gains a mechanical location for most of the differences it
certifies, which a reader still has to confirm. Restricting the proofs an accept may use, as
SemantEq does, and certifying differences with an unrestricted prover answer complementary
halves of the question. What remains in that set, constant conclusions and absorption, is a limit of every
extensional check: a component accept is a necessary condition and a reject a certified
difference, and deciding faithfulness still needs a reader. The same goals applied to a statement and its own
weakenings find errors that survive manual correction in every benchmark we examined.

\appendix
\section{Theory}\label{app:theory}

Fix a consistent, recursively axiomatized $\T$ interpreting $\mathsf{PA}$. Statements about
provability are metatheorems; the non-metamathematical parts are formalized in Lean~4
(\texttt{lean/ExperimentsMathlib/T1--T9}, standard axioms only). Statements are aligned, so they
share a context $\Gamma$. None of the results is deep; they fix the vocabulary of the main text.
Table~\ref{tab:theory} lists, for each result, its Lean counterpart and the measurement that
corresponds to it, and Table~\ref{tab:scope} its assumptions, its scope and what it does not
show.

\begin{table}[h]
\centering\small
\begin{tabular}{p{0.2\linewidth}p{0.33\linewidth}p{0.38\linewidth}}
\toprule
result & Lean (\texttt{lean/ExperimentsMathlib}) & measured in\\
\midrule
Prop.~\ref{prop:cert} certificates & \texttt{T9.cert\_accept}, \texttt{cert\_closed\_refute}, \texttt{cert\_robust\_reject} & every decided verdict carries both regrouping certificates (Table~\ref{tab:dist})\\
Prop.~\ref{thm:char} characterization & \texttt{T6} (\texttt{le\_*}, \texttt{strict\_*}) & closed accepts on unfaithful pairs (Tables~\ref{tab:full}, \ref{tab:sweep}), and their concentration on rejects that a certificate locates in the hypotheses (\S\ref{sec:measure})\\
Prop.~\ref{prop:vac} rejects, direction & \texttt{T6.vacuity\_only\_reject}, \texttt{compEquiv\_equivalence} & the $22+8+13$ rejects of equivalent pairs (\S\ref{sec:measure})\\
Prop.~\ref{thm:mono} prover strength & \texttt{T7.closedFA\_mono}, \texttt{closedFA\_limit}, \texttt{compFA\_bounded} & nested budgets and their cost (Tables~\ref{tab:budget}, \ref{tab:cost})\\
Prop.~\ref{prop:residue} residual collapse & \texttt{T5.residue\_comp} & the \V{residual} columns\\
Prop.~\ref{prop:informative} constancy, absorption & \texttt{T5.informative\_rejects\_constant}, \texttt{comp\_absorb\_iff}, \texttt{hypExt\_absorb} & Tables~\ref{tab:levels}, \ref{tab:hack}\\
Prop.~\ref{thm:repr} representation & \texttt{T4.ReprInvariant.collapse}, \texttt{alignedComp\_not\_reprInvariant} & none; it fixes the role of the alignment\\
Prop.~\ref{prop:split} split dependence & \texttt{T9.regroup\_eq\_pointImp}, \texttt{robust\_reject}, \texttt{vacuity\_only\_curry\_accept} & robust rejects (\S\ref{sec:measure}); rule switches (Appendix~\ref{app:sens})\\
Prop.~\ref{thm:nogo} no-go & \texttt{T6.no\_condExt} & \V{off-domain} (none on the full set)\\
Prop.~\ref{prop:dir}, Thm.~\ref{thm:sep} separation & \texttt{T2.refute\_iff\_proves}, \texttt{flt\_refute\_iff}; (ii) is a metatheorem & with Opus, 111 inequivalent pairs rejected with closed equivalence undecided, 29 the other way (refuted, but a regrouping certificate is missing)\\
Prop.~\ref{thm:transfer}, Cor.~\ref{cor:fragments} & \texttt{T3.sound\_checker\_has\_unknown}; the rest are metatheorems & fragments of the benchmark (Table~\ref{tab:frag})\\
Prop.~\ref{prop:canon} canonical forms & not formalized (proof sketch) & 299 re-presentations, \texttt{CanonV2Var} (Appendix~\ref{app:more})\\
Prop.~\ref{prop:probe} probes & \texttt{T8.implied\_iff\_comp}, \texttt{flag\_iff} & reference audit (Table~\ref{tab:audit})\\
\bottomrule
\end{tabular}
\caption{Results, their Lean counterparts and the measurements they explain.}
\label{tab:theory}
\end{table}

\begin{table}[h]
\centering\footnotesize\setlength{\tabcolsep}{3pt}
\begin{tabular}{p{0.14\linewidth}p{0.24\linewidth}p{0.2\linewidth}p{0.08\linewidth}p{0.26\linewidth}}
\toprule
result & setting and assumptions & scope & certs. & establishes, and does not\\
\midrule
Prop.~\ref{prop:cert} & Lean with Mathlib, standard axioms (classical) & the exhibited telescopes, whatever code produced them & both & a relation between the telescopes, and for an accept $\Cl R\leftrightarrow\Cl S$; not faithfulness\\
Prop.~\ref{thm:char} & any consistent $\T$ & relations over one aligned context & -- & which relation is greatest under which closure properties\\
Prop.~\ref{prop:vac} & any consistent $\T$ & aligned telescopes & both, to speak of $R,S$ & vacuity-only, weaker, stronger; with \texttt{npimp}, that one side is false; not which side is unfaithful\\
Prop.~\ref{thm:mono} & a prover is a set $O\subseteq\mathrm{Thm}(\T)$ & nested provers $O\subseteq O'$ only & -- & false-accept sets grow under nesting; across different models only a tendency\\
Prop.~\ref{prop:residue} & any consistent $\T$ & alignments that end early, provable remainders & -- & a \V{residual} certifies only the aligned hypotheses\\
Prop.~\ref{prop:informative} & $\T$ consistent & (iii): relations invariant under provably equivalent hypotheses & -- & constancy and absorption defeat such relations; not checks that restrict proofs\\
Prop.~\ref{thm:repr} & any $\T$ & relations closed under currying, folding, provable equivalence & -- & they collapse to $\Ecl$; the alignment is what separates\\
Prop.~\ref{prop:split} & any consistent $\T$ & equivalence relations containing $\Ecomp$, aligned context & -- & which rejects depend on the split (vacuity-only) and which survive regrouping\\
Prop.~\ref{thm:nogo} & any $\T$ & contexts with a point outside a hypothesis & -- & \V{off-domain} cannot be folded into accept or reject\\
Prop.~\ref{prop:dir}, Thm.~\ref{thm:sep} & $\T\supseteq\mathsf{PA}$ consistent, recursively axiomatized; (ii) metatheorem & closed against component refutation & -- & component refutations can be much shorter; not that they are found\\
Prop.~\ref{thm:transfer}, Cor.~\ref{cor:fragments} & decidable or arithmetical fragments & components in the fragment & -- & where \V{unknown} is avoidable and where it is not\\
Prop.~\ref{prop:canon} & the rules of \S\ref{sec:framework}; not formalized & re-presentations (a)--(c) only & -- & invariance under those changes, not under others (Appendix~\ref{app:more})\\
Prop.~\ref{prop:probe} & any consistent $\T$ & one statement and its weakenings & -- & idle or implied hypotheses as formal facts; an error needs a reader\\
\bottomrule
\end{tabular}
\caption{Scope of each result: what it assumes, what it quantifies over, whether it needs the two
regrouping certificates to speak about the original statements, and what it does not show. No
result shows faithfulness to an informal statement.}
\label{tab:scope}
\end{table}

\subsection{Closure properties and what they force}\label{sec:axioms}

For a statement $S=(\Gamma,H,C)$ and a predicate $P$ on $\Gamma$ write
$S{\restriction}P:=(\Gamma,H\wedge P,C)$ and $S^{\circ}:=(\Gamma,\top,C)$. A relation $\sim$ on
statements over $\Gamma$ is
(S) \emph{sound} if $R\sim S$ implies $\Cl R\leftrightarrow\Cl S$;
(Res) \emph{restriction-closed} if $R\sim S$ implies $R{\restriction}P\sim S{\restriction}P$ for every $P$;
(Vac) \emph{vacuity-respecting} if $R\sim S$ implies that $H_R$ and $H_S$ are both satisfiable or
both unsatisfiable; and
(Weak) \emph{weakening-closed} if $R\sim S$ implies $R^{\circ}\sim S^{\circ}$.
Pointwise implication is $R\Epi S:\iff\forall\bar x.\,(H_R\to C_R)\leftrightarrow(H_S\to C_S)$.
(Res) ranges over all predicates on $\Gamma$; the characterization below uses only the point
predicates $\bar x=\bar a$, so it holds for any class of restrictions that contains them.

\begin{proposition}[Characterization]\label{thm:char}
Over any context $\Gamma$, $\Ecl$, $\Epi$, $\Econd$ and $\Ecomp$ are the greatest relations
satisfying (S); (S) and (Res); (S), (Res) and (Vac); and all four, respectively. The inclusions
$\Ecomp\subsetneq\Econd\subsetneq\Epi\subsetneq\Ecl$ are strict already over $\Gamma=\mathsf{Bool}$.
\end{proposition}
\begin{proof}
Each relation satisfies its axioms by inspection. Let $\sim$ satisfy the axioms and $R\sim S$.
Restricting to $P:=(\cdot=a)$ and applying (S) gives
$(H_R(a)\to C_R(a))\leftrightarrow(H_S(a)\to C_S(a))$ for every $a$. Applying (Vac) to the same
restriction gives $H_R(a)\leftrightarrow H_S(a)$. Applying the first step to $R^{\circ}\sim S^{\circ}$
gives $C_R(a)\leftrightarrow C_S(a)$. Strictness: $(b{=}\mathsf{true},\top)$ versus
$(b{=}\mathsf{true},b{=}\mathsf{true})$; $(\bot,\bot)$ versus $(\top,\top)$; $(\top,b{=}\mathsf{true})$
versus $(\top,b{=}\mathsf{false})$.
\end{proof}

In particular $\Ecl$ and $\Epi$ relate any two true theorems, which is the folklore
collapse~\cite{beq,tla}; for true $R,S$, $\neg(R\Econd S)$ iff
$\exists a.\,\neg(H_R(a)\leftrightarrow H_S(a))$. Consequently every relation satisfying (Res) and
(Vac) separates $(\Gamma,H,C)$ from $(\Gamma,H\wedge D,C)$ whenever $D$ fails at a point of $H$, even
if $C$ holds there (a decorative hypothesis), and from its contrapositive $(\Gamma,\neg C,\neg H)$
whenever $H$ and $\neg C$ differ at a point.

\begin{proposition}[Rejects and direction]\label{prop:vac}
(i) If $R\Epi S$ but $\neg(R\Econd S)$, the hypotheses differ at some point while the claims $H\to C$
agree at every point (a \emph{vacuity-only} reject). (ii) If
$\forall\bar x.\,(H_R\to C_R)\to(H_S\to C_S)$ then $\Cl R\to\Cl S$ (the candidate is \emph{weaker});
symmetrically \emph{stronger}. (iii) $\Ecomp$ and $\Econd$ are equivalence relations, so replacing
$R$ by $R'\Ecomp R$ changes no verdict.
\end{proposition}
The checker refines every \V{reject} by the goals \texttt{pimp} ($R\Epi S$), \texttt{npimp} and the
two directions \texttt{rs}, \texttt{sr} of (ii).

\begin{proposition}[Prover strength]\label{thm:mono}
Model a prover as a set $O\subseteq\mathrm{Thm}(\T)$ and let $O\subseteq O'$. (i) All decided verdict
sets grow and the \V{unknown} set shrinks. (ii) The closed false-accept set grows, and for
$O=\mathrm{Thm}(\T)$ it contains every unfaithful pair whose closures are both provable or both
refutable in $\T$. (Agreement in truth value in an intended model is not enough when $\T$ is
incomplete.)
(iii) For every $O$ the component false-accept set is contained in
$\{(R,S)\text{ unfaithful}: R\Ecomp S\}$.
\end{proposition}

\begin{proposition}[Residual collapse]\label{prop:residue}
Let the shared prefix be $\bar x$ and write $C_R=\forall\bar y.\,H'_R\to C'_R$ and
$C_S=\forall\bar z.\,H'_S\to C'_S$ for the unaligned remainders. If $\T\vdash\forall\bar x.\,C_R$ and
$\T\vdash\forall\bar x.\,C_S$, then $\T\vdash R\Ecomp S$ iff $\T\vdash\forall\bar x.\,H_R\leftrightarrow H_S$.
\end{proposition}

\begin{proposition}[Constancy and absorption]\label{prop:informative}
Call $R$ \emph{informative} if $\T\vdash\neg\forall\bar x.\,C_R$ and $\T\vdash\neg\forall\bar x.\,\neg C_R$.
(i) An informative $R$ is certifiably not $\Ecomp$-equivalent to any $S$ over the same context whose
conclusion is constant, and every $S\Ecomp R$ is informative; since $\T$ is consistent, an
informative accept stays informative under any stronger prover. (ii) For
$R^{\mathrm{abs}}=(\Gamma,H\wedge C,C)$, $R\Ecomp R^{\mathrm{abs}}$ iff $\forall\bar x.\,H\to C$;
$\Cl R^{\mathrm{abs}}$ is provable by projection; and $R^{\mathrm{abs}}$ is informative iff $R$ is.
(iii) If $\sim$ relates $(\Gamma,H,C)$ and $(\Gamma,H',C)$ whenever $\T\vdash\forall\bar x.\,H\leftrightarrow H'$,
then $R\sim R^{\mathrm{abs}}$ for every true $R$.
\end{proposition}
Part (ii) shows that informativeness does not exclude absorption, which is why the main text
uses a separate triviality certificate. Part (iii) is immediate from $H\leftrightarrow H\wedge C$ under
$\forall\bar x.\,H\to C$.

\subsection{Representation collapse and the no-go result}

The collapse also follows from invariance under re-presentation; the two views agree because
currying hypotheses into the conclusion violates (Vac).

\begin{proposition}[Representation collapse]\label{thm:repr}
Call $\sim$ representation invariant if
(curry) $(\Gamma,H,C)\sim(\Gamma,\top,H\to C)$,
(fold) $(\Gamma,\top,\varphi)\sim((),\top,\forall\bar x{:}\Gamma.\,\varphi)$, and
(ext) $\T\vdash p\leftrightarrow q$ implies $((),\top,p)\sim((),\top,q)$.
If $\sim$ is representation invariant then $\Cl R\leftrightarrow\Cl S$ implies $R\sim S$;
$\Ecl$ is the finest such relation, and aligned $\Econd$, $\Ecomp$ are not representation
invariant.
\end{proposition}
\begin{proof}
$S\sim(\Gamma,\top,H_S\to C_S)\sim((),\top,\Cl S)$ by (curry) and (fold), and
$((),\top,\Cl R)\sim((),\top,\Cl S)$ by (ext). For the last claim take $\Gamma=\mathsf{Bool}$,
$H(b)=(b=\mathsf{true})$, $C=\top$.
\end{proof}

The proposition folds the context away. Keeping the alignment fixed and allowing only the
split to move already costs everything above $\Epi$.

\begin{proposition}[Split dependence]\label{prop:split}
Over an aligned context, call $\sim$ \emph{regrouping invariant} if
$(\Gamma,H\wedge P,C)\sim(\Gamma,H,P\to C)$ for all $H,P,C$. (i) If $\sim$ is an equivalence
relation containing $\Ecomp$ and is regrouping invariant, then $R\Epi S$ implies $R\sim S$; if
$\sim$ also satisfies (S) and (Res), then $\sim\,=\,\Epi$, which relates any two true statements
over $\Gamma$ (\texttt{T9.regroup\_eq\_pointImp}). $\Epi$ itself is regrouping invariant. (ii) If $\neg(R\Epi S)$, then
$\neg(R'\Ecomp S')$ for all $R'\Epi R$ and $S'\Epi S$, in particular for every regrouping of
either side. (iii) If $R\Epi S$, the fully curried telescopes $(\Gamma,\top,H_R\to C_R)$ and
$(\Gamma,\top,H_S\to C_S)$ are $\Ecomp$-equivalent. (iv) If $\neg(R\Epi S)$ then
$\neg\Cl R\vee\neg\Cl S$; if moreover $\forall\bar x.\,(H_R\to C_R)\to(H_S\to C_S)$, then $\neg\Cl R$.
\end{proposition}
\begin{proof}
(i) $R\Ecomp(\Gamma,\top\wedge H_R,C_R)\sim(\Gamma,\top,H_R\to C_R)$, likewise for $S$, and the two
curried telescopes are $\Ecomp$-equivalent by (iii); the converse inclusion is
Proposition~\ref{thm:char}. Regrouping preserves $H\to C$ pointwise. (ii) $\Ecomp\subseteq\Epi$ and
$\Epi$ is an equivalence relation. (iii) Both hypotheses are $\top$ and the conclusions agree
pointwise by assumption. (iv) If both closures hold, both pointwise claims hold everywhere and
are equivalent; if $\Cl R$ holds and the direction holds, the same follows.
\end{proof}

So a criterion that separates true statements over a fixed alignment cannot be invariant
under regrouping: the split is information the criterion uses, not noise it could remove.
Regrouping here is the one operation above, over a fixed context and alignment. It is not
invariance under re-presentations that a canonicalizer normalizes (Proposition~\ref{prop:canon}),
which collapses nothing; the proposition does not say that canonicalization must degrade. The
dependence is confined to one kind of verdict. A reject that is not vacuity-only
($\neg(R\Epi S)$, certified by \texttt{npimp}) survives every regrouping of either side and, by
(iv), certifies a false statement: with the regrouping certificates, the goals \texttt{npimp} and
\texttt{rs} prove $\neg R$ (\texttt{T9.robust\_reject}, \texttt{T9.not\_closed\_of\_rs}). A
vacuity-only reject instead becomes an \V{accept} once both sides are fully curried
(\texttt{T9.vacuity\_only\_curry\_accept}). Accepts depend on the split too: absorption
(Proposition~\ref{prop:informative}) is a regrouping of the conclusion into the hypotheses.

Cross-provability with an unrestricted prover, fingerprints, shadow benchmarks and
downstream-compilation tests are all $\Ecl$-invariant with a complete oracle; whatever they
separate between true statements is an artefact of oracle incompleteness. BEq+ is not: its
proofs must use the other statement, a syntactic restriction under which recall depends on how
close the candidate is to the reference in form. The no-go theorem for the third value is the other
face of the characterization: the component axioms are incompatible with extensionality of
the conclusion relative to the hypotheses.

\begin{proposition}[No-go for conditional extensionality]\label{thm:nogo}
Call $\sim$ conditionally extensional if $\forall\bar x.\,H\to(C\leftrightarrow C')$ implies
$(\Gamma,H,C)\sim(\Gamma,H,C')$. On any context with a point outside some hypothesis, no
relation is simultaneously (S), (Res), (Weak) and conditionally extensional
(\texttt{T6.no\_condExt}). $\Econd$ is conditionally extensional, so dropping (Weak) restores
consistency.
\end{proposition}
\begin{proof}
Let $H(x_0)$ fail. Inside $H$ the conclusions $\top$ and $H$ agree, so conditional
extensionality gives $(\Gamma,H,\top)\sim(\Gamma,H,H)$, and (Weak) gives
$(\Gamma,\top,\top)\sim(\Gamma,\top,H)$. By (S) and (Res), as in the proof of
Proposition~\ref{thm:char}, related statements have equivalent claims at every point; at $x_0$
this is $\top\leftrightarrow H(x_0)$, a contradiction. $\Econd$ satisfies (S), (Res) and
conditional extensionality by definition.
\end{proof}

So a criterion that is robust to Lean's junk values (total division, truncated subtraction),
which differ only outside the hypotheses, cannot also respect the raw claim. Whether a difference
located entirely outside $H$ is a faithfulness error therefore cannot in general be determined
from the two formal statements alone. This is exactly what \V{off-domain} reports.

\subsection{Refutation separation and the placement of \V{unknown}}

\begin{proposition}\label{prop:dir}
If $\T\vdash\neg\Cl S$, then $\T\vdash\neg(\Cl R\leftrightarrow\Cl S)$ iff $\T\vdash\Cl R$.
\end{proposition}

Rejecting a false candidate at the closed level thus requires proving the reference. For a
reference such as Fermat's last theorem and the candidate $a^n+b^n\ne c^n+1$, closed
refutation is as hard as FLT, while component refutation is the numeral check
$9^3+10^3=12^3+1$.

\begin{theorem}[Separation]\label{thm:sep}
(i) There are $R,S$ with $\Delta_0$ components such that $\neg(R\Econd S)$ has a constant-size
proof while $\T\nvdash\neg(\Cl R\leftrightarrow\Cl S)$ (take $\Cl R=\mathrm{Con}(\T)$).
(ii) For every computable $f$ there are infinitely many pairs whose component refutation has
size $O(|R|+|S|)$ and whose closed refutation has size $>f(|R|+|S|)$.
\end{theorem}
\begin{proof}
(i) Let $\Gamma=\mathbb N$, $C_R(x):=\neg\mathrm{Prf}_\T(x,\ulcorner\bot\urcorner)$, a $\Delta_0$
formula with $\Cl R=\mathrm{Con}(\T)$, and $C_S(x):=C_R(x)\wedge x\ne 5$, both without hypotheses.
Since $5$ codes no proof, $C_R(5)\wedge\neg C_S(5)$ is a numeral check and refutes $R\Econd S$.
$\T\vdash\neg\Cl S$ (take $x=5$), so by Proposition~\ref{prop:dir} a closed refutation is a proof
of $\mathrm{Con}(\T)$, which does not exist by the second incompleteness theorem. (ii) For a
$\Delta_0$ formula $D$ let $C_R(x):=D(x)\vee x=0$ and $C_S(x):=D(x)\wedge x\ne0$. At $x=0$,
$C_R(0)\wedge\neg C_S(0)$ has a proof of size $O(|R|+|S|)$, and $\neg\Cl S$ has one too, so a
closed refutation yields a proof of $\psi_D:=\forall x.\,(D(x)\vee x=0)$ of size larger by
$O(|R|+|S|)$. It therefore suffices that for every computable $g$ infinitely many provable
$\psi_D$ have no proof of size $\le g(|\psi_D|)$; taking $g$ to be $f$ shifted by the constant
overhead gives (ii). Suppose instead that for some computable $g$ all but finitely many provable
$\psi_D$ have a proof of size $\le g(|\psi_D|)$. Enlarging $g$ at finitely many arguments, all of
them do, and then $\T$-provability of $\psi_D$ is decidable: enumerate the finitely many
candidate proofs of size $\le g(|\psi_D|)$ and check each. We show that it is not. Let $A$, $B$
be recursively inseparable sets of indices $e$ whose computation on input $e$ halts with output
$0$, respectively $1$, and let $D_e(x)$ say that this computation has not halted with output $1$
within $x$ steps. If $e\in A$, $\T$ proves by $\Sigma_1$-completeness that the computation halts
with output $0$ at some step $t$, and, since computations are deterministic provably in $\T$,
that it never halts with output $1$; hence $\T\vdash\psi_{D_e}$. If $e\in B$, $\T$ proves in
the same way that the computation halts with output $1$ at some step $t>0$, hence
$\neg D_e(t)$ and $\neg\psi_{D_e}$, so by consistency $\T\nvdash\psi_{D_e}$. A decision
procedure for provability of $\psi_{D_e}$ would separate $A$ from $B$. All sizes here are
lengths of shortest formal proofs, not running times of a proof search.
\end{proof}

\subsection{Complexity: a transfer theorem}\label{sec:complexity}

The complexity of the three-valued question is the complexity of the component language,
no more and no less. Let $\mathfrak M$ be a structure with domain $\Gamma$ and $\mathcal C$ a
class of formulas over $\mathfrak M$ closed under Boolean combinations. Write
$\mathrm{SAT}_{\mathcal C}$ for the set of $\mathcal C$-formulas with a satisfying assignment
in $\mathfrak M$ and $\mathrm{VAL}_{\mathcal C}$ for the valid ones, those whose negation is not
in $\mathrm{SAT}_{\mathcal C}$. In this subsection the relations are read in $\mathfrak M$, the
arity $k$ of $\Gamma=\mathfrak M^k$ is part of the input, and statements are encoded as the
syntax trees of their components.

\begin{proposition}[Transfer]\label{thm:transfer}
For statements with components in $\mathcal C$, the problems $\neg(R\Ecomp S)$ and
$\neg(R\Econd S)$ are many-one equivalent to $\mathrm{SAT}_{\mathcal C}$, and $R\Ecomp S$,
$R\Econd S$ to $\mathrm{VAL}_{\mathcal C}$, by polynomial-time syntactic reductions.
\end{proposition}
\begin{proof}
$\neg(R\Ecomp S)$ is $\exists\bar x.\,\neg(H_R\leftrightarrow H_S)\vee\neg(C_R\leftrightarrow C_S)$
and $\neg(R\Econd S)$ is $\exists\bar x.\,\neg(H_R\leftrightarrow H_S)\vee(H_R\wedge\neg(C_R\leftrightarrow C_S))$,
both satisfiability of a Boolean combination of components. Conversely, for $\varphi\in\mathcal C$
take $R=(\Gamma,\top,\neg\varphi)$ and $S=(\Gamma,\top,\top)$: then $\neg(R\Ecomp S)$ and
$\neg(R\Econd S)$ both hold iff $\varphi$ is satisfiable.
\end{proof}

The complexity statements concern the mathematical decision problems under this encoding, not
the running time or completeness of the Lean proof search used in the experiments.

\begin{corollary}\label{cor:fragments}
(i) If $\mathrm{SAT}_{\mathcal C}$ is decidable, the three-valued verdict is decidable and
\V{unknown} never occurs: Presburger arithmetic over $\mathbb N$ or $\mathbb Z$, polynomial
arithmetic over $\mathbb R$ (Tarski), linear real arithmetic, dense linear orders.
(ii) For $\Gamma=\mathbb N^k$ and $\Delta_0$ components, $\neg(R\Ecomp S)$ is
$\Sigma_1$-complete and refutable by a numeral instance, $R\Ecomp S$ is $\Pi_1$-complete, and
every sound computable checker answers \V{unknown} on some true equivalence
(\texttt{T3.sound\_checker\_has\_unknown}). For $\Sigma_n\cup\Pi_n$ components, $n\ge1$,
non-equivalence is $\Sigma_{n+1}$-complete; a quantifier swap inside the conclusion raises $n$
by one.
\end{corollary}

In the decidable fragments refutation certificates need not be points: over $\mathbb R$ a
witness may be algebraic, and the Lean certificate is a quantifier-elimination or
\texttt{nlinarith} proof rather than a numeral. Comparing closures, the reject side is
incomplete even for false candidates (Theorem~\ref{thm:sep}). Comparing components,
rejection is complete for decidable components and the unavoidable \V{unknown} sits on the
accept side, where it costs coverage instead of correctness. Section~\ref{sec:exp} reports
which fragment the benchmark statements actually fall in.

\subsection{Canonical forms}

\begin{proposition}[Re-presentation invariance]\label{prop:canon}
Let $S'$ be obtained from $S$ by (a) swapping two adjacent binders neither of which occurs in
the type of the other, (b) replacing adjacent hypothesis binders $A\to B\to X$ in the shared
prefix, where $B$ and $X$ do not mention the proof of $A$, by $A\wedge B\to X$ or conversely, or (c) writing a final $\neg P$ as $P\to\mathrm{False}$
or conversely. Then the canonical telescopes of $S$ and $S'$ have the same shared prefix up
to the permutation induced by (a), $\alpha$-equivalent conclusions, and hypothesis conjunctions
equal up to associativity and commutativity of $\wedge$; so $S\Ecomp S'$ is provable by
rewriting with \texttt{and\_comm} and \texttt{and\_assoc} alone.
\end{proposition}

\begin{proof}[Proof sketch]
(c) $P\to\mathrm{False}$ is a $\Pi$ with a Prop domain, so peeling adds the binder $P$ and
leaves the body $\mathrm{False}$, which the \texttt{unpeelBot} rule rewrites to $\neg P$;
the result is the telescope of $\neg P$. (a) Peeling is order-preserving, the dependency cut
and the alignment are computed from types and names, and the prefix of $S'$ is aligned to the
prefix of $S$ binder by binder; a swap of two data binders is absorbed by the alignment, a
swap of two hypotheses permutes two conjuncts. (b) One Prop binder $A\wedge B$ and two Prop
binders $A$, $B$ are folded into right-nested conjunctions that differ by one application of
associativity; in the shared prefix no hoisting decision depends on them. Outside the shared
prefix (b) can fail: if $A$ is hoistable and $B$ is not, $A\to B\to X$ hoists $A$ into $H$ and
keeps $B$ in $C$, whereas $A\wedge B\to X$ keeps the conjunction in $C$, and the two telescopes
are then not $\Ecomp$-equivalent. This is a genuine boundary effect of
Proposition~\ref{thm:repr}, not a defect of the rules, and we report how often it occurs.
\end{proof}

The proposition is about the rules as stated; that \texttt{CanonV2} implements them is tested, not
proved, by the 299 re-presentations of Appendix~\ref{app:more}.

\section{Further experiments}\label{app:more}

\paragraph{Re-presentation (Proposition~\ref{prop:canon}).} Of the 129 distinct ProofNetVerif
references whose canonicalization succeeds, 299 surface re-presentations apply (adjacent
independent binders swapped, adjacent free hypotheses conjoined or split, a final negation
written as $C\to\mathrm{False}$). The \texttt{acc} goal of each pair $(R,R')$ is proved by
rewriting alone in 218 cases (canonical forms syntactically equal) and by
\texttt{and\_comm}/\texttt{and\_assoc} in a further 65; 9 more close under propositional
reasoning (\texttt{tauto}), all of them $\neg\exists$ rewritten as $(\exists)\to\mathrm{False}$
or one-point substitutions whose intermediate hypothesis mentioned the eliminated binder.
7 remain open, all of them the predicted failure of (b) outside the shared prefix
(Proposition~\ref{prop:canon}): a hoistable conjunct is joined to a non-hoistable one, or an
instance binder inserted by elaboration breaks the positional alignment. No applied
re-presentation produced a contradictory verdict. Table~\ref{tab:repr} splits the outcomes by
kind: a swap of binders leaves the canonical form unchanged in 178 of 218 cases, while conjoined
hypotheses need conjunct reordering in 29 of 69 and account for 5 of the 7 open ones.

\begin{table}[h]
\centering\small
\begin{tabular}{lrrrrr}
\toprule
re-presentation & applied & identical & reordered & logical & open\\
\midrule
adjacent independent binders swapped & 218 & 178 & 36 & 2 & 2\\
adjacent free hypotheses conjoined & 69 & 32 & 29 & 3 & 5\\
final negation written as $C\to\mathrm{False}$ & 11 & 7 & 0 & 4 & 0\\
conjunctive hypothesis split & 1 & 1 & 0 & 0 & 0\\
\bottomrule
\end{tabular}
\caption{Surface re-presentations of 129 ProofNetVerif references and how the \texttt{acc} goal
of each pair closes: canonical forms identical, equal up to conjunct order, or equal up to
propositional logic (\texttt{tauto}); open: not closed by any of the three
(\texttt{e5\_meta.py}).}
\label{tab:repr}
\end{table}

\paragraph{Rule sensitivity.}\label{app:sens} A copy of the canonicalizer with each rule behind
a flag (\texttt{CanonV2Var.lean}) reproduces the canonical goals of all 300 pairs of the subset
with every rule on. We switched off, one at a time, the $\neg\exists$ rewrite, the
\texttt{unpeelBot} rule ($P\to\mathrm{False}$ read as $\neg P$), one-point elimination and
hoisting, and reran the strong Opus pipeline on the pairs whose telescopes changed: 13, 8, 29
and 8 pairs, and 54 with all four off. Most of their verdicts are unchanged or move between a
decided verdict and \V{unknown}. The decided verdicts that change are of the two kinds
Proposition~\ref{prop:split} allows, both on the split. Four rejects become accepts (358 without
the $\neg\exists$ rule; 11, 189 and 997 without one-point elimination), all vacuity-only.
Seven accepts on labelled-equivalent pairs become rejects without \texttt{unpeelBot}; in each,
one side states a negation as an implication into $\mathrm{False}$ (pair 889:
$\neg\exists x\,y,\ 3x^2+2=y^2$ against \texttt{(h : 3*m\^{}2+2 = n\^{}2) : False}), so the
pointwise claims agree and the new reject is vacuity-only. No reject certified by
\texttt{npimp} changes under any variant. A more liberal one-point rule changes 7 telescopes and
no decided verdict. The rules matter, but only where Proposition~\ref{prop:split} says they
can.

\paragraph{Two rules from the audit.}\label{app:k19} After the blind audit we added two rules to
a copy of the canonicalizer: a hypothesis $\exists x,\ P\,x$ that nothing mentions is read as a
binder $x$ and a hypothesis $P\,x$, and a data or instance binder of type $T$ that nothing
mentions, and that is not known to be nonempty, is read as a hypothesis \texttt{Nonempty}~$T$.
They change the goals of 116 of the 1,140 pairs, which we reproved with the budget of
Table~\ref{tab:full}. Both artefacts they target become accepts (Putnam 2018 B4 and pair 1135),
and two of the three Dummit--Foote 1.6.23 candidates are now located in the hypotheses, where
the audit puts the difference; the third becomes \V{unknown}. Two false accepts become
non-accepts, Herstein 2.2.3 and 2.1.26, both judged unfaithful by the audit, and no pair becomes
a false accept. The three rejects of audited labelled-equivalent pairs that disappear are all
judged faithful. But the rules create artefacts of the same kind: three Artin 10.1.13 candidates
that write $\exists n,\ x^n=0$ where the reference writes \texttt{IsNilpotent x}, the same
existential behind a definition, lose their accept, since only the explicit existential is
unfolded; both statements are true, so these rejects are vacuity-only, as
Proposition~\ref{prop:split} predicts. Local rules trade one presentation choice for another, so
we keep the rules of \S\ref{sec:framework} and report this variant only as a measurement
(\texttt{k19\_rules.py}).

\paragraph{Where \V{unknown} comes from.}\label{app:unk} The strong run leaves 83 of the 300
pairs \V{unknown} (Table~\ref{tab:g1}). In 26 of them a regrouping certificate is missing, so no
component goal counts; 3 of these have a refuted closed equivalence, which without the
certificates says nothing about the telescopes. In 15 both certificates hold but no binder is
aligned, so every component goal compares the residual statements and is as hard as closed
equivalence. In 16 the alignment is partial. Only the remaining 26 are fully
aligned with every component goal open, 16 of them with a proved closed equivalence. Three
quarters of the coverage gap therefore lies in canonicalization and alignment, the same place
the \V{residual} verdicts on unfaithful pairs come from (\S\ref{sec:survive-exp}), rather than in proof search on
aligned components.

\paragraph{Fragments (Table~\ref{tab:frag}).} Corollary~\ref{cor:fragments} predicts no
\V{unknown} when the components lie in a decidable fragment. A Lean metaprogram classifies each
canonical statement by the types of its data binders, its instance binders, the constants of
its components and the presence of nonlinear operations
(\texttt{Fragment.lean}). Few benchmark statements are purely arithmetic:
on ProofNetVerif 9 of 300 pairs have Presburger or real-closed components, and all 9 are
decided (6 accept, 3 reject, 0 \V{unknown}); the 3 pairs with nonlinear integer components are
all \V{unknown}, and the remaining 288 involve functions, sets, structures or type classes,
where the theorem makes no promise and 89 of 264 structural pairs are \V{unknown}. The
competition benchmarks are different: 158 of 488 miniF2F statements (32\%) are Presburger or
real-closed and a further 53 are nonlinear integer arithmetic, while PutnamBench and
ProofNet\# are almost entirely structural.

\begin{table}[t]
\centering\small
\begin{tabular}{lrrrrrr}
\toprule
set & $n$ & Presburger & real-closed & nonlinear $\mathbb N,\mathbb Z$ & other arith. & structural\\
\midrule
ProofNetVerif (pairs) & 300 & 7 & 2 & 3 & 24 & 264\\
\quad of which \V{unknown} & & 0 & 0 & 3 & 14 & 89\\
ProofNet-Verified (pairs) & 316 & 4 & 1 & 1 & 32 & 278\\
ProofNet\# & 335 & 4 & 1 & 1 & 32 & 297\\
miniF2F & 486 & 66 & 92 & 53 & 100 & 175\\
PutnamBench & 654 & 1 & 2 & 4 & 54 & 593\\
\bottomrule
\end{tabular}
\caption{Syntactic fragment of the canonical components. A pair takes the larger class of
its two statements. ``Other arithmetic'': number types with functions outside the fragments
(primes, factorials, sums, transcendental functions) or mixed $\mathbb N$/$\mathbb R$.
Statements whose canonicalization fails are omitted from the counts.}
\label{tab:frag}
\end{table}

\begin{table}[t]
\centering\small
\begin{tabular}{lrrrc}
\toprule
injected drift & pairs & strong prover & weak prover & contradictions\\
\midrule
identity & 129 & accept 129 & accept 129 & 0 / 0\\
drop a hypothesis ($\Delta_H$) & 93 & reject 83 (89\%) & reject 57 (61\%) & 0 / 0\\
negate the conclusion ($\Delta_C$) & 71 & reject 66 (93\%) & reject 47 (66\%) & 0 / 0\\
$C\mapsto C\vee\neg H_i$ (outside $H$) & 93 & off-domain 72 (77\%) & off-domain 35 (38\%) & 0 / 0\\
\bottomrule
\end{tabular}
\caption{Drift injection on 129 ProofNetVerif references. The mutant is built by a Lean
metaprogram, so its class is known. The last drift is closed-equivalent to the reference whenever
the reference is true, so every complete closed oracle accepts it. Remaining verdicts are
\V{unknown}; no verdict contradicts the predicted class.}
\label{tab:g2}
\end{table}

\paragraph{Drift injection (Table~\ref{tab:g2}).} The characterization predicts the verdict
of each drift class. Across all 386 mutants (257 non-identity) and both provers, no verdict
contradicts the prediction; the weak prover halves coverage without introducing a single
error. Only 4 of the 221 non-accept verdicts of the strong run come from fixed tactics.

\paragraph{Ranking (F3c).} Re-scoring the autoformalization systems of
ProofNetVerif with component verdicts preserves the ranking obtained with closed
equivalence (Kendall $\tau=0.966$), but absolute scores rise by $1.2$--$1.8\times$, and three
\V{off-domain} cases (Dummit--Foote 4.5.13 and 4.5.15: $p\mid|G|$ against $p\mid 351$) appear.

\begin{table}[t]
\centering\small
\begin{tabular}{lcc}
\toprule
method (eval204, intended statement given) & relation certified / predicted correctly & certified but wrong\\
\midrule
ours, tactics + bounded witness search (no LLM) & 169 (82.8\%) & 0\\
ours, + GPT-4.1-mini & 181 (88.7\%) & 0\\
ours, + Claude Opus 5.5 & 193 (94.6\%) & 0\\
\midrule
FormalRx-CC, GPT-4.1-mini: relation correct & 114 (55.9\%) & --\\
FormalRx-CC, GPT-4.1-mini: witness compiles & 21 (10.3\%) & --\\
\quad of which still compiles when $C:=I$ or $I:=C$ & 19 & --\\
\bottomrule
\end{tabular}
\caption{External stress test on the FormalRx-CC benchmark (204 items, 121 distinct pairs,
all small definitions over $\mathbb N$, $\mathbb Z$, $\mathbb R$, \texttt{Bool} and finite
relations). Our numbers count items whose semantic relation is fully determined by
kernel-checked goals. Our re-check reproduces the published FormalRx-CC numbers (21 and 114)
exactly. The comparison is for the setting in which the intended statement is given.}
\label{tab:cc}
\end{table}

\paragraph{External stress test: FormalRx-CC (Table~\ref{tab:cc}).} FormalRx-CC~\cite{formalrxcc}
is the closest prior work. It asks a model to classify the relation between a candidate and an
intended definition and to write a Lean witness of the difference. When the intended statement is
given, its benchmark tests exactly our setting. The benchmark is small-scale: most items are
separated by a small numeral instance, so we use it to test whether free-form witnesses are
reliable, not as evidence of performance at research scale, for which ProofNetVerif
(Table~\ref{tab:g1}) is our main evidence. We state the directional goals $\forall\bar x.\,C\to I$, $\neg\forall\bar x.\,C\to I$, their
converses, term equality, and premise vacuity, and read the relation off the proved goals.
Fixed tactics plus a bounded search over small values of base-typed arguments certify 83\% of the
relation labels with no language model; GPT-4.1-mini, the model of the FormalRx-CC main experiment,
raises this to 89\%, against 10\% Lean-passing witnesses with the original protocol. Our
certificates also cover both directions of \emph{weaker} and \emph{stronger}, where the
benchmark's gold witnesses prove one.

The gap is not only in coverage. A witness is free-form Lean, so compiling it does not show
that it separates anything. We redefine the candidate with the intended body (and conversely)
and recompile. A witness that still compiles once the two objects are identical provably fails
to discriminate them; a witness that stops compiling passes only a necessary condition, so the
counts of discriminating witnesses below are upper bounds. 19 of the 21 passing witnesses still
compile after the collapse. Across all 22 archived FormalRx-CC runs on its eval51 split (five
models, with and without the intended statement and repair), 157 of 503 compiling witnesses
(31.2\%) still compile after it and hence fail to discriminate; 35 further witnesses are
complete files that re-import Mathlib and redefine both objects, which do not compile after the
definitions as the protocol requires, and counting them either way gives 29--36\%. The best configuration (GPT-4.1,
intended given, one repair) has at most 30 of 51 items with a correct relation and a
discriminating witness, where our pipeline certifies 46 without a language model. Typical
non-discriminating witnesses are definitions such as
\texttt{def witness (x y : Nat) : Prop := Intended x y → Candidate x y} or \texttt{\#eval}
commands. The same test flags one gold witness: item \texttt{dev\_0016} pairs
\texttt{b - c} with \texttt{b - x} over alpha-renamed binders, our pipeline proves the two
equal, and the gold witness compares \texttt{Candidate 0 1 3} with \texttt{Intended 0 1 0}.
Fixing the goals in advance, as we do, rules out both failure modes.

FormalRx-CC also evaluates a setting in which the model must write the intended statement
itself. Our goals certify such corrections directly: comparing GPT-4.1-mini's corrections with
the gold intended definitions using only tactics and bounded search, 132 of 204 (65\%) are proved
equivalent and 19 proved wrong. Letting the model write only the correction and certifying both
the correction and the relation with our goals yields 107 fully certified items (52\%), or 118
with GPT-4.1-mini as prover, where the original protocol obtains 13 (6.4\%) correct relations
with Lean-passing witnesses. The low scores of witness-based evaluation thus come mostly from
writing witnesses, not from diagnosis.

\section{ProofNet\# reference errors}\label{app:refs}

Table~\ref{tab:refs} lists the 24 errors of Table~\ref{tab:audit}: fourteen found by probes or by
the closed refutation on ProofNet-Verified, and ten found by pairwise rejects.

\begin{table}[h]
\centering\footnotesize
\begin{tabular}{lllp{6.2cm}}
\toprule
statement (ProofNet\# id) & split & certified probe & defect\\
\midrule
Munkres \texttt{exercise\_24\_3a} & test & $H_0$ implied & \texttt{Icc 0 1} taken in $\mathbb N$; statement false (cert.)\\
Ireland--Rosen \texttt{exercise\_4\_4} & test & $H_1$ idle & order of the primitive root written $p$; vacuous\\
Rudin \texttt{exercise\_5\_3} & test & $H_0$ idle & $\exists N$ without $N>0$; conclusion trivially true\\
Herstein \texttt{exercise\_2\_5\_43} & test & $H_0$ idle & conclusion \texttt{Nonempty (CommGroup G)} holds for every group\\
Herstein \texttt{exercise\_4\_2\_5} & valid & $H_0$ idle & conclusion \texttt{Nonempty (CommRing R)} likewise\\
Herstein \texttt{exercise\_2\_5\_37} & valid & $H_0$ implied & \texttt{IsEmpty (CommGroup G)} contradictory; vacuous\\
Herstein \texttt{exercise\_2\_8\_12} & valid & $H_0$--$H_3$ implied & same as 2.5.37\\
Ireland--Rosen \texttt{exercise\_2\_21} & valid & $H_1$ implied & case $n=0$ makes the hypotheses inconsistent\\
Rudin \texttt{exercise\_2\_24} & valid & $H_0$ implied & ``limit point'' written as ``in the closure''; false\\
Rudin \texttt{exercise\_3\_5} & test & $H_2$ implied & \texttt{limsup} without a filter compares functions\\
Rudin \texttt{exercise\_4\_8b} & valid & $H_0$ idle & informal statement inverted; formal statement false\\
Pugh \texttt{exercise\_3\_63a} & test & $H_0$ idle & series written as sequence; $p>1$ redundant\\
Munkres \texttt{exercise\_24\_2} & valid & $H_0$ implied & circle written as the unit sphere of $\mathbb R$, two points; continuity automatic; statement false (cert.)\\
Munkres \texttt{exercise\_21\_6b} & test & closed refutation (Table~\ref{tab:g4}) & \texttt{Icc 0 1} taken in $\mathbb N$; $x^n$ eventually constant; statement false (cert.)\\
\midrule
\multicolumn{4}{l}{\emph{Found by pairwise rejects on ProofNetVerif (\texttt{npimp} and \texttt{rs}; pair id)}}\\
Dummit--Foote \texttt{exercise\_2\_1\_13} & test & 311 and 7 more & hypothesis for one $x$ instead of all; $\mathbb Z\le\mathbb Q$ refutes it\\
Pugh \texttt{exercise\_2\_92} & test & 1114 and 4 more & no Hausdorff assumption; nested compact sets may have empty intersection\\
Herstein \texttt{exercise\_2\_3\_16} & test & 533, 847, 1324 & trivial group allowed; its order 1 is not prime\\
Herstein \texttt{exercise\_2\_7\_7} & test & 366 & $\varphi$ not assumed surjective\\
Munkres \texttt{exercise\_31\_1} & test & 1351 & \texttt{RegularSpace} without $T_1$; fails for two indistinguishable points\\
Rudin \texttt{exercise\_1\_8} & test & 438 & every set carries a linear order; the field structure is missing\\
Rudin \texttt{exercise\_4\_15} & test & 1328, 1145, 923 & \texttt{Monotone} for ``monotonic''; $x\mapsto -x$ refutes it\\
Rudin \texttt{exercise\_2\_29} & test & 1006 & intervals bounded; $\mathbb R$ has no such decomposition\\
Ireland--Rosen \texttt{exercise\_4\_6} & test & 1190, 557 & order of the primitive root written $p$ instead of $p-1$\\
Putnam \texttt{exercise\_2014\_a5} & test & 1259 & summand $X^n$ for $X^i$; $P_1=2X$, $P_2=6X^2$\\
\bottomrule
\end{tabular}
\caption{The 24 ProofNet\# reference errors. ``$H_k$ implied'' is a kernel proof that the other
hypotheses imply the $k$-th; ``$H_k$ idle'' is a kernel proof that the conclusion holds wherever
$H_k$ fails, together with a proof that $H_k$ is not implied; ``cert.'': \texttt{BenchErrors.lean}
proves the negation. In the lower block the negation follows from the regrouping certificate
\texttt{split\_r} and the goals \texttt{npimp} and \texttt{rs} of the listed pairs
(Proposition~\ref{prop:split}(iv)); \texttt{BenchErrors.lean} also proves each negation directly,
from the ProofNet\# statement with \texttt{Type*} instantiated at \texttt{Type}. ProofNet-Verified~\cite{pnv} independently marks all rows
except Munkres 24.3a and 21.6b unfaithful or the informal text wrong. All probe certificates are regenerated by
\texttt{scripts/real/g3\_probe.py}; the 39 flagged statements and our verdict on each are in
\texttt{data/real/g3\_flagged.txt} and \texttt{g3\_verify.txt}.}
\label{tab:refs}
\end{table}

\section{Rejects on labelled-equivalent pairs}\label{app:rejects}

All ids refer to the ProofNetVerif test split.
\emph{Strong prover (14).} The reference is false: 311 (Dummit--Foote 2.1.13), 253 and 1114
(Pugh 2.92), each certified by \texttt{npimp} and \texttt{rs}; 992 shares the reference of 253, and
its candidate, which adds \texttt{[CompactSpace $\alpha$]}, is closed-equivalent to it and hence
false as well. The candidate is false: 695, 738, 1026 and 428 (the candidate generalizes to
semirings, where it is false). A hypothesis difference the annotator tolerated: 365 (extra
\texttt{[Fintype G]}, direction \texttt{rs} only), and the vacuity-only rejects 576 (extra
\texttt{Continuous g}), 693 (\texttt{a $\bullet$ v = 0 $\to$} replaced by $\leftrightarrow$) and 913
(\texttt{map T U = U} replaced by \texttt{T '{}' U $\le$ U}). Canonicalization artefacts, both
vacuity-only: 1135 (\texttt{[Fintype G]} is data and \texttt{[Finite G]} a proposition, placed on
different sides of the boundary) and 358 (a de Morgan step plus currying).
\emph{Weak prover (3).} 1135 and 358 (artefacts) and 576 (tolerated), as above.
GPT-4.1-mini rejects no labelled-equivalent pair and Qwen3.8-27B only 358.
The label audit of the inequivalent side covered 20 pairs left \V{unknown}: 19 labels were
correct and 1 disputed.

\section{Checking the checker}\label{app:repro}

% generated by scripts/real/paper_check.py from DEFINITIONS; do not edit
\begin{table}[h]
\centering\footnotesize
\begin{tabular}{p{0.2\linewidth}p{0.74\linewidth}}
\toprule
quantity & definition\\
\midrule
full set & the 1,140 labelled ProofNetVerif test pairs whose statements elaborate in our Mathlib; 694 labelled inequivalent (U), 446 labelled equivalent (F); published ProofNetVerif labels\\
subset & 300 pairs of the full set stratified 150 / 150 (published labels); experiments too costly for the full set\\
ProofNet-Verified & 317 ProofNet\# problems with the corrected references of ProofNet-Verified; 75 originals unfaithful, 242 faithful by its own labels\\
shared SemantEq set & the 1,135 full-set pairs that SemantEq-Eval also scores; published labels and SemantEq's adjudicated labels (102 originally equivalent pairs relabelled inequivalent)\\
false accept (FA) & an accept of a pair labelled inequivalent under the label set named with it; published labels unless stated\\
false reject & a reject of a pair labelled equivalent under the named labels; a certified formal difference, not a proof of unfaithfulness\\
decided & accept, reject or off-domain with both regrouping certificates; residual is reported separately and counted as decided only where stated\\
residual & a proof of component equivalence on an alignment that ends early; never an accept in v2, while v2.1 accepts it when the two sides are the same proposition up to reducible unfolding\\
held-out set & the 1,478 labelled ProofNetVerif valid pairs whose statements elaborate (953 U, 525 F), from 183 problems none of which is in the test split; published labels; read only after the protocol was frozen\\
v2, v2.1 & v2: the criterion of Section 2, used in every development table; v2.1: v2 with regrouping certificates completed for every pair and the reducible-identity accept above, fixed on the full set before the freeze\\
located reject & a reject whose certificate is a component goal (nhyp or ncin) on telescopes with a nonempty prefix or a hoisted hypothesis; a mechanical location, counted per pair\\
component-goal reject & a reject decided by a component goal rather than by a refuted closed equivalence; counted per pair (43 on labelled-equivalent pairs, resting on 30 distinct statements)\\
blind audit & 249 items shown as 226 distinct statements; a verdict counts only when both annotators agree; rates count distinct statements\\
selection & 317 problems, 7,608 draws, 2,879 distinct candidates; classes from kernel facts against the corrected reference\\
reference error & a benchmark statement that the probes flag and that is unfaithful to its informal text after review; counted per statement\\
\bottomrule
\end{tabular}
\caption{Units, denominators and label versions. Every number in the main tables and in the checked prose is computed from the stored results under these definitions by \texttt{paper\_check.py}, which also writes a machine-readable list of the checked numbers (\texttt{paper\_manifest.json}) and of the numbers it does not yet check (\texttt{paper\_coverage.json}).}
\label{tab:units}
\end{table}


\paragraph{Units and the source of every number.} Table~\ref{tab:units} fixes what each count
means, its denominator and the labels that decide right and wrong; a false accept or false
reject is always relative to the label set named with it, and a pair that a method does not
accept is not thereby shown to differ. \texttt{paper\_check.py} recomputes the main tables and the
checked prose numbers from the stored results and fails on any difference; the numbers it does
not yet check are listed, by section, in \texttt{paper\_coverage.json}.

\paragraph{What was fixed in advance.} No analysis was pre-registered, and the canonicalization
rules were developed on the benchmarks they are evaluated on, so no number here is a held-out
estimate. The least exposed are ProofNet-Verified (Table~\ref{tab:g4}), whose labels are
independent of ProofNetVerif's, and the blind audit, whose items and questions were fixed before
annotation (\S\ref{sec:blind}). Defined after the data they are scored on had been seen, and
therefore exploratory, are the accept rules that add filters or BEq+ (Table~\ref{tab:levels}), the
combination with SemantEq (Table~\ref{tab:common}) and the two rules of Appendix~\ref{app:k19}.

All Lean checks run against one pinned Mathlib commit, and all model calls are cached by
their full request, so every table is regenerated from the cache by one script per table. Three
checks guard the batch harness itself.
\emph{Truncation.} Lean stops reading a file at a stray \texttt{import}, at \texttt{\#exit} or
inside an unterminated comment, and then reports no error for anything after it. Each batch
therefore ends with a random marker that Lean must print, and a model proof containing such a
command fails. Compiled alone, none of the 122 distinct model outputs containing a stray
\texttt{import} is accepted by Lean.
\emph{Layout.} Lean's tactic blocks are indentation-sensitive, so a model proof is dedented as a
whole and re-indented under the goal; stripping only its first line silently changes the
structure of nested blocks.
\emph{Standalone recheck.} The accepted model proofs behind the reported tables are 9,983
distinct proofs (\texttt{accepted\_proofs.jsonl.gz}). We compile each alone, in a fresh file
without the result cache, followed by \texttt{\#print axioms}; a proof passes if the file reaches
its end marker with no error and no axiom beyond \texttt{propext}, \texttt{Classical.choice} and
\texttt{Quot.sound}. All 6,544 of the first recheck passed. The 3,439 that later runs added (the
budget experiment of Table~\ref{tab:budget} and the final reruns of Table~\ref{tab:full}) were
checked by the same script, and all of them passed (\texttt{k5\_standalone\_audit.py}). One
further proof, for a closed goal of an inequivalent pair, calls \texttt{sorryAx} by name, which
the text filter does not catch; the checker flags it as using \texttt{sorry}, so the run counted
the goal as unproved and no table contains it, and alone it fails on the same axiom
(\texttt{accepted\_proofs\_sorry.json}).
\emph{Replay.} Replaying those later runs from the cache, to record their proofs, reproduced
Table~\ref{tab:budget} for both models and all 1,140 verdicts of Qwen and GPT-4.1-mini in
Table~\ref{tab:full} (\texttt{k20\_dump.py}). The versions of Lean, Mathlib and every
third-party source are in \texttt{data/real/sources.json}.
Every experiment was finally rerun from the cache: calls that had failed with a gateway error
were retried, and a single verdict changed: one inequivalent pair left \V{unknown} by DeepSeek
became a reject. Calls that fail deterministically, because the model exhausts its token budget
without an answer, count as failed proofs, which can only move a pair to \V{unknown}.

\section{Case study: a Lean-verified research counterexample}\label{app:kaplansky}

OpenAI's \texttt{openai/math} release contains Lean proofs of counterexamples to Kaplansky's
direct finiteness conjecture in characteristic~2 and in one odd characteristic. We rebuilt both
against the pinned Mathlib commit, checked by \texttt{rfl} that the proved statements are the
published challenge statements, and confirmed that they use only the three standard axioms.
In characteristic~2 the manuscript's main theorem asserts a counterexample over a
\emph{finitely presented} group with an element of odd prime order
(\texttt{KaplanskyFinitelyPresented}); the headline challenge \texttt{KaplanskyDirectFiniteness}
only requires a finitely generated group. Both are proved, so every closed metric declares them
equivalent. Their telescopes share the prefix and the conclusion $\exists a\,b,\ ab=1\wedge
ba\ne1$ and differ in the hypotheses on $G$; the trivial group is finitely generated and has no
element of odd prime order, a one-line certificate of \V{reject}. By
\texttt{true\_true\_reject\_iff} this is the only kind of reject between true statements: the
headline challenge is a strictly weaker existence claim. In odd characteristic the challenge
fixes one prime $p_0$; against a universal reading of the title the criterion reports agreement
on the image of $()\mapsto p_0$ and leaves the other primes \V{off-image}, while the manuscript
text makes clear that the displayed prime is meant. Neither result is wrong; the case shows
that ``the Lean proof checks'' certifies the challenge, not the claim.

\begin{thebibliography}{99}\small
\bibitem{beq} A.~Poiroux, G.~Weiss, V.~Kun\v{c}ak, A.~Bosselut. Reliable evaluation and benchmarks for statement autoformalization. EMNLP 2025. arXiv:2406.07222.
\bibitem{semanteq} X.~Li, P.~S.~Ammanamanchi, S.~Bhat, R.~Jamdar, Y.~Lin, S.~Biderman. SemantEq: faithful bidirectional equivalence of formal statements. arXiv:2610.32919, 2026.
\bibitem{beqliu} Q.~Liu, X.~Zheng, X.~Lu, Q.~Cao, J.~Yan. Rethinking and improving autoformalization: towards a faithful metric and a dependency retrieval-based approach. ICLR 2025.
\bibitem{symeq} Z.~Li et al. Autoformalize mathematical statements by symbolic equivalence and semantic consistency. NeurIPS 2024.
\bibitem{eucgeo} L.~Murphy et al. Autoformalizing Euclidean geometry. ICML 2024.
\bibitem{paraphrase} W.~Feng, E.~Lou, A.~Sharma. Characterizing paraphrase-induced failures in Lean~4 autoformalization. arXiv:2604.23135, 2026.
\bibitem{bpf} N.~I.~S.~Mohammad, T.~Sheikh. The faithfulness gap: certifying semantic equivalence between natural-language and formal mathematical statements (BPF, DriftBench). arXiv:2606.16541, 2026.
\bibitem{shadowbench} H.~Han, J.~Kim, S.~Park, et al. ShadowBench: toward reliable automatic evaluation of semantic alignment in autoformalization. EMNLP 2026. arXiv:2608.29270.
\bibitem{t2} J.~Kim, H.~Han, S.~Hwang. Benchmarking testing in automated theorem proving. ACL 2026 (Industry). arXiv:2604.23698.
\bibitem{formalrxcc} Counterexample-Carrying FormalRx: Lean-checkable witnesses for semantic misalignment in autoformalization. Anonymous artifact \url{https://anonymous.4open.science/r/ccf-F4C7/}, 2026.
\bibitem{gted} Generalized tree edit distance (GTED): a faithful evaluation metric for statement autoformalization. arXiv:2507.07399, 2025.
\bibitem{assess} ASSESS: a semantic and structural evaluation framework for statement similarity. arXiv:2509.22246, 2025.
\bibitem{formalalign} J.~Lu, Y.~Wan, Y.~Huang, J.~Xiong, Z.~Liu, Z.~Guo. FormalAlign: automated alignment evaluation for autoformalization. ICLR 2025.
\bibitem{beyondcomp} K.~Zhang, P.~Gallardo~Candela, S.~Murthy, Y.~Xie, Z.~Wang, M.~Raissi. Beyond compilation: evaluating faithful natural-language-to-Lean statement formalization. arXiv:2606.31002, 2026.
\bibitem{atf} Autoformalizer with tool feedback. arXiv:2510.06857, 2025.
\bibitem{faults} P.~S.~Ammanamanchi, S.~Bhat, S.~Biderman. Faults in our formal benchmarking: dataset defects and evaluation failures in Lean theorem proving. PMLR 306, 2026. arXiv:2606.29493.
\bibitem{proofnet} Z.~Azerbayev et al. ProofNet: autoformalizing and formally proving undergraduate-level mathematics. arXiv:2302.12433, 2023.
\bibitem{proofnetverif} ProofNetVerif and ProofNet\#, released with~\cite{beq}. Hugging Face datasets \texttt{PAug/ProofNetVerif}, \texttt{PAug/ProofNetSharp}, 2024--2025.
\bibitem{pnv} Y.~Min, M.~An, C.~Si, O.~Bastani. Ground False: uncovering errors in formal mathematics benchmarks. NeurIPS 2026 Evaluations and Datasets Track. Data: ProofNet-Verified, \url{https://github.com/marcusm117/ProofNet-Verified}.
\bibitem{meidouglas} A.~Mei, M.~R.~Douglas. Beyond LLM judge: mechanical verification for agentic autoformalization. arXiv:2610.20573, 2026.
\bibitem{onlinerl} S.~Sorg, W.~Li, S.~Banerjee. Online reinforcement learning for autoformalization. NeurIPS 2025 Workshop on Multimodal Algorithmic Reasoning.
\bibitem{verina} Z.~Ye et al. VERINA: benchmarking verifiable code generation. ICLR 2026. arXiv:2505.23135.
\bibitem{clever} A.~Thakur et al. CLEVER: a curated benchmark for formally verified code generation. NeurIPS 2025 Datasets and Benchmarks. arXiv:2505.13938.
\bibitem{mathesis} X.~Yu et al. Mathesis: towards formal theorem proving from natural languages. ICLR 2026. arXiv:2506.07047.
\bibitem{beer} I.~Beer, S.~Ben-David, C.~Eisner, Y.~Rodeh. Efficient detection of vacuity in temporal model checking. Formal Methods in System Design 18:141--163, 2001.
\bibitem{kupfermanvardi} O.~Kupferman, M.~Y.~Vardi. Vacuity detection in temporal model checking. STTT 4(2):224--233, 2003.
\bibitem{gurfinkel} A.~Gurfinkel, M.~Chechik. How vacuous is vacuous? TACAS 2004.
\bibitem{namjoshi} K.~S.~Namjoshi. An efficiently checkable, proof-based formulation of vacuity in model checking. CAV 2004.
\bibitem{minif2f} K.~Zheng, J.~M.~Han, S.~Polu. miniF2F: a cross-system benchmark for formal Olympiad-level mathematics. ICLR 2022. Lean~4 port of DeepSeek-Prover-V1.5.
\bibitem{minif2fv2} A.~Ospanov, F.~Farnia, R.~Yousefzadeh. miniF2F-Lean revisited: reviewing limitations and charting a path forward. NeurIPS 2025. arXiv:2511.03108.
\bibitem{minif2fdm} Google DeepMind. miniF2F, Lean~4 fork with corrected formalizations. \url{https://github.com/google-deepmind/miniF2F}, 2025.
\bibitem{putnambench} G.~Tsoukalas et al. PutnamBench: evaluating neural theorem-provers on the Putnam mathematical competition. NeurIPS 2024 Datasets and Benchmarks.
\bibitem{tla} M.~J.~Min, M.~He, Z.~Li, Z.~Yi, S.~Malik, A.~Gupta, X.~Si, O.~Bastani. Theory-level autoformalization: from isolated statements to unified formal knowledge bases. ICML 2026 (position paper). arXiv:2607.13292.
\bibitem{stepfun} StepFun-Formalizer: unlocking the autoformalization potential of LLMs through knowledge-reasoning fusion. arXiv, 2025.
\bibitem{formarl} FormaRL: enhancing autoformalization with no labeled data. arXiv, 2025.
\bibitem{reform} G.~Chen et al. ReForm: reflective autoformalization with prospective bounded sequence optimization. ICLR 2026. arXiv:2510.24592.
\bibitem{criticlean} Z.~Peng et al. CriticLean: critic-guided reinforcement learning for mathematical formalization. ACL 2026. arXiv:2507.06181.
\end{thebibliography}

\end{document}
